% concatenated from main_modular.tex, 01_introduction.tex, 02_section_preliminaries.tex, 03_section_curvature_phase.tex, 04_section_barrier_method.tex, 05_section_bowl_asymptotics.tex, 06_section_catenoidal.tex, 07_section_uniqueness_barriers.tex, bibliography_amsrefs.tex
% !TeX encoding = UTF-8
\documentclass{amsart}

\usepackage[T1]{fontenc}
\usepackage[utf8]{inputenc}
\usepackage[english]{babel}

\usepackage{amsthm,amssymb,verbatim}
\usepackage{mathtools}
\usepackage{mathrsfs}
\usepackage[shortlabels]{enumitem}
\usepackage{graphicx}
\usepackage{subcaption}
\usepackage{xcolor}
\usepackage{tikz}
\usetikzlibrary{arrows.meta,calc}
\usepackage{pgfplots}
\pgfplotsset{compat=1.16}
\usepackage{chngcntr}
\usepackage{apptools}
\usepackage{ragged2e}
\usepackage{amscd}
\usepackage{hyperref}
\usepackage[msc-links,backrefs]{amsrefs}
\usepackage[foot]{amsaddr}
\emergencystretch=2em

% Image fallback: the manuscript source refers to external figure files.
% If the files are present in the same folder, they are included normally.
% If not, LaTeX inserts a framed placeholder so the article still compiles.
\makeatletter
\let\oldincludegraphics\includegraphics
\renewcommand{\includegraphics}[2][]{%
  \IfFileExists{#2}{\oldincludegraphics[#1]{#2}}{%
    \fbox{\parbox[c][0.28\textwidth][c]{0.68\textwidth}{\centering Missing figure file:\par\texttt{\detokenize{#2}}}}%
  }%
}
\makeatother

\newcommand{\rr}{{\mathbb R}}
\newcommand{\zz}{{\mathbb Z}}
\newcommand{\Sp}{{\mathbb S}}
\newcommand{\cc}{{\mathbb C}}
\newcommand{\Tr}[2]{{\mbox{Trace}}_{#1}\left(#2\right)}
\newcommand{\limit}[2]{\lim_{#1\rightarrow #2}}
\newcommand{\eps}{\varepsilon}
\newcommand{\h}{h^{\omega}}
\newcommand{\g}{g^{\omega}}
\newcommand{\f}{f^{\omega}}
\newcommand{\s}{s^{\omega}}
\newcommand{\set}[1]{\left\{#1\right\}}
\newcommand{\pare}[1]{\left(#1\right)}
\newcommand{\abs}[1]{\left|#1\right|}
\newcommand{\norm}[1]{\abs{\abs{#1}}}
\newcommand{\dive}[2]{\mbox{div}_{_{#1}}\pare{#2}}
\newcommand{\W}{{\mathcal W}}
\newcommand{\A}{{\mathcal A}}
\newcommand{\K}{{\mathcal K}}
\newcommand{\vv}{{\mathbf v}}
\newcommand{\ww}{{\mathbf w}}
\newcommand{\uu}{{\mathbf u}}
\newcommand{\barphi}{{\boldsymbol{\varphi}}}
\newcommand{\bpsi}{{\boldsymbol{\psi}}}
\newcommand{\hh}{{\mathbf h}}
\newcommand{\FF}{{\mathbf F}}
\newcommand{\dist}[1]{\mbox{dist}\left({#1}\right)}
\newcommand{\pI}[1]{\left <{#1}\right >}
\newcommand{\restri}[2]{\left.#1\right|_{#2}}
\newcommand\restr[2]{{%
        \left.\kern-\nulldelimiterspace
        #1
        \vphantom{\big|}
        \right|_{#2}%
}}

\newcommand{\calcheading}[1]{%
    \medskip\noindent\textbf{#1}\par\nobreak\smallskip%
}

\theoremstyle{plain}
\newtheorem{theorem}{Theorem}[section]
\newtheorem{lemma}[theorem]{Lemma}
\newtheorem{proposition}[theorem]{Proposition}
\newtheorem{corollary}[theorem]{Corollary}
\newtheorem{claim}[theorem]{Claim}
\newtheorem{step}[theorem]{Step}

\newtheoremstyle{mainintrothm}
{\topsep}
{\topsep}
{\itshape}
{}
{\bfseries}
{.}
{.5em}
{#1}
\theoremstyle{mainintrothm}
\newtheorem{maintheorema}{Theorem A}
\renewcommand{\themaintheorema}{A}
\newtheorem{maintheoremb}{Theorem B}
\renewcommand{\themaintheoremb}{B}
\newtheorem{maintheoremc}{Theorem C}
\renewcommand{\themaintheoremc}{C}
\newtheorem{maintheoremd}{Theorem D}
\renewcommand{\themaintheoremd}{D}

\theoremstyle{definition}
\newtheorem{definition}[theorem]{Definition}
\newtheorem{assumption}[theorem]{Assumption}
\newtheorem{notation}[theorem]{Notation}
\newtheorem{remark}[theorem]{Remark}
\newtheorem{op}[theorem]{Open Problem}
\newtheorem{example}[theorem]{Example}

\title[Rotationally symmetric translating solitons]%
{Rotationally symmetric translating solitons of fully nonlinear extrinsic geometric flows: Classification and Applications}
\author{Jose Torres Santaella}
\keywords{Translating solitons; Fully nonlinear extrinsic flow; Maximum principles; Elliptic PDEs.}
\subjclass[2010]{53A10, 53C21, 53C42, 53E10, 35J60, 58J70}

\begin{document}

\begin{abstract}
We develop a rotational theory for translating solitons of fully nonlinear
extrinsic curvature flows in Euclidean space.  Furthermore, we obtain fine asymptotic
expansions for bowl-type translators in nondegenerate and degenerate regimes.
On the other hand, we also introduce a signed-neck framework which yields the construction and
classification of catenoidal-type translators, distinguishing complete embedded
families from maximal admissible pieces according to the selected signed
branch.  As applications, we prove uniqueness results for strictly convex
entire graphical translators with prescribed bowl-type asymptotics and obtain
catenoidal-barrier nonexistence results for bounded graphical translators.
\end{abstract}

\maketitle

\section{Introduction}
\label{sec:introduction}

Extrinsic geometric flows in Euclidean space form a broad class of geometric
evolution equations in which immersed hypersurfaces are deformed by prescribing
their normal velocity in terms of the principal curvatures.  These flows provide
a natural framework for studying convexity, curvature concentration, and
self-similar models for singularity formation; see, for instance,
\cite{mantegazza2011lecture,andrews-mccoy-zheng-contracting}.

\medskip

The reference case is the mean curvature flow, where the normal velocity is the
mean curvature and the evolution is the negative gradient flow of the area
functional \cite{mantegazza2011lecture}.  Fully nonlinear extrinsic curvature
flows enlarge this framework by replacing the mean curvature with a symmetric
curvature function $\gamma$ of the principal curvatures.  In this article, we
consider curvature functions $\gamma:\overline{\Gamma}\to[0,\infty)$, where
\begin{align*}
	\Gamma_+=\{\lambda\in\mathbb R^n:\lambda_i>0\text{ for every }i\}
	\subset
	\Gamma
	\subset
	\mathbb R^n
\end{align*}
is a fixed symmetric cone.  The cone $\Gamma$ determines the admissible
curvature regime and controls the ellipticity of the associated stationary and
self-similar equations.

\medskip

More precisely, let $F_0:\Sigma^n\to\mathbb R^{n+1}$ be an oriented immersion,
and let $\gamma:\overline\Gamma\to[0,\infty)$ be a curvature function in the
sense of Definition~\ref{def:alpha-positive-curvature-function}.  The initial
hypersurface $\Sigma_0=F_0(\Sigma)$ evolves by the $\gamma$-flow when there is a
one-parameter family of immersions
$F:\Sigma\times[0,T)\to\mathbb R^{n+1}$ satisfying
\begin{align}
	\label{eq:flow}
	\begin{cases}
		\langle \partial_tF(x,t),\vec{\nu}(x,t)\rangle
		=
		-\gamma(\lambda(x,t)),
		& (x,t)\in\Sigma\times(0,T),\\
		F(x,0)=F_0(x),
		& x\in\Sigma.
	\end{cases}
\end{align}
Here $\Sigma_t=F(\Sigma,t)$, $\vec{\nu}(x,t)$ is the chosen unit normal along
$\Sigma_t$, and
$\lambda(x,t)=(\lambda_1(x,t),\ldots,\lambda_n(x,t))\in\Gamma$ denotes the
principal curvature vector with respect to $\vec{\nu}(x,t)$.

\medskip

The solutions studied in this article are translating solitons.  These are
solutions of \eqref{eq:flow} whose evolution is, up to tangential
reparametrization, translation in a fixed direction.  After a rotation of the
ambient space, we take this direction to be $\vec e_{n+1}$ and write
\begin{align*}
	F(x,t)=F_0(x)+t\vec e_{n+1},
	\qquad
	(x,t)\in\Sigma\times\mathbb R.
\end{align*}
Each time slice then satisfies the translator equation
\begin{align}
	\label{eq:gamma-translator-equation}
	\gamma(\lambda)=\langle \vec{\nu},\vec e_{n+1}\rangle .
\end{align}
Thus a $\gamma$-translator may be viewed simultaneously as a self-similar
solution of the geometric evolution equation and as a hypersurface satisfying a
Weingarten-type equation, where a function of the principal curvatures is
prescribed by the vertical angle function.  Throughout the paper, a hypersurface
is called $\Gamma$-admissible, with respect to the chosen normal, when its
principal curvature vector belongs to $\Gamma$ at every point where the equation
is considered.

\medskip

In the mean-curvature case, translating solitons have been studied extensively
both as models for type-II singularities and through their connection with
minimal surface theory.  The complete graphical theory in $\mathbb R^3$ was
developed through the work of Wang, Spruck--Xiao, and
Hoffman--Ilmanen--Mart\'in--White
\cite{wang2011convex,spruck-xiao,hoffman2019graphical,HIMW-notes}.  In
particular, Hoffman--Ilmanen--Mart\'in--White classified complete translating
graphs and constructed the $\Delta$-wing translators over strips
\cite{hoffman2019graphical,HIMW-notes}.  Further noncompact examples, including
Scherk-like translators, tridents, semigraphical translators, and annular
translators, were constructed and analyzed by Hoffman--Mart\'in--White and
collaborators
\cite{HMW-scherk,HMWS-tridents,HMW-translating-annuli,HMWS-deltawings}.  The
classification of semigraphical translators was completed by
Mart\'in--S\'aez--Tsiamis--White \cite{martin2024classification}.  In higher
dimensions, Bourni--Langford--Tinaglia constructed convex translators in slab
regions and studied their convexity, regularity, asymptotics, and reflection
symmetry \cite{bourni2020existence}.

\medskip

For fully nonlinear curvature functions, the translator theory is more recent.
A basic class is given by the elementary symmetric functions
\begin{align*}
	S_r(\lambda)
	=
	\sum_{1\leq i_1<\cdots<i_r\leq n}
	\lambda_{i_1}\cdots\lambda_{i_r},
	\mbox{ with }
	1\leq r\leq n,
\end{align*}
defined on the G{\aa}rding cone $\Gamma_r=\{\lambda\in\mathbb R^n:S_i(\lambda)>0\text{ for }i=1,\ldots,r\}$. The corresponding $r$-mean-curvature translators are the translating solitons
associated with $\gamma=S_r$.  Translators for these higher-order mean curvature
flows were studied by de Lima--Pipoli in product ambient spaces, where they
constructed rotational bowl-type and catenoid-type examples
\cite{lima2025translators}.  In Euclidean space, Rengaswami classified
rotational bowl-type translators for general fully nonlinear curvature functions
by reducing the translator equation to an implicit branch problem determined by
the rotational principal curvatures \cite{rengaswami2024classification}.
Half-space theorems for translating solitons of the $r$-mean-curvature flow
were later obtained by Alencar--Bessa--Silva Neto
\cite{alencar2026halfspace}.  In previous work, we developed
maximum-principle methods for $\gamma$-translators, including tangency
principles, nonexistence consequences, and convexity-type results
\cite{Yo,Yo2}.  The present article continues this program by studying
rotationally symmetric $\gamma$-translators beyond the bowl case, with emphasis
on catenoidal families, phase continuation, asymptotics, and applications of the
comparison theory.

\medskip

For the rotational reduction used throughout the paper, we consider a
hypersurface generated by rotating a curve $c(s)=(r(s),u(s))$, with $r(s)>0$,
around the vertical axis.  Away from the axis, it is parametrized by
\begin{align*}
	F(s,\omega)=(r(s)\omega,u(s)),
	\qquad
	\omega\in\mathbb S^{n-1}.
\end{align*}
With respect to a chosen unit normal, the principal curvature vector has the
form $(x,y,\ldots,y)$, where $x$ is the meridional curvature and $y$ is the
rotational curvature.  This reduces the curvature function to two variables by
setting
\begin{align*}
	\widetilde\gamma(x,y)=\gamma(x,y,\ldots,y).
\end{align*}
On each portion where the generating curve is written as a vertical graph
$u=u(r)$, with $v(r)=u'(r)$, the translator equation becomes an implicit
first-order equation for $v$ governed by the level structure of
$\widetilde\gamma$.  Positive graphical ends are controlled by the admissible
positive branch of the level set $\widetilde\gamma=1$, which is the branch used
in Rengaswami's classification of rotational bowl-type translators
\cite{rengaswami2024classification}.

\medskip

The behavior of this positive branch depends on the cylindrical curvature
direction.  In reduced variables this direction is $(0,1)$.  The normalized
$1$-nondegenerate case corresponds to $\widetilde\gamma(0,1)=1$ and leads to an
entire positive branch approaching the cylindrical endpoint.  The
$1$-degenerate case places the cylindrical direction on the boundary of the
admissible cone and changes the asymptotic behavior of the branch.  This
distinction produces the different positive-end regimes treated in
Section~\ref{sec:positive-branch-asymptotics}.

\medskip

Our first result refines the rotational bowl classification at the level of
positive-end asymptotics.  Starting from Rengaswami's branch formulation, we
obtain higher-order expansions in the nondegenerate case and leading or
second-order expansions in the degenerate regimes.  First-order asymptotics for
fully nonlinear bowls also appear in the ancient-solution work of
Cogo--Lynch--Vi\v{c}'anek Mart'inez
\cite[Appendix~A, Proposition~A.1]{cogo-lynch-vicanek2023rotational}.  The
results below refine this picture in the branch variables used here, and provide
the quantitative separation needed later in the Alexandrov argument.

\medskip

\begin{maintheorema}
	\label{Thm A}
	Let $B={(x,u_B(|x|)):x\in\mathbb R^n}$ be a smooth $\Gamma$-admissible
	rotational bowl-type $\gamma$-translator in one of the regimes covered by
	Rengaswami's classification, and set $r=|x|$. In the normalized $1$-nondegenerate case $\widetilde\gamma(0,1)=1$, the
	radial derivative satisfies
	\begin{align*}
		u_B'(r)
		=
		r^\alpha
		-
		\frac{a}{r^\alpha}
		+
		\frac{b}{r^{3\alpha}}
		+
		o(r^{-4\alpha})
		\qquad
		\text{as }r\to\infty,
	\end{align*}
	where $a$ and $b$ are determined by the expansion of the positive branch at
	the cylindrical endpoint. On the other hand, in the $1$-degenerate case, assume that the positive branch satisfies
	\begin{align*}
		g_+(y)
		=
		c_\gamma y^{-k_\gamma}(1+o(1))
		\qquad
		\text{as }y\to\infty,
	\end{align*}
	with $c_\gamma>0$ and $k_\gamma>0$.  When $k_\gamma>2\alpha-1$, one has
	\begin{align*}
		u_B'(r)
		=
		A_\gamma r^{d_\gamma}(1+o(1))
		\qquad
		\text{as }r\to\infty,
	\end{align*}
	where $d_\gamma=\frac{\alpha(k_\gamma+1)}{k_\gamma+1-2\alpha}$, $A_\gamma^{\frac{k_\gamma+1-2\alpha}{\alpha}}=\frac{c_\gamma}{d_\gamma}$. When $k_\gamma=2\alpha-1$, one has
	\begin{align*}
		\log u_B'(r)
		=
		\frac{c_\gamma}{2\alpha}r^{2\alpha}
		+
		o(r^{2\alpha})
		\qquad
		\text{as }r\to\infty.
	\end{align*}
	Under a balanced second-order decay assumption on $g_+$, the power-growth
	case also admits a second-order correction, including a resonant logarithmic
	correction in the critical algebraic case.
\end{maintheorema}

\medskip

The full statement is proved in Theorem~\ref{thm:positive-bowl-asymptotics}.  The argument combines the barrier method for the graphical phase variable with uniform expansions of the positive branch.  In the nondegenerate case, the phase remains near the cylindrical endpoint; in the degenerate case, the decay of $g_+(y)$ as $y\to\infty$ determines the positive-end growth.

\medskip

The asymptotic theory also yields a uniqueness result for entire graphical
translators with prescribed bowl-type behavior at infinity.  The moving-plane
argument follows the Alexandrov strategy used by
Mart'in--Savas-Halilaj--Smoczyk for translators asymptotic to the translating
paraboloid \cite{Paco_2014}.  The expansions provide the separation at infinity
needed to start the reflection, while the tangency principle for
$\gamma$-translators \cite[Theorem~1.4]{Yo} prevents a first contact unless the
translators coincide.

\medskip

\begin{maintheoremb}
	Let $B={(x,u_B(|x|)):x\in\mathbb R^n}$ be a smooth strictly convex
	$\Gamma$-admissible rotational bowl-type $\gamma$-translator.  Assume that
	the positive end of $B$ belongs to one of the regimes where the Alexandrov
	moving-plane argument applies: the normalized $1$-nondegenerate case with
	$\alpha\ge1$, the $1$-degenerate power-growth case with $d_\gamma\ge1$, or
	the critical $1$-degenerate case.  Let
	$\Sigma={(x,u(x)):x\in\mathbb R^n}$ be a complete strictly convex entire
	$\Gamma$-admissible $\gamma$-translator.  Assume that there exists
	$c\in\mathbb R$ such that $\Sigma$ is $C^2$-asymptotic to
	$B+c\vec e_{n+1}$.  Then $\Sigma=B+c\vec e_{n+1}$.
\end{maintheoremb}

\medskip

The second part of the paper develops the signed rotational theory needed for
catenoidal-type translators.  At a catenoidal neck, the rotational curvature is
positive while the meridional curvature has the opposite sign.  Thus the positive
branch of $\widetilde\gamma$ does not describe the full phase diagram.  We
therefore introduce a signed rotational extension $\widehat\gamma$ and organize
the construction around three levels.  The level $\widetilde\gamma=1$ controls
positive graphical ends, the signed level $\widehat\gamma=0$ fixes the neck,
and the level $\widehat\gamma=-1$ describes the reoriented lower side.  The
formal signed data, their regularity, and the model phase diagrams are given in
Section~\ref{sec:curvature-functions-branches}. 

\medskip

The resulting catenoidal classification separates complete embedded families
from maximal selected-component pieces.  The local neck is produced by the
zero level $\widehat\gamma=0$.  Its upper side enters the positive branch and
therefore has a positive bowl-type end.  After reorientation, the lower side may
return to the positive branch, approach the selected zero trace, remain in a flat
complete fixed-level component, or stop at the boundary of the selected lower
component.

\medskip

\begin{maintheoremc}
	Assume that the signed neck data are regular, that the upper-entry sector is
	compatible with the positive rotational trace, and that $R$ is an admissible
	neck radius.  Then the signed zero trace produces a local catenoidal neck of
	radius $R$.  Its upper side reaches the positive branch at finite radius and
	continues as a positive bowl-type end.
		\medskip
	
	After reorienting the lower side, the continuation is governed by the
	selected fixed-level diagram.  The following alternatives occur.
	
	\begin{enumerate}[label=\textup{(\roman*)}]
		\item The lower side reaches the positive branch at finite radius and
		continues as a second positive bowl-type end.  This gives a complete
		embedded catenoidal translator with two positive bowl-type ends.
		
		\item The lower side approaches the selected zero ray at the origin of the
		signed phase plane.  This gives a complete embedded catenoidal translator
		with one positive bowl-type end and one zero-trace radial end.
		
		\item The lower side remains on a flat complete fixed-level component.
		When the reoriented output level belongs to that component, this gives a
		complete embedded catenoidal translator.
		
		\item Endpoint and pole-limited lower components give maximal admissible
		rotational pieces inside the selected signed component, rather than
		complete catenoidal translators.
	\end{enumerate}
	
\end{maintheoremc}

\begin{figure}[htbp]
	\centering
	\begin{minipage}{0.48\textwidth}
		\centering
		\includegraphics[width=\textwidth]{classic_catenoid_mcmanus.png}
	\end{minipage}
	\hfill
	\begin{minipage}{0.48\textwidth}
		\centering
		\includegraphics[width=\textwidth]{nobel_catenoid_mcmanus.png}
	\end{minipage}
	\caption{\footnotesize Schematic catenoidal profiles and their rotational models for two
		representative signed-neck continuations.  The left panels show the
		generating curves in the $(r,u)$-plane, and the right panels show the
		corresponding rotational hypersurfaces.  Images courtesy of Ignacio
		McManus.}
	\label{fig:intro-catenoidal-profiles}
\end{figure}


The complete catenoidal families obtained in Theorem~C yield barriers for
graphical translators over bounded domains.  For a complete embedded
catenoidal family ${W_R:R>0}$, let $r_R^{\mathrm{ent}}>R$ be the first radius
at which the upper side enters the positive branch, and let
$W_{R,\mathrm{cvx}}^+$ denote the corresponding positive convex tail.  The
family satisfies the small-neck entry condition when $\lim\limits_{R\to0^+} r_R^{\mathrm{ent}}=0$. This condition supplies positive elliptic tails whose entry radii are arbitrarily small, which allows the comparison to be placed over any bounded graphical domain.

\medskip

\begin{maintheoremd}
	Assume that there exists a complete embedded catenoidal family
	${W_R:R>0}$ satisfying the small-neck entry condition.  Then there is no
	complete convex $\Gamma$-admissible graphical $\gamma$-translator
	$\Sigma=\operatorname{graph}(u)$, with $u:\Omega\to\mathbb R$, over a
	bounded connected domain $\Omega\subset\mathbb R^n$, satisfying $u(x)\to+\infty$ as $x\to\partial\Omega$.
\end{maintheoremd}

\medskip

The proof centers a small-neck catenoidal tail outside the cylinder over
$\Omega$ and moves the bounded-domain graph vertically until first contact.  The
boundary blow-up forces the contact to occur at an interior point, and the
tangency principle for $\gamma$-translators then gives the contradiction.

\calcheading{Organization of the paper:}
Section~\ref{sec:preliminaries} fixes the geometric notation and derives the
rotational form of the translator equation in the vertical and horizontal
graphical regimes.  Section~\ref{sec:curvature-functions-branches} introduces
the curvature-function classes, the positive rotational branch, and the signed
phase data needed for the catenoidal construction.  Section~\ref{sec:barrier-method}
develops the comparison tools for the first-order graphical equation, including
level barriers and conical phase barriers.  Section~\ref{sec:positive-branch-asymptotics}
proves the positive-end asymptotics stated in Theorem~A.
Section~\ref{sec:catenoidal-translators} constructs the local catenoidal necks
and classifies their global continuations according to the selected signed
branch, as summarized in Theorem~C.
Section~\ref{sec:uniqueness-catenoidal-barriers} contains the Alexandrov
reflection uniqueness theorem and the catenoidal-barrier nonexistence result
stated in Theorems~B and~D.  Finally, Section~\ref{sec:conclusions} records
conclusions, open problems, and further directions.

\medskip

\noindent\textbf{Acknowledgements.}
The author is grateful to S. Rengaswami, M. S\'aez, and S. Vald\'es for their
support and for many fruitful conversations related to this work.  The author
also thanks Ignacio McManus for preparing the catenoidal-profile illustrations
used in the introduction and in the barrier section.

\section{Geometric preliminaries and rotational ansatz}
\label{sec:preliminaries}

In this section, we fix the geometric notation and derive the rotational form of
the translator equation.  The ambient space is $\mathbb R^{n+1}$, and the
orientation conventions stated below will be used throughout the paper.

\subsection{Background in extrinsic geometry}\,
\label{subsec:extrinsic-notation}

Let $F_0:\Sigma^n\to\mathbb R^{n+1}$ be an oriented immersion of a smooth
manifold $\Sigma$, and write $\Sigma_0=F_0(\Sigma)$.  Let $\vec\nu$ be a chosen
unit normal along $\Sigma_0$.  The induced metric is given by $g(\vec v,\vec w)=\langle dF_0(\vec v),dF_0(\vec w)\rangle$, and $\nabla$ denotes the Euclidean connection.

\medskip

At each point $p\in\Sigma_0$, the Weingarten operator is $\mathcal W_p:T_p\Sigma_0\to T_p\Sigma_0$, given by  $\mathcal W_p(\vec{X})=-\nabla_{\vec{X}}\vec\nu$. Given a local tangent frame $\vec\tau_1,\ldots,\vec\tau_n$, one has
$g_{ij}=\langle\vec\tau_i,\vec\tau_j\rangle$, and the second fundamental form is
$h_{ij}=-\langle\nabla_{\tau_i}\vec\nu,\vec\tau_j\rangle=\langle\nabla_{\tau_i}\vec\tau_j,\vec\nu\rangle $. Moreover, with Einstein summation convention, the local matrix of $\mathcal W$ is $h^i_{\ j}=g^{ik}h_{kj}$.  The eigenvalues of $\mathcal W_p$, computed with respect to $\vec\nu$, are the principal curvatures and will be denoted by $\lambda(p)=(\lambda_1(p),\ldots,\lambda_n(p))$.


\subsection{Rotationally symmetric hypersurfaces}\,
\label{subsec:rotational-generating-curves}

The rotational ansatz reduces the hypersurface geometry to the geometry of a
plane curve.  Let $\mathtt c(s)=(r(s),u(s))$ be a regular curve in
$\mathbb R^2$ with $r(s)>0$, parametrized by arclength on an interval
$(s_-,s_+)$.  Thus $\dot r(s)^2+\dot u(s)^2=1$, where $\dot{(\cdot)}=\frac{d}{ds}$. 

\medskip

 A rotational hypersurface generated by a regular
$\mathtt c$ around the vertical axis is parametrized by $F(s,\omega)=(r(s)\omega,u(s))$, $s\in(s_-,s_+)$, $\omega\in\mathbb S^{n-1}$, and we denote its image by $\Sigma(\mathtt c)$.

\medskip

To express the curvature in scalar form, the angle function $\theta$  measured with respect the horizontal axis is defined
by $\dot r=\cos\theta$, $\dot u=\sin\theta$. Then, with this convention, the unit normal along the generating curve is $\vec N(s)=(-\dot u(s),\dot r(s))=(-\sin\theta(s),\cos\theta(s))$. The corresponding unit normal along $\Sigma(\mathtt c)$ has vertical component $\langle\vec\nu,\vec e_{n+1}\rangle=\cos\theta$.

\medskip

Whenever $\mathtt c$ is of class $\mathcal C^2$, its signed curvature is
$\lambda_{\mathtt c}=\dot\theta=\ddot u\dot r-\dot u\ddot r$. The principal curvatures of $\Sigma(\mathtt c)$ are the meridional curvature $\lambda_1=\lambda_{\mathtt c}=\dot\theta$ and the rotational curvature, with multiplicity $n-1$, $\lambda_2=\cdots=\lambda_n=\frac{\sin\theta}{r}$. Consequently, $\Sigma(\mathtt c)$ is a $\gamma$-translator in the direction
$\vec e_{n+1}$ precisely when the profile curve satisfies
\begin{align}
	\label{eq:rotational-translator-arclength}
	\widetilde\gamma
	\left(
	\dot\theta, r^{-1}\sin\theta
   \right)
	=
	\cos\theta,
\end{align}
where $\widetilde\gamma(x,y)=\gamma(x,y,\ldots,y)$ and  $\widetilde\Gamma=\{(x,y)\in\mathbb R^2:(x,y,\ldots,y)\in\Gamma\}$.

\medskip

For later use, we record the two graphical forms of
\eqref{eq:rotational-translator-arclength}.  First, on an interval where
$\dot r\ne0$, the profile curve can be written as a vertical graph
$\mathtt c(r)=(r,u(r))$.  Setting $v(r)=u'(r)$, the principal curvatures are
\begin{align*}
	\lambda_1
	=
	\frac{v'}{(1+v^2)^{\frac{3}{2}}},
	\qquad
	\lambda_2=\cdots=\lambda_n
	=
	\frac{v}{r\sqrt{1+v^2}}.
\end{align*}
Using the $\alpha$-homogeneity of $\gamma$, the translator equation becomes
\begin{align}
	\label{eq:vertical-graph-translator}
	\widetilde\gamma
	\left(
	\frac{v'}{(1+v^2)^{1+\beta}},
	\frac{v}{r(1+v^2)^\beta}
	\right)
	=
	1,
	\qquad
	\beta=\frac{\alpha-1}{2\alpha}.
\end{align}

\medskip

Similarly, on an interval where $\dot u\ne0$, the same curve can be written as a
horizontal graph $\mathtt c(u)=(r(u),u)$.  Using
$r'(u)=u'(r)^{-1}$ and
$r''(u)=-\frac{u''(r)}{u'(r)^3}$, the translator equation takes the form
\begin{align}
	\label{eq:horizontal-graph-translator}
	\widetilde\gamma
	\left(
	\frac{-r''}{(1+(r')^2)^{\beta+1}},
	\frac{1}{r(1+(r')^2)^\beta}
	\right)
	=
	r'.
\end{align}

\section{Curvature functions and phase diagrams}
\label{sec:curvature-functions-branches}

In this section, we introduce the curvature-function classes and the rotational
phase diagrams used throughout the paper.  Recall from the previous section that
level equation
\begin{align}
	\label{sec3:level-set-equation}
	\widetilde\gamma(x,y)=z,
	\qquad
	(x,y)\in\widetilde\Gamma
	=
	\{(x,y)\in\mathbb R^2:(x,y,\ldots,y)\in\Gamma\},
\end{align}
where $x$ denotes the meridional curvature, $y$ denotes the rotational
curvature,  $\widetilde\gamma(x,y)=\gamma(x,y,\ldots,y)$, and $z$ denotes the angle function.

\begin{definition}
	\label{def:alpha-positive-curvature-function}
	Let $\Gamma\subset\mathbb R^n$ be an open symmetric cone containing the
	positive cone
	$\Gamma_+=\{\lambda\in\mathbb R^n:\lambda_i>0\text{ for every }i\}$.
	A function $\gamma:\overline\Gamma\to[0,\infty)$ is called an
	$\alpha$-positive curvature function if:
	\begin{enumerate}[label=\textup{(\alph*)}]
		\item $\gamma$ is symmetric, i.e., invariant under permutations of its
		entries, and continuous on $\overline\Gamma$;
		\item $\gamma\in\mathcal C^1(\Gamma,(0,\infty))$ and it is elliptic on
		$\Gamma$, i.e., $\partial_i\gamma>0$ for every $i$;
		\item $\gamma$ is $\alpha$-homogeneous, i.e.,
		$\gamma(c\lambda)=c^\alpha\gamma(\lambda)$ for every $c>0$ and every
		$\lambda\in\overline\Gamma$.
	\end{enumerate}
\end{definition}

The level $z$, the ellipticity assumptions, and the homogeneity of $\gamma$
lead to different rotational translator theories through
\eqref{sec3:level-set-equation}.  At the level $z=1$, this equation was studied
by Rengaswami in \cite{rengaswami2024classification}, where rotational
bowl-type translators were constructed and classified.  These solutions are
written as rotational graphs $\mathtt c(r)=(r,u(r))$, with either
$u\in\mathcal C^\infty([0,R))$ or
$u\in\mathcal C^1([0,R))\cap\mathcal C^\infty((0,R))$, and with
$R=\infty$ or $R=\widetilde\gamma(1,1)^{-\frac{1}{\alpha}}$.  The value of $R$
is governed by the degeneracy of the curvature function, which we now recall.

\begin{definition}
	\label{def:one-nondegenerate-and-degenerate}
	Let $\gamma:\overline\Gamma\to[0,\infty)$ be an $\alpha$-positive
	curvature function.
	\begin{enumerate}[label=\textup{(\roman*)}]
		\item $\gamma$ is $1$-nondegenerate if $(0,1)\in\widetilde\Gamma$.  In
		this case, we use the normalization $\widetilde\gamma(0,1)=1$.

		\item $\gamma$ is $1$-degenerate if $(0,1)\in\partial\widetilde\Gamma$
		and $\lim\limits_{t\to0^+}\widetilde\gamma(t,1)=0$.
	\end{enumerate}
\end{definition}

This terminology is tied to cylindrical directions. Round cylinders
$\mathbb R\times\mathbb S^{n-1}$ collapse to their axis under the
$\gamma$-flow in the $1$-nondegenerate case, whereas in the $1$-degenerate case
the cylinder is a stationary solution under the $\gamma$-flow.

\medskip

For bowl-type translators, Rengaswami's classification shows that the positive
branch of \eqref{sec3:level-set-equation} determines the global graphical rotational
behavior.  In the normalized $1$-nondegenerate case, the bowl solution is entire and its slope satisfies
$v(r)=r^\alpha(1+o(1))$ as $r\to\infty$
\cite{rengaswami2024classification}.  In the degenerate case, the same branch
formulation separates entire solutions from finite-radius solutions that are asymptotic to a round cylinder according to the behavior of the implicit function $x=x(y)$ from $\widetilde\gamma(x,y)=1$ as $y\to\infty$.  Below, we keep this branch
formulation and develop the sharper asymptotics needed later for the
Alexandrov-reflection principle.

\medskip

Recall that the level $\widetilde\gamma=1$ controls positive graphical ends, while the homogeneous structure also requires positive levels $z>0$.  For this reason,
we record the maximal positive branch $x=g_+(y,z)$ associated with
$\widetilde\gamma(x,y)=z$.  Its level $z=1$ is the branch used in
\cite[Definition~6 and Section~2.3.2]{rengaswami2024classification}, and its
homogeneous extension to arbitrary positive levels will be used below.

\begin{proposition}
	\label{prop:positive-branch}
	Let $\gamma:\overline\Gamma\to[0,\infty)$ be an $\alpha$-positive
	curvature function.  Define
	\begin{align*}
		\mathcal D_+
		=
		\begin{cases}
			\left\{(y,z)\in(0,\infty)^2:
				\frac{z}{\widetilde\gamma(1,1)}<y^\alpha<z\right\},
			& \text{if $\gamma$ is $1$-nondegenerate},\\[6pt]
			\left\{(y,z)\in(0,\infty)^2:
				\frac{z}{\widetilde\gamma(1,1)}<y^\alpha\right\},
			& \text{if $\gamma$ is $1$-degenerate}.
		\end{cases}
	\end{align*}
	Then there exists a unique function
	$g_+\in\mathcal C^1(\mathcal D_+)$ such that
	\begin{align*}
		0<g_+(y,z)<y,
		\qquad
		\widetilde\gamma(g_+(y,z),y)=z
	\end{align*}
	for every $(y,z)\in\mathcal D_+$.  Moreover,
	\begin{align*}
		\partial_y g_+(y,z)
         =
		-\frac{\partial_y\widetilde\gamma(g_+(y,z),y)}
		{\partial_x\widetilde\gamma(g_+(y,z),y)}
		<0,\,
		\partial_z g_+(y,z)
		=
		\frac{1}
		{\partial_x\widetilde\gamma(g_+(y,z),y)}
		>0.
	\end{align*}
	At the level $z=1$, we write $g_+(y)=g_+(y,1)$ and
	\begin{align*}
		I_+
		=
		\begin{cases}
			(y_*,1), & \text{if $\gamma$ is $1$-nondegenerate},\\
			(y_*,\infty), & \text{if $\gamma$ is $1$-degenerate},
		\end{cases}
		\mbox{ with }
		y_*=\widetilde\gamma(1,1)^{-\frac{1}{\alpha}} .
	\end{align*}
\end{proposition}

\begin{proof}
Firstly, we	fix $y>0$. Then,  by ellipticity, we have that the map $x\mapsto\widetilde\gamma(x,y)$ is strictly increasing on the positive sheet. Next, we assume that $\gamma$ is $1$-nondegenerate.  The normalization
	$\widetilde\gamma(0,1)=1$ and homogeneity give
	$\widetilde\gamma(0,y)=y^\alpha$ and
	$\widetilde\gamma(y,y)=\widetilde\gamma(1,1)y^\alpha$.  Hence the level
	$z$ is crossed by a unique $x\in(0,y)$ precisely when $y^\alpha<z<\widetilde\gamma(1,1)y^\alpha$. This gives the stated domain $\mathcal D_+$ in the $1$-nondegenerate case. On the other hand, if  $\gamma$ is $1$-degenerate we have that  the lower endpoint of the positive sheet is reached as a boundary limit.  Therefore, by homogeneity,
	degeneracy, and ellipticity hypotheses, the level $z>0$ is crossed by a unique $x\in(0,y)$ precisely when $0<z<\widetilde\gamma(1,1)y^\alpha$. This gives the stated domain $\mathcal D_+$ in the $1$-degenerate case.

	\medskip
	
	Now we show that the implicit assignment is continuous.  Let
	$(y_j,z_j)\to(y_0,z_0)$ in $\mathcal D_+$, and let $x_j\in(0,y_j)$ be the
	unique solution of $\widetilde\gamma(x_j,y_j)=z_j$.  The inequalities
	$0<x_j<y_j$ and the convergence $y_j\to y_0$ show that the sequence
	$\{x_j\}$ is bounded.  Passing to a subsequence, assume that $x_j\to\bar x$,
	with $0\leq\bar x\leq y_0$. Then, in the $1$-nondegenerate case, continuity gives
	$\widetilde\gamma(\bar x,y_0)=z_0$.  The alternatives $\bar x=0$ and
	$\bar x=y_0$ would imply, respectively,
	$z_0=y_0^\alpha$ and
	$z_0=\widetilde\gamma(1,1)y_0^\alpha$, both contradicting the strict
	inequalities defining $\mathcal D_+$.  Hence $\bar x\in(0,y_0)$. In the $1$-degenerate case, the endpoint $\bar x=y_0$ is excluded in the
	same way.  When $\bar x=0$, then
	\begin{align*}
		x_jy_j^{-1}\to0,
		\qquad
		z_j
		=
		\widetilde\gamma(x_j,y_j)
		=
		y_j^\alpha
		\widetilde\gamma\left(x_jy_j^{-1},1\right)
		\to0,
	\end{align*}
	which contradicts $z_0>0$.  Thus again $\bar x\in(0,y_0)$. Consequently, we have that in both cases, the limit $\bar x$ satisfies
	$\widetilde\gamma(\bar x,y_0)=z_0$.  By uniqueness at $(y_0,z_0)$,
	$\bar x=x_0$, where $x_0$ is the unique solution of
	$\widetilde\gamma(x_0,y_0)=z_0$.  Every convergent subsequence has the
	same limit, so $x_j\to x_0$.  Hence the map $(y,z)\mapsto x$ is continuous on
	$\mathcal D_+$.

	\medskip

	Finally, we observe that the local regularity follows from the implicit function theorem.  Indeed, we take $(y_0,z_0)\in\mathcal D_+$, and let $x_0\in(0,y_0)$ be the unique solution
	of $\widetilde\gamma(x_0,y_0)=z_0$.  The inequality
	$\partial_x\widetilde\gamma(x_0,y_0)>0$ allows the implicit function theorem to give
	a local $\mathcal C^1$ branch $x=g_+(y,z)$ near $(y_0,z_0)$.  By uniqueness,
	these local branches agree on overlaps, and therefore define a unique
	function $g_+\in\mathcal C^1(\mathcal D_+)$ satisfying
	$0<g_+(y,z)<y$ and $\widetilde\gamma(g_+(y,z),y)=z$. Moreover, we notice that differentiating the identity $\widetilde\gamma(g_+(y,z),y)=z$ with respect
	to $y$ and $z$ gives
	\begin{align*}
		\partial_y g_+(y,z)
		=
		-\frac{\partial_y\widetilde\gamma(g_+(y,z),y)}
		{\partial_x\widetilde\gamma(g_+(y,z),y)},
		\qquad
		\partial_z g_+(y,z)
		=
		\frac{1}
		{\partial_x\widetilde\gamma(g_+(y,z),y)}.
	\end{align*}
	The inequalities $\partial_x\widetilde\gamma>0$ and
	$\partial_y\widetilde\gamma>0$ on the positive sheet yield
	$\partial_y g_+<0$ and $\partial_z g_+>0$.

	\medskip

	In particular, restricting to the level $z=1$, we write $g_+(y)=g_+(y,1)$.  Therefore, the condition $(y,1)\in\mathcal D_+$ gives
	\begin{align*}
		&\widetilde\gamma(1,1)^{-1}<y^\alpha<1
		\quad\text{ in the $1$-nondegenerate case},\\
		&\widetilde\gamma(1,1)^{-1}<y^\alpha
		\text{ in the $1$-degenerate case}.
	\end{align*}
	The identity
	$y_*=\widetilde\gamma(1,1)^{-\frac{1}{\alpha}}$ gives the stated
	interval $I_+$.

\end{proof}

The positive branch information needed later is different in the two regimes.
In the $1$-nondegenerate case, the endpoint $y=1$ is regular under the
corresponding regularity of $\widetilde\gamma$.  In the $1$-degenerate case,
the relevant data are the asymptotics of the positive branch as $y\to\infty$.

\begin{remark}
	\label{rem:positive-endpoint-regularity}
	Let $\gamma$ be a $1$-nondegenerate $\alpha$-positive curvature function,
	and assume that $\widetilde\gamma$ is of class $\mathcal C^m$ near $(0,1)$,
	with $m\geq1$.  The conditions $\widetilde\gamma(0,1)=1$ and
	$\partial_x\widetilde\gamma(0,1)>0$ give through the implicit function theorem  and the uniqueness of monotone level equations a $\mathcal C^m$ extension of
	the positive branch to $y=1$ after setting $g_+(1)=0$ since  
   Moreover, by homogeneity,
	\begin{align*}
		g_+'(1)
		=
		-\frac{\partial_y\widetilde\gamma(0,1)}
		{\partial_x\widetilde\gamma(0,1)}
		=
		-\frac{\alpha}{\partial_x\widetilde\gamma(0,1)}.
	\end{align*}
\end{remark}

\begin{definition}
	\label{def:degenerate-positive-branch-asymptotics}
	Let $\gamma$ be a $1$-degenerate $\alpha$-positive curvature function, and
	let $g_+(y)=g_+(y,1)$ be the positive branch.

	\begin{enumerate}[label=\textup{(\roman*)}]
		\item The function $\gamma$ has a power-decay positive branch if there
		are constants $c_\gamma>0$ and $k_\gamma>0$ such that
		\begin{align*}
			g_+(y)
			=
			c_\gamma y^{-k_\gamma}(1+o(1))
			\qquad\text{as }y\to\infty .
		\end{align*}

		\item The function $\gamma$ has a second-order power-decay positive
		branch if it has a power-decay positive branch and there are constants
		$c_{\gamma,1}\in\mathbb R$ and $\sigma_\gamma>0$ such that
		\begin{align*}
			g_+(y)
			=
			c_\gamma y^{-k_\gamma}
			+
			c_{\gamma,1}y^{-k_\gamma-\sigma_\gamma}
			+
			o(y^{-k_\gamma-\sigma_\gamma})
			\qquad\text{as }y\to\infty .
		\end{align*}

		\item The second-order power-decay branch is balanced if
		$\sigma_\gamma=k_\gamma+1$, that is,
		\begin{align*}
			g_+(y)
			=
			c_\gamma y^{-k_\gamma}
			+
			c_{\gamma,1}y^{-(2k_\gamma+1)}
			+
			o(y^{-(2k_\gamma+1)})
			\qquad\text{as }y\to\infty .
		\end{align*}
	\end{enumerate}


\end{definition}

\begin{notation}
	\label{not:asymptotic-notation}
	For real-valued functions $f$ and $h$ defined near a limiting point
	$\ell$, where $\ell$ may be finite or infinite, the notation $f=O(h)$ means
	that $|f(x)h(x)^{-1}$ is bounded in a punctured neighborhood of $\ell$,
	whenever $h(x)$ is nonzero near the limit. The notation $f=o(h)$ means that
	$\lim\limits_{x\to\ell}f(x)h(x)^{-1}=0$. All asymptotic symbols are
	taken with respect to the limit explicitly stated in the corresponding
	result. Endpoint limits are one-sided from within the relevant domain.
\end{notation}


	\subsection{Signed extensions and maximal implicit branches}\,
	\label{subsec:signed-level-sheets}

	In this subsection, we introduce the signed rotational phase structure used in
	the catenoidal construction.  
	
	\medskip
	
	Firstly, we notice that catenoidal-type translators force us to enlarge this positive picture to a signed rotational phase diagram.  Indeed, the neck is encoded by the zero level of the signed profile, $\widehat\gamma(x,y)=0$, with $y>0$ and $x<0$, while the reoriented lower graphical branch is read on the fixed level $\widehat\gamma(x,y)=-1$. Thus the construction requires a signed rotational extension $\widehat\gamma$ of $\widetilde\gamma$, together with a selected component on
	which the levels $0$ and $-1$ can be followed by implicit branches.

	\medskip

	For a catenoidal profile, the neck is represented by $(r,u,\theta)= \left(R,0,\frac{\pi}{2}\right)$, where $R>0$ is the neck radius.  At this point, the signed rotational phase
	equation reads
	\begin{align}
		\label{sec3:neck-zero-level}
		\widehat\gamma\left(\dot\theta(0),R^{-1}\right)=0.
	\end{align}
	In the local horizontal chart, the identity $r''(0)=-\dot\theta(0)$ shows that a
	strict neck corresponds to $\dot\theta(0)<0$, or equivalently $r''(0)>0$.
	Thus the meridional curvature has the opposite sign from the rotational
	curvature, and the zero level must be read on a mixed-sign curvature sheet.

	\medskip

	The definitions below isolate the phase data needed for this construction.
	First, the signed rotational profile $\widehat\gamma$ provides the mixed-sign
	extension of the reduced curvature function.  The selected angular datum
	$(A_g,\rho_0)$ chooses the conical component of the signed phase diagram and
	the zero ray $x=-\rho_0y$.  Once this angular datum is fixed, homogeneity and
	monotonicity produce the maximal implicit chart $x=g(y,z)$ on the selected
	component.  The barrier method in Section~\ref{sec:barrier-method} later
	determines which part of this selected phase diagram is followed by the
	rotational profile.



	\begin{definition}
		\label{def:signed-rotational-extension}
		Let $\gamma:\overline\Gamma\to[0,\infty)$ be an $\alpha$-positive
		curvature function.  A signed rotational level profile for $\gamma$ is a
		real-valued function
		$\widehat\gamma:\widehat\Gamma\to\mathbb R$, defined on an open set
		$\widehat\Gamma\subset\mathbb R^2\setminus\{(0,0)\}$, such that
		$\widehat\gamma=\widetilde\gamma$ on
		$\widehat\Gamma\cap\widetilde\Gamma$.
	\end{definition}

	\begin{definition}
		\label{def:selected-signed-neck-datum}
		Let $\widehat\gamma:\widehat\Gamma\to\mathbb R$ be a signed rotational
		level profile for $\gamma$.  A selected angular neck data consist of an open
		interval $A_g\subset\mathbb R$ and a number $\rho_0>0$ such that
		$-\rho_0\in A_g$ and $\widehat\gamma(-\rho_0,1)=0$.
		After setting $a=xy^{-1}$ for the angular phase variable, the conical
		sheet selected by $A_g$ is
		\begin{align*}
			U_g
			=
			\{(ay,y):a\in A_g,\ y>0\}.
		\end{align*}
		The selected data are required to satisfy:
		\begin{enumerate}[label=\textup{(H\arabic*)}, ref=H\arabic*]
			\item \label{hyp:signed-domain}
			$U_g\subset\widehat\Gamma$ and $-U_g\subset\widehat\Gamma$.


			\item \label{hyp:signed-homogeneity}
			The function $\widehat\gamma$ is positively $\alpha$-homogeneous on
			$U_g\cup(-U_g)$, namely $\widehat\gamma(tx,ty)=t^\alpha\widehat\gamma(x,y)$
			for every $t>0$ and every $(x,y)\in U_g\cup(-U_g)$.

			\item \label{hyp:signed-oddness}
			The function $\widehat\gamma$ is odd under central reflection between
			the two selected sheets, that is  $\widehat\gamma(-x,-y)=-\widehat\gamma(x,y)$ for every $(x,y)\in U_g$.

			\item \label{hyp:signed-angular-monotonicity}
			The angular profile $\Psi_g(a)=\widehat\gamma(a,1)$ with $a\in A_g$ is continuous and strictly increasing on $A_g$.
		\end{enumerate}


	\end{definition}

	\begin{definition}
		\label{def:regular-signed-neck-datum}
		The selected angular neck data from
		Definition~\ref{def:selected-signed-neck-datum} are called regular when there
		exists an open set $\Omega_g\subset\widehat\Gamma$ such that:
		\begin{enumerate}[label=\textup{(R\arabic*)}, ref=R\arabic*]
			\item \label{hyp:regular-neighborhood}
			$U_g\subset\Omega_g$ and $\widehat\gamma$ is of class
			$\mathcal C^1$ on $\Omega_g$.


			\item \label{hyp:regular-ellipticity}
			The ellipticity inequalities $\partial_x\widehat\gamma>0$, $\partial_y\widehat\gamma>0$ hold on $U_g$.
		\end{enumerate}


	\end{definition}

	The zero condition in Definition~\ref{def:selected-signed-neck-datum} fixes the
	angular direction of the neck, while the monotonicity of $\Psi_g$ selects a
	single conical component of the signed phase diagram.  Regularity is the
	ellipticity hypothesis on this selected component; it is precisely what allows
	the signed level equation to be written as an implicit graph $x=g(y,z)$ in the
	next proposition.


	\begin{proposition}
		\label{prop:conical-sheet-maximal-chart}
		Let $(A_g,\rho_0)$ be selected angular neck data, and define
		\begin{align*}
			\mathcal D_g
			=
			\left\{
			(y,z)\in(0,\infty)\times\mathbb R:
			zy^{-\alpha}\in\Psi_g(A_g)
			\right\}.
		\end{align*}
		The selected signed sheet is represented on $\mathcal D_g$ by the
		unique function
		\begin{align*}
			g(y,z)
			=
			y\,\Psi_g^{-1}\!\left(zy^{-\alpha}\right).
		\end{align*}
		Equivalently, $(g(y,z),y)\in U_g$ and
		$\widehat\gamma(g(y,z),y)=z$ for every $(y,z)\in\mathcal D_g$.  The map $g$
		is continuous, and $\mathcal D_g$ is the maximal domain on which the
		selected sheet can be written as a graph over $(y,z)$.  Moreover, when the
		selected angular neck data are $\mathcal C^k$-regular, the function $g$ is
		of class $\mathcal C^k$ locally on $\mathcal D_g$.
	\end{proposition}

	\begin{proof}
	Firstly, we	fix $y>0$ and notice that  every point of the selected sheet $U_g$ has the form $(ay,y)$ with $a\in A_g$. Then,  by Hypothesis~\ref{hyp:signed-homogeneity}, we have $\widehat\gamma(ay,y)=y^\alpha\widehat\gamma(a,1)
		=y^\alpha\Psi_g(a)$.  Hence, we observe that solving
		$\widehat\gamma(x,y)=z$ in $U_g$ is equivalent to solving $\Psi_g(a)=zy^{-\alpha}$ for $a\in A_g$.
		
		\medskip

		Moreover, by Hypothesis~\ref{hyp:signed-angular-monotonicity}, we have that 
		$\Psi_g$ is continuous and strictly increasing on the interval $A_g$.
		Therefore the angular equation has at most one solution, and it has one
		precisely when $zy^{-\alpha}\in\Psi_g(A_g)$.  This gives the domain
		$\mathcal D_g$ and defines the angular solution $a(y,z)\in A_g$, and then we set $g(y,z)=y\,a(y,z)$.
		
		\medskip

		Next, we prove continuity on the map. Indeed,  let $(y_j,z_j)\to(y_0,z_0)$ in $\mathcal D_g$,
		and set $q_j=z_jy_j^{-\alpha}$ and $q_0=z_0y_0^{-\alpha}$.
		Then, we have $q_j\to q_0$, and $q_j,q_0\in\Psi_g(A_g)$. Let $a_j,a_0\in A_g$ be the corresponding angular solutions, namely $\Psi_g(a_j)=q_j$, $\Psi_g(a_0)=q_0$. Therefore,  by fixing $\varepsilon>0$ such that
		$a_0-\varepsilon,a_0+\varepsilon\in A_g$,  we see that the strict monotonicity of $\Psi_g$ gives $\Psi_g(a_0-\varepsilon)	<q_0<\Psi_g(a_0+\varepsilon)$. Then, for $j$ sufficiently large, $q_j$ lies between these two values. Hence
		$a_0-\varepsilon<a_j<a_0+\varepsilon$.  As $\varepsilon$ is arbitrary, we
		obtain $a_j\to a_0$.  In particular, $g(y_j,z_j)=y_j a_j\to y_0 a_0=g(y_0,z_0)$ and $g$ is continuous on $\mathcal D_g$.
		
		\medskip

		Assume now that the selected angular neck data are regular.  By
		Hypothesis~\ref{hyp:regular-neighborhood}, the function $\widehat\gamma$ is
		of class $\mathcal C^1$ on an open neighborhood of $U_g$. Then, we consider
		\begin{align*}
			\mathcal G(a,y,z)
			=
			\Psi_g(a)-zy^{-\alpha},
			\quad
			a\in A_g,
			\quad
			y>0.
		\end{align*}
		Therefore, by hypothesis~\ref{hyp:regular-ellipticity}  we have 
		\begin{align*}
			\partial_a\mathcal G(a,y,z)
			=
			\Psi_g'(a)
			=
			\partial_x\widehat\gamma(a,1)
			>
			0.
		\end{align*}
		We notice that the classical implicit function theorem gives a local $\mathcal C^1$ regularity of the angular solution $a=a(y,z)$, and therefore of
		$g(y,z)=y\,a(y,z)$.  When the data are $\mathcal C^k$-regular, the same
		argument gives local $\mathcal C^k$ regularity. Furthermore, the local implicit solutions agree on overlaps by uniqueness of the angular
		solution.  Finally, when $(y,z)\notin\mathcal D_g$, $zy^{-\alpha}\notin\Psi_g(A_g)$.  Hence no point of the selected sheet
		$U_g$ solves $\widehat\gamma(x,y)=z$.  This proves that $\mathcal D_g$ is
		the maximal domain on which the selected sheet can be represented as a graph
		$x=g(y,z)$ over $(y,z)$.
	\end{proof}

	Below, we write $\mathcal D=\mathcal D_g$ and denote by
	$g:\mathcal D\to\mathbb R$ the selected signed neck chart, so that
	$\widehat\gamma(g(y,z),y)=z$ and $(g(y,z),y)\in U_g$.  Equivalently,
	\begin{align*}
		g(y,z)
		=
		y\,\Psi_g^{-1}\left(zy^{-\alpha}\right),
	\end{align*}
	where the inverse is taken on the selected interval $A_g$.  Moreover, because
	$-\rho_0\in A_g$ and $\Psi_g(-\rho_0)=0$, the zero level belongs to the
	selected chart for every $y>0$.  By the strict monotonicity of $\Psi_g$, we obtain
	\begin{align}
		\label{eq:zero-trace-conical}
		g(y,0)
		=
		-\rho_0 y,
		\mbox{ with } 
		y>0.
	\end{align}
	Thus, when a rotational profile has a neck of radius $R>0$ whose phase is read
	from this selected zero level, we have $\dot\theta(0)=g\left(R^{-1},0\right)=-\rho_0R^{-1}$ and 
	the identities $\theta(0)=\frac{\pi}{2}$ and $\ddot r=-\sin\theta\,\dot\theta$ give $\ddot r(0)=\rho_0R^{-1}>0$.


	\subsection{Selected conical phase data}\,
	\label{subsec:selected-conical-phase-data}

	Next, we record the angular data associated with the selected zero ray
	$x=-\rho_0y$ in the phase plane $(x,y)$. The centered angular coordinate is
	\begin{align*}
		a=-\rho_0+\chi,
		\quad
		\chi=\frac{x+\rho_0 y}{y},
		\quad
		y>0,
	\end{align*}
	and we set $\mathcal X_g=A_g+\rho_0$.  Thus $\chi=0$ corresponds to the selected
	zero direction. In addition, for $y>0$ and $\chi\in\mathcal X_g$, we observe that  homogeneity allow us to write
	\begin{align}
		\label{eq:angular-profile-homogeneity}
		\widehat\gamma((-\rho_0+\chi)y,y)
		=
		y^\alpha\Psi_g(-\rho_0+\chi).
	\end{align}
	Then, by the continuity of $\Psi_g$ on $A_g$ with $-\rho_0\in A_g$, we obtain  that $\widehat\gamma((-\rho_0+\chi)y,y)=O(y^\alpha)$ as $y\to0^+$, uniformly in $\chi$ on compact subsets of the corresponding neighborhood
	of $0$ in $\mathcal X_g$.

	\begin{remark}
		\label{rem:conical-decay-not-origin-continuity}
		The centered variable $\chi$ parametrizes angular rays issuing from the
		selected zero direction. Through the homogeneous relation
		\eqref{eq:angular-profile-homogeneity}, these rays give us the conical level
		structure that is later converted into slope barriers through
		\eqref{eq:conical-slope-barrier} and
		\eqref{eq:explicit-conical-slope-barrier} in the next section. 
	\end{remark}

	\begin{definition}
		\label{def:upper-entry-sector-condition}
		The upper-entry sector is $\mathcal S_+(\rho_0)=\{(ay,y):y>0,\ -\rho_0<a<0\}$.
		The selected angular data satisfy the upper-entry sector condition when
		$\mathcal S_+(\rho_0)\subset U_g$, or equivalently, $(-\rho_0,0)\subset A_g$.
	\end{definition}

	The upper-entry sector condition records that the selected angular component
	contains the positive-side directions between the zero ray $x=-\rho_0y$ and the
	rotational axis $x=0$.  This condition is used with the barrier arguments in
	Section~\ref{sec:barrier-method} and in the construction of the catenoidal
	profiles in Section~\ref{sec:catenoidal-translators}.


	\subsection{Reflection and lower fixed-level components}\,
	\label{subsec:reflected-lower-components}

	Next, we describe the reoriented lower graphical branch and the fixed-level
	components used later in the barrier method.  
	\medskip

	Let $\mathtt c(s)=(r(s),u(s))$ be an arclength parametrized rotational profile
	with $\dot{\mathtt c}(s)=(\cos\theta(s),\sin\theta(s))$.  On the lower side we
	set
	\begin{align*}
		\overline{\mathtt c}(s)=\mathtt c(-s),
		\quad
		\overline\theta(s)=\theta(-s)+\pi .
	\end{align*}
	Then, we have the following identities $\dot{\overline\theta}(s)=-\dot\theta(-s)$,
	$\sin\overline\theta(s)=-\sin\theta(-s)$, and
	$\cos\overline\theta(s)=-\cos\theta(-s)$.  Consequently, the phase variables
	$x=\dot\theta$, $y=\frac{\sin\theta}{r}$, and $z=\cos\theta$ transform as
	$(x,y,z)\mapsto(-x,-y,-z)$.

	\medskip

	In this sense, we  define the reflected chart by
	\begin{align*}
		g^\sharp(y,z)
		=
		-g(-y,-z),
		\mbox{ with }(y,z)\in
		\mathcal D^\sharp
		=
		\{(y,z):(-y,-z)\in\mathcal D\}.
	\end{align*}
	We notice that by the oddness condition, when $(y,z)\in\mathcal D^\sharp$,
	$\widehat\gamma(g^\sharp(y,z),y)=z$, the reflected chart satisfies
	$(g^\sharp(y,z),y)\in -U_g$.  


	\begin{definition}
		\label{def:selected-lower-branch}
		A selected fixed-level lower branch consists of an interval
		$I_-\subset(-\infty,0)$ and a function $g_-:I_-\to\mathbb R$ such that $\widehat\gamma(g_-(y),y)=-1$ for every $y\in I_-$. Then we say:
		\begin{enumerate}[label=\textup{(\roman*)}]
			\item The branch is induced by the reflected chart if
			$g_-(y)=g^\sharp(y,-1)$, for every $y\in I_-\subset\{y:(y,-1)\in\mathcal D^\sharp\}$.

			\item The branch is regular if there exists an open set
			$\Omega_-\subset\widehat\Gamma$ containing
			$\{(g_-(y),y):y\in I_-\}$ such that
			$\widehat\gamma\in\mathcal C^1(\Omega_-)$ and $\partial_x\widehat\gamma>0$, $\partial_y\widehat\gamma>0$ on $\{(g_-(y),y):y\in I_-\}$.
		\end{enumerate}

	\end{definition}

In particular, we observe that	the monotonicity of $g_-(y)$ agrees with its upper branch, since by differentiating the level $z=-1$ equation we have  $g_-'(y)=-\frac{\partial_y\widehat\gamma(g_-(y),y)}{\partial_x\widehat\gamma(g_-(y),y)}<0$.

	\begin{definition}
		\label{def:lower-continuation-outcomes}
		Let $(I_-,g_-)$ be a selected fixed-level lower branch. Then we say:
		\begin{enumerate}[label=\textup{(\roman*)}]
			\item The pair $(I_-,g_-)$ is called a flat complete fixed-level branch if
			$I_-=(\underline y,0)$, with $\underline y<0$, and there are constants
			$\ell<0$ and $\mu>0$ such that  $g_-(y)=\ell y+O(|y|^{1+\mu})\mbox{ as } y\to0^-$.

		
			\item The pair $(I_-,g_-)$ is called a normalized endpoint branch if
			$I_-=(-\infty,-1)$ and $\lim\limits_{y\to-1^-}g_-(y)=0$.

			\item The pair $(I_-,g_-)$ is called a pole-limited branch if
			$I_-=(-1,\lambda_+)$, with $-1<\lambda_+<0$, and $\lim\limits_{y\to\lambda_+^-}\left|g_-(y)y^{-1}\right|=\infty$.
		\end{enumerate}

	\end{definition}

	These alternatives are used as phase alternatives for the lower continuation.
	Their geometric consequences are proved later through the barrier and
	continuation arguments.  The angular formulation also shows that no additional
	endpoint case is needed: blow-up of the angular quotient
	$g_-(y)y^{-1}$ is already encoded by the endpoint behavior of the selected
	angular component and is represented by the pole-limited alternative.


	\subsection{Model examples and verification of the phase data}\,
	\label{subsec:model-examples-branches}

	Finally, we verify the signed phase data in the model families used later.  For each
	family we record the selected angular interval, the angular profile, the zero
	direction, the maximal chart, and the fixed-level lower components.  In the
	quotient examples, the lower component selected after reorientation depends on
	the neck radius through the reoriented lower rotational level
	$y_R=-R^{-1}$.

	\subsubsection{Adjacent Hessian quotients}\,
	\label{subsubsec:adjacent-hessian-quotients}

	Firstly, we consider the adjacent Hessian quotients. Indeed, we  set
	$A_j=\binom{n-1}{j}$ for $j=0,\ldots,n-1$, with $A_{-1}=0$. Then, for a rotational
	curvature vectors, we have $S_j(x,y,\ldots,y)=A_jy^j+A_{j-1}xy^{j-1}$.
	Then, since  $S_j(-\lambda)=(-1)^jS_j(\lambda)$, the adjacent quotient
 $\mathcal Q_{k,k-1}=S_kS_{k-1}^{-1}$ changes sign under central reflection on its real signed branch, and hence it is compatible with the odd reflection condition.

	\medskip

	Next, we write $a=xy^{-1}$.  On the angular component
	the normalized profile is
	\begin{align}
		\label{eq:adjacent-hessian-quotient-angular-profile}
		\Psi_k(a)
		=
		\frac{A_{k-1}}{A_k}
		\frac{(A_k+A_{k-1}a)}{(A_{k-1}+A_{k-2}a)}.
	\end{align}
	Then, by letting $\rho_j=A_jA_{j-1}^{-1}$, the pole corresponds to $a=-\rho_{k-1}$ and the zero ray is $a=-\rho_k$.  Furthermore, by the Newton's inequalities we have  $\rho_k<\rho_{k-1}$, and thus the choice $A_g=(-\rho_{k-1},\infty)$ with $\rho_0=\rho_k$ contains the upper-entry sector $(-\rho_0,0)$.

	\medskip

	On $A_g$, the profile $\Psi_k$ is real analytic and strictly increasing.  Indeed, by the Newton's inequalities we have
	\begin{align*}
		\Psi_k'(a)=\frac{A_{k-1}(A_{k-1}^2-A_kA_{k-2})}
		{A_k(A_{k-1}+A_{k-2}a)^2}>0.
	\end{align*}
 Moreover, the range is $\Psi_k(A_g)=(-\infty,\tau_k^{-1})$,  with $\tau_k=A_kA_{k-2}A_{k-1}^{-2}$. The $1$-homogeneity of $\mathcal Q_{k,k-1}$ gives the maximal selected chart $\mathcal D=\{(y,z):y>0,\ z<\tau_k^{-1}y\}$. Consequently, by solving $\Psi_k(a)=zy^{-1}$ we obtain
	\begin{align*}
		g(y,z)
		=
		\frac{A_kA_{k-1}y(z-y)}
		{A_{k-1}^2y-A_kA_{k-2}z},
	\end{align*}
and, $g(y,0)=-\rho_ky$ as required by \eqref{eq:zero-trace-conical}.  

	\medskip

	Next, the reflected fixed level $z=-1$ is now read from the same chart.  Then, we have by  the formula above that
	\begin{align}
		\label{eq:adjacent-hessian-lower-branch}
		g_-^{(k)}(y)
		=
		-\frac{A_kA_{k-1}y(y+1)}
		{A_{k-1}^2y+A_kA_{k-2}},
		\qquad
		y<0,
		\qquad
		y\ne-\tau_k .
	\end{align}
	Thus the pole separates the fixed level into the components
	\begin{align*}
		(-\infty,-1),
		\qquad
		(-1,-\tau_k),
		\qquad
		(-\tau_k,0),
	\end{align*}
	and the reoriented lower level $y_R=-R^{-1}$ selects the component followed by
	the lower side.  The identity $R_0=\tau_k^{-1}$ gives the following alternatives:
	\begin{itemize}
		\item 	When $R>R_0$, $y_R\in(-\tau_k,0)$, and the lower branch lies in the flat
		complete fixed-level component $I_-^\flat=(-\tau_k,0)$. Then, as $y\to0^-$, \eqref{eq:adjacent-hessian-lower-branch} gives
		\begin{align*}
			g_-^{(k)}(y)
			=
			-\frac{A_{k-1}}{A_{k-2}}y+O(y^2),
			\quad
			\lim\limits_{y\to0^-}
			\frac{g_-^{(k)}(y)}{y}
			=
			-\frac{A_{k-1}}{A_{k-2}}
			=
			-\frac{(n-k+1)}{(k-1)}<0.
		\end{align*}
		Thus this component is a flat complete fixed-level branch in the sense of
		Definition~\ref{def:lower-continuation-outcomes}.
		
		\item When $1<R<R_0$, $y_R\in(-1,-\tau_k)$, and the lower branch lies in the
		pole-limited component $I_-^{\mathrm{pole}}=(-1,-\tau_k)$. The denominator in \eqref{eq:adjacent-hessian-lower-branch} vanishes at
		$y=-\tau_k$, and hence $\lim\limits_{y\to-\tau_k^-}
			\left|g_-^{(k)}y^{-1}\right|=\infty$. This is the pole-limited alternative of Definition~\ref{def:lower-continuation-outcomes}.
		\item When $0<R<1$, $y_R\in(-\infty,-1)$, and the lower branch lies in the
		normalized endpoint component $I_-^{\mathrm{end}}=(-\infty,-1)$, with 
	$\lim\limits_{y\to-1^-}g_-^{(k)}(y)=0$. Thus the same adjacent Hessian quotient phase diagram produces different lower
		continuations according to the neck radius.  
	\end{itemize}
	In each component away from the
	pole, the branch is regular in the sense of
	Definition~\ref{def:selected-lower-branch}; the monotonicity follows from the
	ellipticity of the quotient on the corresponding selected component.


	\subsubsection{Polynomial Hessian speeds}\,

	Next, consider the polynomial speed $S_k$.  For rotational curvature vectors,
	$S_k(x,y,\ldots,y)=A_{k-1}y^{k-1}(x+\rho_ky)$. For odd $k$, this polynomial changes sign under central reflection and gives
	a signed rotational profile satisfying the oddness condition.  The angular
	profile is linear, up to normalization: $\Psi_g(a)=A_{k-1}(a+\rho_k)$. The selected zero direction is $-\rho_0=-\rho_k$.  The selected data are regular
	on the corresponding signed sheet, and the conical decay in
	\eqref{eq:angular-profile-homogeneity} is immediate.  This model therefore
	verifies the selected angular neck data and the zero-trace structure directly.

	\subsubsection{Odd root models}\,

	Next, consider $S_k^{\frac{1}{k}}$ with $k$ odd, using the real-root branch.
	After the normalization $\Psi_g(0)=1$, the angular profile is $\Psi_g(a)=\rho_k^{-\frac{1}{k}}(a+\rho_k)^{\frac{1}{k}}$. The selected zero direction is again $-\rho_0=-\rho_k$.  The profile is
	continuous and strictly increasing, so it defines selected angular neck data.
	For $k>1$, however, it is not of class $\mathcal C^1$ at $a=-\rho_k$.  Thus the
	conical decay statement applies, while the regular selected-data hypotheses at
	the zero ray do not.

	\medskip

	On the reflected fixed level $z=-1$, we select the normalized endpoint component
    $I_-=(-\infty,-1)$, with  $\lim\limits_{y\to-1^-}g_-(y)=0$. Hence the odd-root models give maximal selected-component pieces in this
	fixed-level chart, rather than complete lower ends inside the same component.

	


\section{Barrier method for the first-order equation}
	\label{sec:barrier-method}
	
	In this section, we develop the barrier method for the first-order equations
	arising from rotationally symmetric translators written as vertical graphs
	$F(r,\omega)=(r\omega,u(r))$, with $v(r)=u'(r)$.  Throughout this section, we
	assume $\alpha>\frac{1}{2}$ and set $\beta=\frac{\alpha-1}{2\alpha}\in\left(-\frac{1}{2},\frac{1}{2}\right)$. Furthermore, the rotational phase variables are denoted by $(x,y)$, and  when the phase variables are attached to a graphical slope $v$, we write
	\begin{align*}
		Q(r,s)
		=
		\frac{s}{r(1+s^2)^\beta},
		\quad
		y_v(r)
		=
		Q(r,v(r)),
		\quad
		x_v(r)
		=
		\frac{v'(r)}{(1+v(r)^2)^{\beta+1}}.
	\end{align*}
	Thus $y_v$ records the level of the implicit branch evaluated along the
	solution, while $x_v$ is the corresponding meridional phase coordinate.
	
	\medskip
	
	The role of the barriers is to turn the geometry of the selected branch into
	control of solutions of the graphical equation.  They will be used to keep
	$y_v$ in the relevant branch interval, to identify when a branch ceases to be
	admissible, and to control the approach to the selected conical phase origin.
	
	\medskip
	
	The comparison arguments use the following selected branches.  The positive
	branch is the branch $(I_+,g_+)$ constructed in
	Proposition~\ref{prop:positive-branch}.  For the lower side, we use a regular
	selected fixed-level lower branch $(I_-,g_-)$ in the sense of
	Definition~\ref{def:selected-lower-branch}; thus
	$\widehat\gamma(g_-(y),y)=-1$, with $y\in I_-$.  In the catenoidal
	construction, the reoriented lower side selects a fixed-level component through
	the initial lower level $-R^{-1}$.  When this component
	is induced by the reflected chart $\mathcal D^\sharp$, we have
	$g_-(y)=g^\sharp(y,-1)$.  Otherwise, we work with the selected fixed-level
	branch in the sense of Definition~\ref{def:selected-lower-branch}, and the
	regularity is checked on that component. From now on, $(I_\sigma,g_\sigma)$ denotes either the positive
	branch $(I_+,g_+)$ or a regular selected lower branch $(I_-,g_-)$.  In both
	cases, $g_\sigma$ is strictly decreasing on $I_\sigma$.
	\medskip
	
	The selected graphical equation associated with $(I_\sigma,g_\sigma)$ is
	\begin{align}
		\label{eq:selected-branch-graphical-equation}
		v'
		=
		(1+v^2)^{\beta+1}g_\sigma(Q(r,v)),
		\mbox{ with }
		Q(r,v)\in I_\sigma.
	\end{align}
	Equivalently, in the phase variables introduced above, this equation reads
	$x_v(r)=g_\sigma(y_v(r))$.

	\medskip
	
	\begin{definition}
		\label{def:branch-admissible-functions}
		Let $J\subset(0,\infty)$ and let $v\in\mathcal C^1(J)$.  The function $v$ is
		called
		branch-admissible with respect to $(I_\sigma,g_\sigma)$ on $J$ if $Q(r,v(r))\in I_\sigma\text{ for every }r\in J$.
	\end{definition}
	
	
	The assumption $\alpha>\frac{1}{2}$ is used through the comparison primitive
	$\int_0^s\frac{dt}{(1+t^2)^{\beta+1}}$, which has finite one-sided limits as
	$s\to\pm\infty$ precisely in this range.  The elementary properties of $Q$,
	$y_v$, and the level functions $w_m$ are collected in
	Appendix~\ref{app:q-eta-properties}.  In this section, we use those identities
	to build the comparison theory for
	\eqref{eq:selected-branch-graphical-equation}.
	
	
	\subsection{Residuals and comparison}\,
	\label{subsec:residuals-comparison}
	
	The first tool is the comparison principle for a fixed selected branch.  
	
	\begin{definition}
		\label{def:residual-barriers}
		Let $J\subset(0,\infty)$ and let $w\in\mathcal C^1(J)$ be
		branch-admissible with respect to $(I_\sigma,g_\sigma)$.  The residual is
		defined by
		\begin{align*}
			\mathcal R_\sigma[w]
			=
			w'
			-
			(1+w^2)^{\beta+1}g_\sigma(Q(r,w)).
		\end{align*}
		The function $w$ is called a supersolution, subsolution, strict supersolution,
		or
		strict subsolution according as
		\begin{align*}
			\mathcal R_\sigma[w]\geq0,
			\quad
			\mathcal R_\sigma[w]\leq0,
			\quad
			\mathcal R_\sigma[w]>0,
			\quad
			\mathcal R_\sigma[w]<0,
		\end{align*}
		respectively.
	\end{definition}
	
	\begin{proposition}
		\label{prop:comparison-principle-ode}
		Let $J=[r_0,R)\subset(0,\infty)$, and let
		$v,w\in\mathcal C^1(J)$ be branch-admissible with respect to the same
		regular selected branch $(I_\sigma,g_\sigma)$.  Assume that
		\begin{align*}
			\mathcal R_\sigma[v]\leq0\leq\mathcal R_\sigma[w]
			\qquad
			\text{on }J,
		\end{align*}
		and that $v(r_0)\leq w(r_0)$.  Then $v\leq w$ on $J$.  Moreover, when $v(r_0)<w(r_0)$ and one of the two
		residual inequalities is strict on
		every compact subinterval of $J$, then $v<w$ on $J$.
	\end{proposition}
	
	\begin{proof}
		Set
		\begin{align*}
			d(r)
			=
			\int_{v(r)}^{w(r)}
			\frac{dt}{(1+t^2)^{\beta+1}}.
		\end{align*}
		Because $v,w\in\mathcal C^1(J)$ and the integrand is smooth, $d$ is of class
		$\mathcal C^1$, and
		\begin{align*}
			d'(r)
			=
			\frac{w'(r)}{(1+w(r)^2)^{\beta+1}}
			-
			\frac{v'(r)}{(1+v(r)^2)^{\beta+1}}.
		\end{align*}
		Using the residual inequalities, we obtain
		\begin{align*}
			d'(r)
			\geq
			g_\sigma(Q(r,w(r)))
			-
			g_\sigma(Q(r,v(r))).
		\end{align*}
		
		Arguing by contradiction, suppose that $v(r_1)>w(r_1)$ for some
		$r_1\in J$.  Then $d(r_1)<0$, while $d(r_0)\geq0$.  Define
		\begin{align*}
			r_2
			=
			\sup\{r\in[r_0,r_1]:d(r)=0\}.
		\end{align*}
		By continuity, $d(r_2)=0$.  Because $d(r_1)<0$, we have $r_2<r_1$, and by
		the definition of $r_2$ we have $d<0$ on $(r_2,r_1]$.  The mean value
		theorem gives a point $r_3\in(r_2,r_1)$ such that $d'(r_3)<0$.
		\medskip
		
		At $r_3$ we have $d(r_3)<0$, hence $v(r_3)>w(r_3)$.  Because both functions
		are branch-admissible for the same branch, the levels
		$Q(r_3,v(r_3))$ and $Q(r_3,w(r_3))$ belong to $I_\sigma$.  The strict
		monotonicity of $Q(r_3,\cdot)$ gives
		\begin{align*}
			Q(r_3,v(r_3))
			>
			Q(r_3,w(r_3)).
		\end{align*}
		Because $g_\sigma$ is decreasing on $I_\sigma$, we obtain
		\begin{align*}
			g_\sigma(Q(r_3,w(r_3)))
			-
			g_\sigma(Q(r_3,v(r_3)))
			\geq0.
		\end{align*}
		Therefore $d'(r_3)\geq0$, contradicting $d'(r_3)<0$.  Hence $v\leq w$ on
		$J$.
		\medskip
		
		When the initial inequality is strict and one residual inequality is strict on
		compact subintervals, the same argument excludes an interior first point
		where $d$ vanishes after being positive.  At such a point, the strict
		residual inequality gives $d'>0$, while first vanishing from positive values
		would force $d'\leq0$.  Hence $v<w$ on $J$.
	\end{proof}	
	
\subsection{Level barriers and admissible trapping}\,
\label{subsec:level-barriers}

In the next step we introduce the level barriers associated with the graphical phase
coordinate $y_v=Q(r,v)$.  By Lemma~\ref{lem:level-set-control}, for every
$m\in\mathbb R$ there is a unique smooth function
$w_m:(0,\infty)\to\mathbb R$ satisfying
\begin{align*}
	Q(r,w_m(r))=m .
\end{align*}
Moreover, if $m_1<m_2$, then $w_{m_1}(r)<w_{m_2}(r)$ for every $r>0$, and the
strict monotonicity of $Q(r,\cdot)$ gives
\begin{align*}
	w_{m_1}(r)\leq v(r)\leq w_{m_2}(r)
	\quad
	\Longleftrightarrow
	\quad
	m_1\leq y_v(r)\leq m_2.
\end{align*}
Thus, we have that the level barriers translate comparison for the slope $v$ into control of the
branch level $y_v$.

\medskip

Let $(I_\sigma,g_\sigma)$ be a selected branch.  For $m\in I_\sigma$, the
function
$w_m$ is branch-admissible with respect to $(I_\sigma,g_\sigma)$.  For such
levels, using \eqref{eq:wm-derivative-level} from
Appendix~\ref{app:q-eta-properties} and the identity $Q(r,w_m(r))=m$, we obtain
\begin{align}
	\label{eq:level-barrier-residual-section}
	\mathcal R_\sigma[w_m]
	=
	(1+w_m^2)^{\beta+1}
	\left[
	\frac{m}{1+(1-2\beta)w_m^2}
	-
	g_\sigma(m)
	\right].
\end{align}
This formula is the sign test for deciding whether an admissible level curve is
a subsolution or a supersolution for the selected branch equation.

\begin{proposition}
	\label{prop:level-barrier-trapping}
	Let $v$ solve \eqref{eq:selected-branch-graphical-equation} on
	$J=[r_0,R)$ and assume that $v$ is branch-admissible with respect to
	$(I_\sigma,g_\sigma)$.  Let $a<b$ be levels in $I_\sigma$.  Suppose that
	$w_a$ is a subsolution and $w_b$ is a supersolution for the same branch, and
	that
	\begin{align*}
		w_a(r_0)\leq v(r_0)\leq w_b(r_0).
	\end{align*}
	Then $a\leq y_v(r)\leq b$ for every $r\in J$. In addition, when both residuals are strict and the initial inequalities are strict, then $a<y_v(r)<b$ for every $r\in J$. In particular, if $[a,b]\Subset I_\sigma$, then the solution cannot lose
	branch-admissibility through $\partial I_\sigma$ on $J$.
\end{proposition}

\begin{proof}
	Because $a,b\in I_\sigma$, both level barriers are branch-admissible for the
	same branch as $v$.  By
	Proposition~\ref{prop:comparison-principle-ode}, we have
	$w_a\leq v\leq w_b$ on $J$.  Applying the pointwise equivalence above with
	$m_1=a$ and $m_2=b$, we obtain $a\leq y_v\leq b$ on $J$.  The strict
	statement follows from the strict comparison principle and the strict
	monotonicity of $Q(r,\cdot)$. Finally, in case  $[a,b]\Subset I_\sigma$, the selected branch is evaluated only at
	levels in a compact subset of its domain.  Hence the solution cannot reach
	$\partial I_\sigma$ through its level variable on the comparison interval.
\end{proof}

\begin{remark}
	\label{rem:endpoint-strict-containment}
	We notice that the preceding proposition applies directly to admissible levels
	$m\in I_\sigma$. Moreover, finite endpoint levels will also be used as limiting
	barriers. Indeed,  let $I_\sigma=(a,b)$ and we assume that $a,b\in\mathbb R$ and
	that $g_\sigma$ has one-sided continuous extensions to $a$ and $b$, and that
	the endpoint levels $w_a$ and $w_b$ give, respectively, a strict limiting
	subsolution and a strict limiting supersolution when
	\eqref{eq:level-barrier-residual-section} is read with these endpoint
	values.
	
	\medskip
	
	For the lower endpoint, suppose that $y_v(r_0)\in(a,b)$ and that
	$y_v(r)\in(a,b)$ on $[r_0,\rho)$, with $\lim\limits_{r\to\rho^-}y_v(r)=a$. Choose $\varepsilon>0$ such that
	$0<\varepsilon<y_v(r_0)-a$.  Taking $\varepsilon$ smaller if necessary, the
	residual formula \eqref{eq:level-barrier-residual-section}, the one-sided
	continuity of $g_\sigma$ at $a$, and the dependence of $w_m$ on the level
	parameter from Lemma~\ref{lem:level-set-control} imply that
	$w_{a+\varepsilon}$ is a strict subsolution on the compact interval under
	consideration.  Because $a+\varepsilon\in I_\sigma$, this is an admissible
	barrier.  Moreover,
	\begin{align*}
		Q(r_0,w_{a+\varepsilon}(r_0))
		=
		a+\varepsilon
		<
		y_v(r_0)
		=
		Q(r_0,v(r_0)),
	\end{align*}
	and hence $w_{a+\varepsilon}(r_0)<v(r_0)$.  The strict comparison principle
	gives $w_{a+\varepsilon}<v$ on $[r_0,\rho)$.  Applying $Q(r,\cdot)$, we get
	\begin{align*}
		a+\varepsilon
		=
		Q(r,w_{a+\varepsilon}(r))
		<
		Q(r,v(r))
		=
		y_v(r)
		\qquad
		\text{on }[r_0,\rho).
	\end{align*}
	Letting $r\to\rho^-$ gives $a+\varepsilon\leq a$, a contradiction.
	
	\medskip
	
	The upper endpoint is treated in the same way, using the interior levels
	$b-\varepsilon$.
\end{remark}

In particular, for the normalized $1$-nondegenerate positive branch
$I_+=(y_*,1)$, the endpoint values are $g_+(y_*)=y_*$ and $g_+(1)=0$.  Formula
\eqref{eq:level-barrier-residual-section} gives
\begin{align*}
	\mathcal R_+[w_{y_*}]
	&=
	(1+w_{y_*}^2)^{\beta+1}
	y_*\left[
	\frac{1}{1+(1-2\beta)w_{y_*}^2}
	-
	1
	\right]
	<0,
	\\
	\mathcal R_+[w_1]
	&=
	\frac{(1+w_1^2)^{\beta+1}}
	{1+(1-2\beta)w_1^2}
	>
	0.
\end{align*}
Thus the endpoint levels are strict limiting barriers, and
Remark~\ref{rem:endpoint-strict-containment} gives strict containment in
$I_+$.
	
	
\subsection{Level barriers on lower fixed-level components}\,
\label{subsec:level-barriers-negative-components}

The lower fixed-level components selected in
Definition~\ref{def:lower-continuation-outcomes} are the first application of
these level barriers.  After reorientation, the
lower graphical equation is read on the fixed level $z=-1$, and the solution is
admissible only while its level variable remains inside the selected interval
$I_-$.  Thus the first question is not yet the geometry of the lower end, but
whether the solution can remain in that component.

\medskip

For a normalized endpoint component $I_-=(-\infty,-1)$, the endpoint $-1$ is
reached in finite radius.  This is detected by a constant barrier.  For
admissible negative levels inside a selected lower component, the level barriers
give one-sided control.  Flat complete components require a different scale
near the endpoint $0$, and will be treated by power barriers in the next
subsection.

\begin{lemma}
	\label{lem:constant-barrier-normalized-endpoint-branch}
	Let $(I_-,g_-)$ be a regular normalized endpoint component in the sense of
	Definition~\ref{def:lower-continuation-outcomes}; that is,
	$I_-=(-\infty,-1)$ and
	$\lim\limits_{y\to-1^-}g_-(y)=0$.  Let $v$ be a solution of
	\eqref{eq:selected-branch-graphical-equation} on an interval $[r_0,T)$ on
	which it is branch-admissible for $I_-$.  Suppose that
	$v(r_0)=v_0<0$ and $Q(r_0,v_0)=-m$, with $m>1$.  Set
	$r_*=mr_0$.  Then
	\begin{align*}
		Q(r,v(r))>-mr_0r^{-1}
		\qquad
		\text{for every }r\in(r_0,\min\{T,r_*\}).
	\end{align*}
	In particular, the solution cannot remain admissible in the component
	$I_-=(-\infty,-1)$ on an interval containing $r_*$.
\end{lemma}

\begin{proof}
	First we record the sign of the branch on this component.  Indeed, since the selected
	lower branch is regular, we have $g_-'(y)<0$ on $I_-$.  Together with
	$\lim\limits_{y\to-1^-}g_-(y)=0$, this gives
	$g_-(y)>0$ for every $y\in(-\infty,-1)$.
	
	\medskip
	
	Next, we set $\underline w(r)=v_0$.  Then, by the level barrier $Q(r_0,v_0)=-m$, the radial scaling of
	$Q$ gives $Q(r,\underline w(r))=-mr_0r^{-1}$. Thus, for $r\in[r_0,r_*)$, we have
	$Q(r,\underline w(r))\in(-\infty,-1)$, and hence $\underline w$ is
	branch-admissible for the same lower component on this interval.  Moreover,
	since $g_-(y)>0$ on $(-\infty,-1)$, we obtain $\mathcal R_-[\underline w](r)=-(1+\underline w^2)^{\beta+1}g_-(Q(r,\underline w(r)))<0$ for $r\in[r_0,r_*)$.  Hence $\underline w$ is a strict subsolution there.
	
	\medskip
	
	Then, since  $v(r_0)=\underline w(r_0)$, we have that  the strict comparison principle gives $\underline w(r)<v(r)$ for $r\in(r_0,\min\{T,r_*\})$. Consequently, by applying the strict monotonicity of $Q(r,\cdot)$, we get
   $Q(r,v(r))>Q(r,\underline w(r))=-mr_0r^{-1}$ on the same interval.
	
	\medskip
	
	Finally, if the solution remained admissible on an interval containing
	$r_*$, then by continuity we would have $Q(r_*,v(r_*))\geq-1$, while admissibility in $I_-=(-\infty,-1)$ requires
	$Q(r_*,v(r_*))<-1$.  This contradiction proves the claim.
\end{proof}

\begin{proposition}
	\label{prop:level-barriers-signed}
	Let $(I_-,g_-)$ be either a regular normalized endpoint component or a
	regular flat complete fixed-level component in the sense of
	Definition~\ref{def:lower-continuation-outcomes}.  Let
	$\overline m\in I_-$.  Then the level function $w_{\overline m}$ is
	branch-admissible for the selected lower branch and is a strict subsolution.
\end{proposition}

\begin{proof}
	Because $\overline m\in I_-$, the level function is branch-admissible for the
	selected lower branch.  In both cases under consideration, regularity gives
	$g_-'(y)<0$, and the right endpoint of the component has limiting value
	$0$.  Therefore $g_-(y)>0$ for every $y\in I_-$.  Then, by
	\eqref{eq:level-barrier-residual-section} together with $\overline m<0$, $1-2\beta=\alpha^{-1}>0$, and
	$g_-(\overline m)>0$, we have
	\begin{align*}
		\mathcal R_-[w_{\overline m}]
		=
		(1+w_{\overline m}^2)^{\beta+1}
		\left[
		\frac{\overline m}
		{1+(1-2\beta)w_{\overline m}^2}
		-
		g_-(\overline m)
		\right]<0,
	\end{align*}
 and so $w_{\overline m}$ is a strict subsolution.
\end{proof}

	
	\subsection{Power barriers on flat lower components}\,
	\label{subsec:power-barriers-flat-components}
	
	Flat complete lower components require a separate treatment.  In this case the
right endpoint
	of the selected interval is $0$, and fixed negative levels do not capture the
	asymptotic scale of the solution near that endpoint.  Therefore, we use power
	barriers, whose levels approach $0^-$ as $r\to\infty$.
	
	Let $(I_-,g_-)$ be a regular flat complete fixed-level component in the sense
	of Definition~\ref{def:lower-continuation-outcomes}; that is,
	$I_-=(\underline y,0)$ and, for some $\ell<0$ and $\mu>0$,
	\begin{align*}
		g_-(y)
		=
		\ell y+O(|y|^{1+\mu})
		\qquad
		\text{as }y\to0^-.
	\end{align*}
	
	\begin{proposition}
		\label{prop:power-barrier-flat-component}
		Let $\overline w(r)=-ar^p$, with $a>0$ and $p<0$.  Then
		$\overline w$ is branch-admissible for $(I_-,g_-)$ for all sufficiently
		large $r$.  More precisely, if
		\begin{align*}
			t(r)
			=
			ar^{p-1}(1+a^2r^{2p})^{-\beta},
		\end{align*}
		then $Q(r,\overline w(r))=-t(r)$, with $t(r)\to0^+$ as $r\to\infty$, and
		for all sufficiently large $r$,
		\begin{align}
			\label{eq:power-barrier-residual}
			\mathcal R_-[\overline w]
			=
			-ar^{p-1}
			\left[
			p
			-
			(1+\overline w^2)
			\frac{g_-(-t(r))}{-t(r)}
			\right].
		\end{align}
		Consequently, when $p<\ell<0$, $\overline w$ is a strict supersolution
		for all sufficiently large $r$.  When $\ell<p<0$, $\overline w$ is a
		strict subsolution for all sufficiently large $r$.
	\end{proposition}
	
	\begin{proof}
		Admissibility comes first.  Because $p<0$, we have
		$\overline w(r)\to0^-$ as $r\to\infty$.  Moreover,
		\begin{align*}
			Q(r,\overline w(r))
			=
			\frac{-ar^p}
			{r(1+a^2r^{2p})^\beta}
			=
			-ar^{p-1}(1+a^2r^{2p})^{-\beta}
			=
			-t(r).
		\end{align*}
		As $p-1<0$, we have $t(r)\to0^+$.  Hence
		$Q(r,\overline w(r))=-t(r)\to0^-$.  Because
		$I_-=(\underline y,0)$, it follows that
		$Q(r,\overline w(r))\in I_-$ for all sufficiently large $r$.  Thus
		$\overline w$ is branch-admissible for the selected lower branch on every
		sufficiently large comparison interval.
		
		\medskip
		
		It remains to compute the residual. For this purpose we compute $\overline w'(r)=-ap r^{p-1}$ and $Q(r,\overline w(r))=-t(r)$. Then, the residual turns to
		\begin{align*}
			\mathcal R_-[\overline w]
			=
			-ar^{p-1}
			\left[
			p
			-
			(1+\overline w^2)
			\frac{g_-(-t(r))}{-t(r)}
			\right],
		\end{align*}
		which proves \eqref{eq:power-barrier-residual}.
		\medskip
		
		Finally, because $g_-(y)=\ell y+O(|y|^{1+\mu})$ as $y\to0^-$ and
		$t(r)\to0^+$, we have
		\begin{align*}
			\frac{g_-(-t(r))}{-t(r)}
			\to
			\ell.
		\end{align*}
		Also, $\overline w(r)\to0$, so $1+\overline w^2\to1$.  Hence the bracket in
		\eqref{eq:power-barrier-residual} tends to $p-\ell$.  Because
		$-ar^{p-1}<0$, the residual has the opposite sign of $p-\ell$ for all
		sufficiently large $r$.  Therefore, when $p<\ell<0$,
		$\mathcal R_-[\overline w]>0$ for all sufficiently large $r$, and
		$\overline w$ is a strict supersolution.  When $\ell<p<0$,
		$\mathcal R_-[\overline w]<0$ for all sufficiently large $r$, and
		$\overline w$ is a strict subsolution.
	\end{proof}	
	
	\subsection{Conical phase barriers}\,
	\label{subsec:conical-phase-barriers}
	
In this subsection, we develop a complementary form of control to the level
barriers introduced in Subsection~\ref{subsec:level-barriers}.  The level
barriers constrain the phase level $y_v$ to remain inside a prescribed branch
interval, but they do not determine the direction in which the phase point
$(x_v,y_v)$ approaches the origin in the phase plane.  For the lower
continuation, this directional information is essential: the phase point must
approach the origin along the selected zero ray $x=-\rho_0 y$.  To obtain this
control, we introduce conical phase barriers.  These barriers restrict the phase
dynamics to prescribed angular regions and force convergence toward the
selected direction.
	
	\medskip
	
	The barriers that we obtain in this subsection are by prescribing nearby conical rays in the phase plane and translating them into slope functions. For this purpose we consider $\chi\in\mathcal X_g=A_g+\rho_0$, and set $\lambda_\chi=-\rho_0+\chi$. Then, the corresponding conical ray is $x=\lambda_\chi y$.  A conical slope barrier
	with parameter $\chi$ is a $\mathcal C^1$ function $\varphi_\chi$ satisfying
	\begin{align}
		\label{eq:conical-slope-barrier}
		x_{\varphi_\chi}
		=
		\lambda_\chi y_{\varphi_\chi}
		\Longleftrightarrow
		\varphi_\chi'
		=
		\lambda_\chi
		\frac{\varphi_\chi(1+\varphi_\chi^2)}{r}.
	\end{align}
	On every interval where $\varphi_\chi$ has fixed sign, integration gives
	\begin{align}
		\label{eq:explicit-conical-slope-barrier}
		\varphi_{\chi,\varepsilon}(r)
		=
		\varepsilon
		\frac{C r^{\lambda_\chi}}
		{\sqrt{1-C^2r^{2\lambda_\chi}}},
		\mbox{ with }
		\varepsilon\in\{-1,1\},
	\end{align}
	whenever $0<C^2r^{2\lambda_\chi}<1$.  The parameter $\chi$ is the centered
	angular coordinate of the conical barrier.  Indeed, wherever
	$y_{\varphi_\chi}\ne0$,
	\begin{align*}
		\frac{x_{\varphi_\chi}+\rho_0y_{\varphi_\chi}}
		{y_{\varphi_\chi}}
		=
		\chi.
	\end{align*}
	

	
	The residual test is the following.  Let $(I_\sigma,g_\sigma)$ be a selected
	branch, and restrict the comparison interval so that
	$y_{\varphi_\chi}(r)\in I_\sigma$.  On this interval the conical barrier is
	branch-admissible for $(I_\sigma,g_\sigma)$.  Because
	$x_{\varphi_\chi}=\lambda_\chi y_{\varphi_\chi}$, we obtain
	\begin{align}
		\label{eq:conical-barrier-residual}
		\mathcal R_\sigma[\varphi_\chi]
		=
		(1+\varphi_\chi^2)^{\beta+1}
		\left[
		\lambda_\chi y_{\varphi_\chi}
		-
		g_\sigma(y_{\varphi_\chi})
		\right].
	\end{align}
	Thus the sign of the residual is determined by the position of the branch graph
	$x=g_\sigma(y)$ relative to the ray $x=\lambda_\chi y$ on the admissible range
	of levels.  When $\lambda_\chi y-g_\sigma(y)>0$ on the relevant admissible level interval, $\varphi_\chi$ is a strict supersolution there.  When the opposite strict inequality holds, $\varphi_\chi$ is a strict subsolution.
	
	\medskip
	
	In Section~\ref{sec:catenoidal-translators}, the level barriers will first be
	used to ensure that the solution and the conical barriers are evaluated in the
	same selected branch interval.  The sign condition in
	\eqref{eq:conical-barrier-residual} then provides admissible conical
	subsolutions and supersolutions.  The comparison principle traps the phase
	point between the corresponding rays, and letting the angular parameters tend
	to $0$ gives the selected zero-origin phase condition, namely convergence to
	the ray $x=-\rho_0y$ in the phase plane.
	
	
\subsection{The selected zero-origin phase condition}\,
\label{subsec:zero-origin-phase}

This subsection records the phase condition associated with a return to the selected
conical phase origin.  In the abstract phase plane, this means convergence to
the origin with direction asymptotic to the selected ray $x=-\rho_0y$.  Along a
graphical solution $v$, this condition is expressed through the graphical phase
variables $(x_v,y_v)$.

\begin{definition}
	\label{def:zero-origin-return}
	Let $v$ be a reoriented graphical solution on an interval
	$J\subset(0,\infty)$, and let $r_*\in\overline J$ be an endpoint.  The
	solution $v$ satisfies the selected zero-origin phase condition at $r_*$ if, as
	$r\to r_*$ from within $J$,
	\begin{align*}
		x_v(r)\to0,
		\qquad
		y_v(r)\to0,
	\end{align*}
	and, along the set where $y_v(r)\ne0$,
	\begin{align*}
		\frac{x_v(r)+\rho_0y_v(r)}
		{y_v(r)}
		\to0.
	\end{align*}
\end{definition}

The last condition says that the centered angular coordinate of the phase point
tends to zero.  Equivalently, the phase point approaches $(0,0)$ with direction
asymptotic to the selected ray $x=-\rho_0y$.  In the continuation arguments
below, this condition will be verified by conical barriers: for each small
$\varepsilon>0$, one traps the phase point between the two rays
\begin{align*}
	x=(-\rho_0-\varepsilon)y,
	\qquad
	x=(-\rho_0+\varepsilon)y,
\end{align*}
on the relevant admissible phase interval, and then lets
$\varepsilon\to0^+$.

	\section{Asymptotics of bowl-type solutions}
	\label{sec:positive-branch-asymptotics}
	
	In this section, we derive the positive-end asymptotics for rotational
	bowl-type solutions.  Throughout the section, we use the positive branch
	$(I_+,g_+)$ constructed in Proposition~\ref{prop:positive-branch}, and we
	write the positive graphical equation as
	\begin{align}
		\label{eq:positive-slope-equation-sec5}
		v'
		=
		(1+v^2)^{\beta+1}g_+(Q(r,v)),
		\qquad
		Q(r,v)\in I_+.
	\end{align}
	A solution of \eqref{eq:positive-slope-equation-sec5} is called admissible
	when it is branch-admissible with respect to $(I_+,g_+)$ in the sense of
	Definition~\ref{def:branch-admissible-functions}.  Throughout this section, we
	assume $\alpha>\frac{1}{2}$, so that the comparison and barrier results from
	Section~\ref{sec:barrier-method} apply.
	
	\medskip
	
	Here we use the barriers to identify which part of the positive branch is sampled by
	the graphical solution.  In the normalized $1$-nondegenerate case, they keep the
	phase level $y_v=Q(r,v)$ inside the interval $(y_*,1)$ and force it to approach
	the cylindrical endpoint $y=1$.  In the $1$-degenerate case, they show that the
	same phase level escapes to the infinite end of the positive branch.  Once this
	branch information is known, we insert the corresponding expansion of $g_+$
	into the graphical equation and extract the slope asymptotics.
	
	\medskip
	
	At this point, we distinguish two regimes.  In the normalized $1$-nondegenerate case,
	$I_+=(y_*,1)$, with
	$y_*=\widetilde\gamma(1,1)^{-\frac{1}{\alpha}}$, and $g_+(1)=0$.  The endpoint
	expansion of $g_+$ gives the successive corrections to the leading scale
	$v(r)=r^\alpha(1+o(1))$.  In the $1$-degenerate case,
	$I_+=(y_*,\infty)$, and the decay of $g_+(y)$ as $y\to\infty$ determines the
	growth of the positive end.  The elementary identities for $y_v$ used below,
	including \eqref{eq:eta-derivative-branch},
	\eqref{eq:eta-large-slope-asymptotic}, and
	\eqref{eq:eta-large-slope-second-order}, are collected in
	Appendix~\ref{app:q-eta-properties}.
	
	\subsection{Statement of the positive-end asymptotics}\,
	\label{subsec:positive-end-asymptotics-statement}
	
	Here we collect the asymptotic alternatives proved in this section.
	
	\begin{theorem}
		\label{thm:positive-bowl-asymptotics}
		Let $v$ be the positive slope of a bowl-type solution of
		\eqref{eq:positive-slope-equation-sec5}.
		
		\begin{enumerate}[label=\textup{(\roman*)}]
			\item Assume that $\gamma$ is normalized $1$-nondegenerate and that the
			positive branch has the endpoint regularity required in
			Lemma~\ref{lem:uniform-expansions-nondegenerate}.  Write the endpoint
			expansion of the positive branch as
			\begin{align*}
				g_+(y)
				=
				a_1(1-y)
				+
				a_2(1-y)^2
				+
				o((1-y)^2)\mbox{ as }y\to1^-,
			\end{align*}
			where $a_1=-g_+'(1)>0$ and $a_2=\frac{1}{2}g_+''(1)$.  Define $a= \alpha\left(\alpha a_1^{-1}-\beta\right)$, $L_1=-a\alpha^{-1}-\beta$,
			and
			\begin{align*}
				\mathfrak P(a,\beta)
				=
				\frac{1}{2}(1-2a)\beta
				\bigl(1+\beta(1-2a)\bigr)
				-
				a^2\beta.
			\end{align*}
			Then
			\begin{align*}
				b
				=
				\alpha
				\left[
				-\frac{a\alpha}{a_1}
				-
				\mathfrak P(a,\beta)
				-
				(\beta+1)(1-2a)L_1
				+
				\frac{a_2}{a_1}L_1^2
				\right],
			\end{align*}
			and
			\begin{align*}
				v(r)
				=
				r^\alpha
				-
				\frac{a}{r^\alpha}
				+
				\frac{b}{r^{3\alpha}}
				+
				o(r^{-4\alpha})
				\qquad
				\text{as }r\to\infty.
			\end{align*}
			
		\item Assume that $\gamma$ is $1$-degenerate and that $g_+$ has a
		power-decay positive branch in the sense of
		Definition~\ref{def:degenerate-positive-branch-asymptotics}, namely
		\begin{align*}
			g_+(y)
			=
			c_\gamma y^{-k_\gamma}(1+o(1))
			\qquad
			\text{as }y\to\infty,
		\end{align*}
		with $c_\gamma>0$ and $k_\gamma>0$.  For $k_\gamma>2\alpha-1$, one has
		\begin{align*}
			v(r)
			=
			A_\gamma r^{d_\gamma}(1+o(1))
			\qquad
			\text{as }r\to\infty,
		\end{align*}
		where
		\begin{align*}
			d_\gamma
			=
			\frac{\alpha(k_\gamma+1)}{k_\gamma+1-2\alpha},
			\qquad
			A_\gamma^{\frac{k_\gamma+1-2\alpha}{\alpha}}
			=
			\frac{c_\gamma}{d_\gamma}.
		\end{align*}
		
		\item Under the hypotheses of \textup{(ii)}, if
		$k_\gamma=2\alpha-1$, then
		\begin{align*}
			\log v(r)
			=
			\frac{c_\gamma}{2\alpha}r^{2\alpha}
			+
			o(r^{2\alpha})
			\qquad
			\text{as }r\to\infty.
		\end{align*}
		
		\item Assume, in addition to \textup{(ii)} with
		$k_\gamma>2\alpha-1$, that the positive branch has balanced second-order
		power decay in the sense of
		Definition~\ref{def:degenerate-positive-branch-asymptotics}, that is,
		\begin{align*}
			g_+(y)
			=
			c_\gamma y^{-k_\gamma}
			+
			c_{\gamma,1}y^{-(2k_\gamma+1)}
			+
			o\left(y^{-(2k_\gamma+1)}\right)
			\qquad
			\text{as }y\to\infty.
		\end{align*}
		Set
		\begin{align*}
			\mathcal C_\gamma
			=
			\bigl(1+\beta(k_\gamma+1)\bigr)A_\gamma^{-2}
			+
			\frac{c_{\gamma,1}}{c_\gamma}
			A_\gamma^{-\frac{k_\gamma+1}{\alpha}}.
		\end{align*}
		For $k_\gamma+1>4\alpha$, one has
		\begin{align*}
			v(r)
			=
			A_\gamma r^{d_\gamma}
			+
			B_\gamma r^{-d_\gamma}
			+
			o(r^{-d_\gamma}),
			\qquad
			B_\gamma
			=
			\frac{\alpha A_\gamma\mathcal C_\gamma}
			{k_\gamma+1-4\alpha}.
		\end{align*}
		For $k_\gamma+1=4\alpha$, one has
		\begin{align*}
			v(r)
			=
			A_\gamma r^{d_\gamma}
			+
			L_\gamma r^{-d_\gamma}\log r
			+
			o(r^{-d_\gamma}\log r),
			\qquad
			L_\gamma
			=
			A_\gamma d_\gamma\mathcal C_\gamma.
		\end{align*}
		
		\end{enumerate}
	\end{theorem}
	
	\begin{remark}
		\label{rem:balanced-second-order-scale}
		The exponent $\sigma_\gamma=k_\gamma+1$ in the balanced second-order
		assumption is the scale at which the second term in the branch expansion
		and the first correction in the large-slope expansion of the phase variable
		contribute to the same order in the slope equation.  This is why the
		coefficient $\mathcal C_\gamma$ contains both the geometric correction
		coming from $Q(r,v)$ and the branch coefficient $c_{\gamma,1}$.  Other
		second-order powers in Definition~\ref{def:degenerate-positive-branch-asymptotics}
		lead to different correction scales and are not used in the second-order
		statement above.
	\end{remark}
	
	
	
	\begin{remark}
		\label{rem:nondegenerate-coefficients-gamma-derivatives}
		Assume that $\widetilde\gamma$ is of class $\mathcal C^2$ near $(0,1)$.
		In the normalized $1$-nondegenerate case, $\widetilde\gamma(0,1)=1$,
		the positive branch satisfies $g_+(1)=0$. Differentiating once gives
		$\partial_x\widetilde\gamma(0,1)g_+'(1)+\partial_y\widetilde\gamma(0,1)=0$.Using homogeneity at $(0,1)$, we have
		$\partial_y\widetilde\gamma(0,1)=\alpha$.  Hence
		\begin{align*}
			g_+'(1)
			=
			-\frac{\partial_y\widetilde\gamma(0,1)}
			{\partial_x\widetilde\gamma(0,1)}
			=
			-\frac{\alpha}{\partial_x\widetilde\gamma(0,1)},
		\end{align*}
		and therefore, we obtain  $a_1=-g_+'(1)=\alpha\partial_x\widetilde\gamma(0,1)^{-1}$. Differentiating twice gives
		\begin{align*}
			\partial_x\widetilde\gamma(0,1)g_+''(1)
			+
			\partial_x^2\widetilde\gamma(0,1)(g_+'(1))^2
			+
			2\partial_{xy}^2\widetilde\gamma(0,1)g_+'(1)
			+
			\partial_y^2\widetilde\gamma(0,1)
			=
			0.
		\end{align*}
		Thus, we have 
		\begin{align*}
			a_2
			=
			-\frac{1}{2\partial_x\widetilde\gamma(0,1)}
			\left[
			\partial_x^2\widetilde\gamma(0,1)(g_+'(1))^2
			+
			2\partial_{xy}^2\widetilde\gamma(0,1)g_+'(1)
			+
			\partial_y^2\widetilde\gamma(0,1)
			\right].
		\end{align*}
		We notice that homogeneity also gives $\partial_y^2\widetilde\gamma(0,1)=\alpha(\alpha-1)$ and  $\partial_{xy}^2\widetilde\gamma(0,1)=(\alpha-1)\partial_x\widetilde\gamma(0,1)$. Therefore the coefficients $a$ and $b$ in
		Theorem~\ref{thm:positive-bowl-asymptotics}\textup{(i)} may be written
		entirely in terms of the derivatives of $\widetilde\gamma$ at $(0,1)$.
	\end{remark}
	
	\begin{remark}
		\label{rem:normalized-mean-curvature-check}
		Let $\alpha=1$ and consider the normalized mean-curvature speed
		$\gamma=(n-1)^{-1}H$. Then $\widetilde\gamma(x,y)=y+(n-1)^{-1}x$, and  the positive branc is $	g_+(y)=(n-1)(1-y).$. Thus, in the notation of
		Theorem~\ref{thm:positive-bowl-asymptotics}\textup{(i)}, we have
		$a_1=n-1$, $a_2=0$, and $\beta=0$.  Consequently $a=(n-1)^{-1}$ and  $b=(n-4)(n-1)^{-2}$. Hence the expansion becomes
		\begin{align*}
			v(r)
			=
			r
			-
			\frac{1}{(n-1)r}
			+
			\frac{n-4}{(n-1)^2r^3}
			+
			o(r^{-4}).
		\end{align*}
		This is the normalized form of the bowl-soliton expansion obtained by
		Clutterbuck--Schn\"urer--Schulze \cite{clutterbuck2007stability}.  Indeed,
		their standard mean-curvature-flow normalization gives
		\begin{align*}
			v_{\mathrm{CSS}}(r)
			=
			\frac{r}{n-1}
			-
			\frac{1}{r}
			+
			\frac{(n-1)(n-4)}{r^3}
			+
			o(r^{-4}),
		\end{align*}
		and the change from the speed $H$ to the normalized speed
		$\frac{1}{n-1}H$ corresponds to the scaling
		$v(r)=v_{\mathrm{CSS}}((n-1)r)$.
	\end{remark}
	
	
	\subsection{The nondegenerate endpoint}\,
	\label{subsec:one-nondegenerate-asymptotics}
	
	In the first regime, we start with the normalized $1$-nondegenerate case.  Here
	$I_+=(y_*,1)$, with $y_*=\widetilde\gamma(1,1)^{-\frac{1}{\alpha}}$, and
	$g_+$ has the endpoint value $g_+(1)=0$.  The role of the barriers is to keep
	the positive solution strictly inside this interval and then force the level
	variable to approach the cylindrical endpoint $1$.
	
	\begin{proposition}
		\label{prop:positive-end-leading-order}
		Assume that $\gamma$ is normalized $1$-nondegenerate.  Let $v$ solve
		\eqref{eq:positive-slope-equation-sec5} on its maximal positive-branch
		admissible interval $[r_0,T)$, and assume that
		\begin{align*}
			v(r_0)>0,
			\qquad
			y_v(r_0)=Q(r_0,v(r_0))\in(y_*,1).
		\end{align*}
		Then $T=\infty$, $y_*<y_v(r)<1$ for every $r\ge r_0$, and $y_v(r)\to1$ with $v(r)=r^\alpha(1+o(1))$ as $r\to\infty$.
	\end{proposition}
	
	\begin{proof}
		The endpoint barriers first give the branch containment.  Since
		$y_v(r_0)\in(y_*,1)$, the strict endpoint-containment statement in
		Remark~\ref{rem:endpoint-strict-containment} gives $y_*<y_v(r)<1$
		on every interval where the positive branch solution is defined.
		
		\medskip
		
		The maximal admissible interval is unbounded.  Suppose, by
		contradiction, that $T<\infty$.  Applying
		Remark~\ref{rem:endpoint-strict-containment} on $[r_0,T)$, the trajectory stays
		a positive distance away from the endpoint levels as $r\to T^-$.  Hence there
		exists $\delta>0$ such that
		\begin{align*}
			y_*+\delta
			\le
			y_v(r)
			\le
			1-\delta
			\qquad
			\text{for }r\in[r_0,T).
		\end{align*}
		The map $Q$ is strictly increasing in the slope variable on the positive
		branch.  Therefore the inequalities above, together with $0<r\le T$, give
		positive constants $m_T$ and $M_T$ such that $m_T\le v(r)\le M_T$ for $r\in[r_0,T)$. Consequently $(r,v(r))$ remains in a compact subset of the domain of the right-hand side of \eqref{eq:positive-slope-equation-sec5}.  By the
		extensibility criterion for ordinary differential equations
		\cite[Corollary~2.15]{teschl2012ordinary}, the solution extends beyond
		$T$, contradicting maximality.  Thus $T=\infty$.
		
		\medskip
		
		Because $g_+>0$ on $(y_*,1)$, equation
		\eqref{eq:positive-slope-equation-sec5} gives $v'>0$.  Thus $v$ is
		increasing.  Were $v$ bounded as $r\to\infty$, we would have
		$Q(r,v(r))\to0$, contradicting $y_v(r)>y_*$.  Hence $v(r)\to\infty$.
		
		\medskip
		
		It remains to identify the limiting phase level.  Using
		\eqref{eq:eta-derivative-branch}, we have
		\begin{align*}
			y_v'
			=
			\frac{1}{r}
			\left(
			\bigl[1+(1-2\beta)v^2\bigr]g_+(y_v)
			-
			y_v
			\right).
		\end{align*}
		Fix $\varepsilon\in(0,1-y_*)$.  As $g_+$ is positive on
		$(y_*,1-\varepsilon]$, there is $c_\varepsilon>0$ such that
		\begin{align*}
			g_+(y)\ge c_\varepsilon
			\qquad
			\text{for }y\in(y_*,1-\varepsilon].
		\end{align*}
		As $v(r)\to\infty$, for all sufficiently large $r$ we have
		$y_v'(r)>0$ whenever $y_v(r)\le1-\varepsilon$.
		
		Suppose that $y_v$ did not eventually cross the level $1-\varepsilon$.  Then
		$y_v(r)\le1-\varepsilon$ for all large $r$.  Because $y_v>y_*$ and
		$v(r)\to\infty$, the large-slope relation
		\eqref{eq:eta-large-slope-asymptotic} gives $v(r)\ge c r^\alpha$ for all
		large $r$.  Therefore, on that tail,
		\begin{align*}
			y_v'(r)
			\ge
			c_1 r^{2\alpha-1}
		\end{align*}
		for some $c_1>0$.  The condition $\alpha>\frac{1}{2}$ makes
		$r^{2\alpha-1}$ non-integrable at infinity, so integration gives an
		unbounded increase of $y_v$, contradicting $y_v\le1-\varepsilon$.
		Thus $y_v$ eventually crosses $1-\varepsilon$.
		
		At every sufficiently large point with $y_v=1-\varepsilon$, the formula
		for $y_v'$ gives $y_v'>0$.  Hence, after the first large upward
		crossing, $y_v$ cannot cross that level downward.  Because
		$\varepsilon>0$ is arbitrary and $y_v<1$, we obtain
		$y_v(r)\to1$.
		
		Finally, we use \eqref{eq:eta-large-slope-asymptotic}, which gives
		$y_v(r)=r^{-1}v(r)^{\frac{1}{\alpha}}(1+o(1))$.  Since
		$y_v(r)\to1$, we conclude that $v(r)=r^\alpha(1+o(1))$.
	\end{proof}
	
	The next lemma is the elementary stability mechanism behind the coefficient
	bootstrap.  Once the leading asymptotic scale is known, the normalized
	remainder satisfies a scalar equation whose linear part has a negative
	coefficient of order $r^{2\alpha-1}$.  Because this weight is non-integrable for
	$\alpha>\frac{1}{2}$, the remainder is forced to converge to the equilibrium
	selected by the lower-order terms.
	
	\begin{lemma}
		\label{lem:stable-scalar-bootstrap}
		Let $\alpha>\frac{1}{2}$ and let $c<0$.  Let
		$Z\in\mathcal C^1([R_0,\infty))$ satisfy
		\begin{align*}
			Z'
			=
			c\,r^{2\alpha-1}(Z-L)
			+
			E(r)
		\end{align*}
		for some $L\in\mathbb R$, where
		\begin{align*}
			E(r)
			=
			o\!\left((1+|Z(r)-L|)r^{2\alpha-1}\right)
			\qquad
			\text{as }r\to\infty
		\end{align*}
		along the trajectory.  Then $Z(r)\to L$ as $r\to\infty$.
	\end{lemma}
	
	\begin{proof}
		Set $Y=Z-L$.  Then
		\begin{align*}
			Y'
			=
			c\,r^{2\alpha-1}Y
			+
			E(r),
			\mbox{ with }
			E(r)
			=
			o\!\left((1+|Y(r)|)r^{2\alpha-1}\right).
		\end{align*}
		It remains to prove that this normalized error tends to zero.
		
		\medskip
		
		Fix $\varepsilon>0$.  By the assumption on $E$, for every $\delta>0$ there
		exists $R_\delta\ge R_0$ such that
		\begin{align*}
			|E(r)|
			\le
			\delta(1+|Y(r)|)r^{2\alpha-1}
			\text{ for every }r\ge R_\delta.
		\end{align*}
		On the region $|Y|\ge\varepsilon$, we have
		$1+|Y|\le(1+\varepsilon^{-1})|Y|$.  Choose $\delta>0$ so small that
		$\delta(1+\varepsilon^{-1})\le-\frac{c}{2}$.  Then, for $r\ge R_\delta$,
		\begin{align*}
			Y(r)\ge\varepsilon
			&\quad\Longrightarrow\quad
			Y'(r)\le\frac{c}{2}r^{2\alpha-1}Y(r)<0,\\
			Y(r)\le-\varepsilon
			&\quad\Longrightarrow\quad
			Y'(r)\ge\frac{c}{2}r^{2\alpha-1}Y(r)>0.
		\end{align*}
		These two inequalities imply that, after $R_\delta$, an upward crossing of
		the level $\varepsilon$ and a downward crossing of the level $-\varepsilon$
		are impossible.
		
		\medskip
		
		It remains to show that the trajectory must enter the strip
		$(-\varepsilon,\varepsilon)$.  Suppose first that $Y(r)\ge\varepsilon$ for all
		$r\ge R\ge R_\delta$.  Then the first differential inequality gives
		\begin{align*}
			\frac{d}{dr}\log Y(r)
			\le
			\frac{c}{2}r^{2\alpha-1}
			\qquad
			\text{for }r\ge R.
		\end{align*}
		After integration, the right-hand side tends to $-\infty$ because
		$\alpha>\frac{1}{2}$, contradicting $Y\ge\varepsilon$.  Similarly, if
		$Y(r)\le-\varepsilon$ for all $r\ge R\ge R_\delta$, then
		$U=-Y$ satisfies
		\begin{align*}
			\frac{d}{dr}\log U(r)
			\le
			\frac{c}{2}r^{2\alpha-1},
		\end{align*}
		which again contradicts $U\ge\varepsilon$ after integration.
		
		Thus, for every $\varepsilon>0$, the function $Y$ eventually enters
		$(-\varepsilon,\varepsilon)$ and cannot leave it through either boundary
		after increasing the entry radius if necessary.  Hence
		$\limsup\limits_{r\to\infty}Y(r)\le\varepsilon$ and
		$\liminf\limits_{r\to\infty}Y(r)\ge-\varepsilon$.  Letting
		$\varepsilon\downarrow0$ gives $Y(r)\to0$, and therefore $Z(r)\to L$.
	\end{proof}
	
	Before applying the bootstrap, we record the uniform expansions of the
	right-hand side of the positive graphical equation near the nondegenerate
	endpoint $y=1$.  These estimates translate the endpoint expansion of $g_+$
	into the normalized scales used for the slope corrections.  The uniformity of
	the remainders will allow us to apply Lemma~\ref{lem:stable-scalar-bootstrap}
	successively to the coefficient remainders.
	
	\begin{lemma}
		\label{lem:uniform-expansions-nondegenerate}
		Let
		\begin{align*}
			\mathcal F(r,s)
			=
			(1+s^2)^{\beta+1}
			g_+\left(\frac{s}{r(1+s^2)^\beta}\right).
		\end{align*}
		Then the following expansions hold as $r\to\infty$.  In each item, the
		remainder is uniform in the following sense: if the corresponding normalized
		parameter is restricted to a bounded set, then the quotient between the
		remainder and the displayed scale tends to zero uniformly on that set. Therefore, by setting $\Theta=\Theta(r)=o(r^\alpha)$, it holds 
		
	\begin{align*}
		\textup{(E1)}\quad
		&\mathcal F(r,r^\alpha+\Theta)-\alpha r^{\alpha-1}
		=
		\frac{g_+'(1)}{\alpha}\Theta r^{2\alpha-1}
		+
		o\!\left((1+|\Theta|)r^{2\alpha-1}\right);
		\\[4pt]
		\textup{(E2)}\quad
		& r^\alpha \mathcal F\left(r,r^\alpha+\frac{\Theta}{r^\alpha}\right)
		-
		\alpha r^{2\alpha-1}
		+
		\frac{\alpha\Theta}{r}
		\\
		& \qquad =
		\frac{g_+'(1)}{\alpha}(\Theta+a)r^{2\alpha-1}
		+
		o\!\left((1+|\Theta+a|)r^{2\alpha-1}\right);
		\\[4pt]
		\textup{(E3)}\quad
		& r^\alpha\left[
		r^\alpha \mathcal F\left(r,r^\alpha-\frac{a}{r^\alpha}
		+
		\frac{\Theta}{r^{2\alpha}}\right)
		-
		\alpha r^{2\alpha-1}
		-
		a\alpha r^{-1}
		\right]
		+
		\frac{2\alpha\Theta}{r}
		\\
		& \qquad =
		\frac{g_+'(1)}{\alpha}\Theta r^{2\alpha-1}
		+
		o\!\left((1+|\Theta|)r^{2\alpha-1}\right);
		\\[4pt]
		\textup{(E4)}\quad
		& r^\alpha\left[
		r^\alpha\left(
		r^\alpha \mathcal F\left(r,r^\alpha-\frac{a}{r^\alpha}
		+
		\frac{\Theta}{r^{3\alpha}}\right)
		-
		\alpha r^{2\alpha-1}
		-
		a\alpha r^{-1}
		\right)
		\right]
		+
		\frac{3\alpha\Theta}{r}
		\\
		& \qquad =
		\frac{g_+'(1)}{\alpha}(\Theta-b)r^{2\alpha-1}
		+
		o\!\left((1+|\Theta-b|)r^{2\alpha-1}\right);
		\\[4pt]
		\textup{(E5)}\quad
		& r^\alpha\left[
		r^\alpha\left[
		r^\alpha\left(
		r^\alpha \mathcal F\left(
		r,
		r^\alpha-
		\frac{a}{r^\alpha}
		+
		\frac{b}{r^{3\alpha}}
		+
		\frac{\Theta}{r^{4\alpha}}
		\right)
		-
		\alpha r^{2\alpha-1}
		-
		a\alpha r^{-1}
		\right)
		\right]
		\right]
		+
		3\alpha b r^{\alpha-1}
		+
		\frac{4\alpha\Theta}{r}
		\\
		& \qquad =
		\frac{g_+'(1)}{\alpha}\Theta r^{2\alpha-1}
		+
		o\!\left((1+|\Theta|)r^{2\alpha-1}\right).
	\end{align*}
	\end{lemma}
	
	\begin{proof}
		Set $q=r^{-2\alpha}$ and write $s=r^\alpha(1+\xi)$, where
		$\xi=o(1)$.  All estimates below are uniform for the normalized variables
		appearing in the statement on bounded sets.
		
		\medskip
		
		Because $1-2\beta=\alpha^{-1}$, we have the exact identity
		\begin{align}
			\label{eq:master-eta-expansion}
			\frac{s}{r(1+s^2)^\beta}
			=
			(1+\xi)^{\frac{1}{\alpha}}
			\left(1+q(1+\xi)^{-2}\right)^{-\beta}.
		\end{align}
		Taylor's theorem gives, uniformly in this regime,
		\begin{align*}
			\frac{s}{r(1+s^2)^\beta}
			=
			1
			+
			\left(\frac{\xi}{\alpha}-\beta q\right)
			+
			O(\xi^2+|\xi|q+q^2).
		\end{align*}
		Similarly,
		\begin{align}
			\label{eq:master-prefactor-expansion}
			(1+s^2)^{\beta+1}
			=
			r^{3\alpha-1}(1+\xi)^{2\beta+2}
			\left(1+q(1+\xi)^{-2}\right)^{\beta+1}.
		\end{align}
		
		\medskip
		
		By the endpoint expansion in
		Theorem~\ref{thm:positive-bowl-asymptotics}\textup{(i)},
		\begin{align*}
			g_+(y)
			=
			a_1(1-y)
			+
			a_2(1-y)^2
			+
			o((1-y)^2)
			\mbox{ as }y\to1^-.
		\end{align*}
		Equivalently,
		\begin{align}
			\label{eq:master-g-expansion}
			g_+(y)
			=
			g_+'(1)(y-1)
			+
			\frac{g_+''(1)}{2}(y-1)^2
			+
			o(|y-1|^2),
		\end{align}
		with one-sided uniformity near $y=1$.  Combining
		\eqref{eq:master-eta-expansion}, \eqref{eq:master-prefactor-expansion}, and
		\eqref{eq:master-g-expansion} gives an expansion of $\mathcal F(r,s)$ whose
		remainders are uniform in each normalized regime listed in the statement.
		
		\medskip
		
		For the first estimate, we take $s=r^\alpha+\Theta$, so
		$\xi=\Theta r^{-\alpha}=o(1)$.  The linear term in $\Theta$ is
		$g_+'(1)\alpha^{-1}\Theta r^{2\alpha-1}$.  The terms independent of
		$\Theta$ are lower order than $r^{2\alpha-1}$, and the quadratic and mixed
		remainders are
		$o((1+|\Theta|)r^{2\alpha-1})$ with the stated uniformity.
		
		\medskip
		
		For the higher normalizations, we successively set
		\begin{align*}
			s
			=
			r^\alpha+
			\frac{\Theta}{r^\alpha},
			\qquad
			s
			=
			r^\alpha-
			\frac{a}{r^\alpha}
			+
			\frac{\Theta}{r^{2\alpha}},
		\end{align*}
		\begin{align*}
			s
			=
			r^\alpha-
			\frac{a}{r^\alpha}
			+
			\frac{\Theta}{r^{3\alpha}},
			\qquad
			s
			=
			r^\alpha-
			\frac{a}{r^\alpha}
			+
			\frac{b}{r^{3\alpha}}
			+
			\frac{\Theta}{r^{4\alpha}}.
		\end{align*}
		In each case $\xi=o(1)$, and the expansion above applies uniformly on the
		corresponding bounded normalized ranges.
		
		\medskip
		
		The constants $a$ and $b$ are those in
		Theorem~\ref{thm:positive-bowl-asymptotics}\textup{(i)}.  With
		$L_1=-a\alpha^{-1}-\beta$, the first cancellation is
		$g_+'(1)L_1=\alpha$.  The definition of $b$ gives the next cancellation,
		whose remaining quadratic contribution is encoded by
		\begin{align*}
			\mathfrak P(a,\beta)
			=
			\frac{1}{2}(1-2a)\beta\bigl(1+\beta(1-2a)\bigr)
			-
			a^2\beta.
		\end{align*}
		After these cancellations, the dependence on the free normalized parameter
		is linear with coefficient $g_+'(1)\alpha^{-1}$ at the scale
		$r^{2\alpha-1}$; in the fourth normalization the centered parameter is
		$\Theta-b$.  The remaining differentiated normalization terms, such as
		$\Theta r^{-1}$ and $r^{\alpha-1}$, are lower order relative to the displayed
		scales.  This proves all five estimates.
	\end{proof}
	
	
	\begin{remark}
		\label{rem:nondegenerate-expansion-organization}
		The estimates above are used in the proof of
		Theorem~\ref{thm:positive-bowl-asymptotics} as follows.  Once the leading
		behavior $v(r)=r^\alpha(1+o(1))$ is known, the associated phase variable
		remains in a fixed neighborhood of the cylindrical endpoint.  Therefore the
		expansions in Lemma~\ref{lem:uniform-expansions-nondegenerate} may be applied
		with remainders that are uniform along the solution.
		
\medskip

		Then we normalize the error successively.  At each stage, the equation for the
		new remainder has the scalar form treated in
		Lemma~\ref{lem:stable-scalar-bootstrap}.  The first normalization identifies
		the correction of order $r^{-\alpha}$, the next one removes the possible order
		$r^{-2\alpha}$ term, and the following normalizations identify the coefficient
		of order $r^{-3\alpha}$ and rule out an order $r^{-4\alpha}$ remainder.  This
		gives
		\begin{align*}
			v(r)
			=
			r^\alpha
			-
			\frac{a}{r^\alpha}
			+
			\frac{b}{r^{3\alpha}}
			+
			o(r^{-4\alpha})
			\qquad\text{as }r\to\infty .
		\end{align*}
	
		
	\end{remark}
	
	
	At this point, we prove the nondegenerate expansion in
	Theorem~\ref{thm:positive-bowl-asymptotics}\textup{(i)}.  We set
	$\mathcal F(r,s)=(1+s^2)^{\beta+1}g_+(Q(r,s))$, so that the positive equation
	is $v'=\mathcal F(r,v)$.  The uniform expansions of $\mathcal F$ near the
	cylindrical endpoint are given by
	Lemma~\ref{lem:uniform-expansions-nondegenerate}, and the convergence of the
	successive normalized coefficients follows from
	Lemma~\ref{lem:stable-scalar-bootstrap}.
	
	\begin{proof}[Proof of Theorem~\ref{thm:positive-bowl-asymptotics}\textup{(i)}]
		By Proposition~\ref{prop:positive-end-leading-order},
		$v(r)=r^\alpha(1+o(1))$.  We start with the normalization
		$v(r)=r^\alpha+\Phi_0(r)$, with $\Phi_0=o(r^\alpha)$.  Because
		$v'=\mathcal F(r,v)$, Lemma~\ref{lem:uniform-expansions-nondegenerate}
		gives
		\begin{align*}
			\Phi_0'
			=
			\frac{g_+'(1)}{\alpha}\Phi_0 r^{2\alpha-1}
			+
			o\!\left((1+|\Phi_0|)r^{2\alpha-1}\right).
		\end{align*}
		As $g_+'(1)<0$, Lemma~\ref{lem:stable-scalar-bootstrap} gives
		$\Phi_0(r)\to0$.
		
		\medskip
		
		The first convergence improves the scale, and we write
		$v(r)=r^\alpha+\Phi_1(r)r^{-\alpha}$.  Then
		$\Phi_1=o(r^\alpha)$, and
		$v'=\alpha r^{\alpha-1}+\Phi_1'r^{-\alpha}-\alpha\Phi_1r^{-(\alpha+1)}$.  The second estimate in Lemma~\ref{lem:uniform-expansions-nondegenerate} gives
		\begin{align*}
			\Phi_1'
			=
			\frac{g_+'(1)}{\alpha}(\Phi_1+a)r^{2\alpha-1}
			+
			o\!\left((1+|\Phi_1+a|)r^{2\alpha-1}\right).
		\end{align*}
		Thus $\Phi_1(r)\to-a$, and
		$v(r)=r^\alpha-\frac{a}{r^\alpha}+o(r^{-\alpha})$.
		
		\medskip
		
		After the coefficient $a$ has been identified, we write
		$v(r)=r^\alpha-ar^{-\alpha}+\Phi_2(r)r^{-2\alpha}$.
		Then $\Phi_2=o(r^\alpha)$, and
		$v'=\alpha r^{\alpha-1}+a\alpha r^{-\alpha-1}+\Phi_2'r^{-2\alpha}-
		2\alpha\Phi_2r^{-(2\alpha+1)}$.  The third estimate gives
		\begin{align*}
			\Phi_2'
			=
			\frac{g_+'(1)}{\alpha}\Phi_2 r^{2\alpha-1}
			+
			o\!\left((1+|\Phi_2|)r^{2\alpha-1}\right).
		\end{align*}
		Hence $\Phi_2(r)\to0$, and
		$v(r)=r^\alpha-ar^{-\alpha}+o(r^{-2\alpha})$.
		
		\medskip
		
		The next normalization tests the coefficient at order $r^{-3\alpha}$.  We write
		$v(r)=r^\alpha-ar^{-\alpha}+\Phi_3(r)r^{-3\alpha}$.
		Then $\Phi_3=o(r^\alpha)$, and
		$v'=\alpha r^{\alpha-1}+a\alpha r^{-\alpha-1}+\Phi_3'r^{-3\alpha}-
		3\alpha\Phi_3r^{-(3\alpha+1)}$.  The fourth estimate gives
		\begin{align*}
			\Phi_3'
			=
			\frac{g_+'(1)}{\alpha}(\Phi_3-b)r^{2\alpha-1}
			+
			o\!\left((1+|\Phi_3-b|)r^{2\alpha-1}\right).
		\end{align*}
		Therefore $\Phi_3(r)\to b$, and
		$v(r)=r^\alpha-ar^{-\alpha}+br^{-3\alpha}+o(r^{-3\alpha})$.
		
		\medskip
		
		Finally, after the coefficient $b$ has been fixed, we write
		$v(r)=r^\alpha-ar^{-\alpha}+br^{-3\alpha}+
		\Phi_4(r)r^{-4\alpha}$.  Then $\Phi_4=o(r^\alpha)$, and
		$v'=\alpha r^{\alpha-1}+a\alpha r^{-\alpha-1}-3\alpha b r^{-3\alpha-1}+
		\Phi_4'r^{-4\alpha}-4\alpha\Phi_4r^{-(4\alpha+1)}$.  The fifth
		estimate gives
		\begin{align*}
			\Phi_4'
			=
			\frac{g_+'(1)}{\alpha}\Phi_4 r^{2\alpha-1}
			+
			o\!\left((1+|\Phi_4|)r^{2\alpha-1}\right).
		\end{align*}
		One last application of Lemma~\ref{lem:stable-scalar-bootstrap} yields
		$\Phi_4(r)\to0$.  Consequently,
		\begin{align*}
			v(r)
			=
			r^\alpha
			-
			\frac{a}{r^\alpha}
			+
			\frac{b}{r^{3\alpha}}
			+
			o(r^{-4\alpha}),
		\end{align*}
		as claimed.
	\end{proof}
	
	\subsection{Degenerate positive branches}\,
	\label{subsec:degenerate-positive-branches}
	
	For the degenerate regime, we assume that $\gamma$ is $1$-degenerate.  Then $I_+=(y_*,\infty)$, and
	the positive branch is controlled by its behavior as $y\to\infty$.  Before
	using the asymptotic expansion of $g_+$, we verify that the positive end
	actually samples this large-level region.
	
	\begin{lemma}
		\label{lem:eta-goes-to-infinity-degenerate}
		Assume that $\gamma$ is $1$-degenerate and that the positive branch is
		defined on $I_+=(y_*,\infty)$.  Let $v$ be a positive solution of
		\eqref{eq:positive-slope-equation-sec5}, defined for all sufficiently large
		$r$, admissible for the positive branch, and satisfying $v(r)\to\infty$.
		Then
		\begin{align*}
			y_v(r)=Q(r,v(r))\to\infty
			\qquad
			\text{as }r\to\infty.
		\end{align*}
	\end{lemma}
	
	\begin{proof}
		Fix $M>y_*$.  We prove that $y_v(r)>M$ for all
		sufficiently large $r$.  The admissibility of the positive solution gives
		$y_v(r)\in(y_*,\infty)$ for all large $r$.  From
		\eqref{eq:eta-derivative-branch} and $1-2\beta=\alpha^{-1}$, we get
		\begin{align*}
			y_v'
			=
			\frac{1}{r}
			\left(
			\left[1+\frac{v^2}{\alpha}\right]g_+(y_v)
			-
			y_v
			\right).
		\end{align*}
		On the region $y_v\le M$, the monotonicity of $g_+$ gives
		$g_+(y_v)\ge g_+(M)>0$.  Therefore, at every point where
		$y_v(r)\le M$,
		\begin{align*}
			y_v'(r)
			&\ge
			\frac{1}{r}
			\left(
			\left[1+\frac{v(r)^2}{\alpha}\right]g_+(M)
			-
			M
			\right).
		\end{align*}
		Because $v(r)\to\infty$, we can choose $R_M$ so large that
		\begin{align*}
			\frac{g_+(M)}{\alpha}v(r)^2
			\ge
			2M
			\qquad
			\text{for every }r\ge R_M.
		\end{align*}
		Hence, for $r\ge R_M$ and $y_v(r)\le M$,
		\begin{align}
			\label{eq:eta-crossing-lower-bound}
			y_v'(r)
			\ge
			\frac{g_+(M)+M}{r}
			>
			\frac{M}{r}
			>
			0.
		\end{align}
				
		The level $M$ must be crossed.  Otherwise, for some $R\ge R_M$ we would have
		$y_v(r)\le M$ for all $r\ge R$.  Integrating
		\eqref{eq:eta-crossing-lower-bound} on $[R,r]$ gives
		\begin{align*}
			y_v(r)
			\ge
			y_v(R)
			+
			M\int_R^r\frac{d\tau}{\tau}
			=
			y_v(R)
			+
			M\log\frac{r}{R},
		\end{align*}
		which eventually exceeds $M$, a contradiction.  Once the level $M$ is
		crossed after $R_M$, no later downward crossing is possible, because
		\eqref{eq:eta-crossing-lower-bound} gives $y_v'>0$ at every later point
		where $y_v=M$.  Thus $y_v(r)>M$ for all sufficiently large $r$.
		Because $M>y_*$ was arbitrary, $y_v(r)\to\infty$.
	\end{proof}
	
	\begin{proof}[Proof of Theorem~\ref{thm:positive-bowl-asymptotics}\textup{(ii)}]
		Lemma~\ref{lem:eta-goes-to-infinity-degenerate} gives
		$y_v(r)\to\infty$.  Hence
		\begin{align*}
			g_+(y_v(r))
			=
			c_\gamma y_v(r)^{-k_\gamma}(1+o(1)).
		\end{align*}
		The large-slope expansion gives $y_v(r)=r^{-1}v(r)^{\frac{1}{\alpha}}(1+o(1))$,
		and therefore $y_v(r)^{-k_\gamma}=r^{k_\gamma}v(r)^{-\frac{k_\gamma}{\alpha}}(1+o(1))$. Consequently, we obtain 
		\begin{align*}
			g_+(y_v)
			=
			c_\gamma r^{k_\gamma}v^{-\frac{k_\gamma}{\alpha}}(1+o(1)).
		\end{align*}
		Moreover, since $v(r)\to\infty$, we have $(1+v^2)^{\beta+1}=v^{2\beta+2}(1+v^{-2})^{\beta+1}=v^{\frac{3\alpha-1}{\alpha}}(1+o(1))$. Thus, the positive equation becomes
		\begin{align}
			\label{eq:degenerate-power-main-ode}
			v'
			=
			c_\gamma r^{k_\gamma}
			v^{\frac{3\alpha-1-k_\gamma}{\alpha}}(1+o(1)).
		\end{align}
		
		Set $\mu_\gamma=\alpha^{-1}(k_\gamma+1-2\alpha)$.  Then, since 
		$k_\gamma>2\alpha-1$, we have $\mu_\gamma>0$.  Multiplying
		\eqref{eq:degenerate-power-main-ode} by
		$\mu_\gamma v^{\mu_\gamma-1}$ gives
		\begin{align*}
			\frac{d}{dr}\left(v^{\mu_\gamma}\right)
			=
			\mu_\gamma c_\gamma r^{k_\gamma}(1+o(1)).
		\end{align*}
		The error integrates to $o(r^{k_\gamma+1})$: for every $\varepsilon>0$, the
		error is bounded by $\varepsilon r^{k_\gamma}$ for large $r$, and
		$k_\gamma>-1$.  Hence
		\begin{align*}
			v(r)^{\mu_\gamma}
			=
			\frac{\mu_\gamma c_\gamma}{k_\gamma+1}r^{k_\gamma+1}
			+
			o(r^{k_\gamma+1}).
		\end{align*}
		Finally, by taking the power $\mu_\gamma^{-1}$, we obtain 
		\begin{align*}
			v(r)
			=
			A_\gamma r^{d_\gamma}(1+o(1)),
			\mbox{ with }
			d_\gamma
			=
			\frac{k_\gamma+1}{\mu_\gamma}
			=
			\frac{\alpha(k_\gamma+1)}{k_\gamma+1-2\alpha},
		\end{align*}
		where $A_\gamma^{\mu_\gamma}=c_\gamma d_\gamma^{-1}$. This proves the power-growth asymptotic.
	\end{proof}
	
	\begin{proof}[Proof of Theorem~\ref{thm:positive-bowl-asymptotics}\textup{(iii)}]
		Lemma~\ref{lem:eta-goes-to-infinity-degenerate} gives
		$y_v(r)\to\infty$.  With $k_\gamma=2\alpha-1$, the branch expansion and
		\eqref{eq:eta-large-slope-asymptotic} give
		\begin{align*}
			g_+(y_v)
			=
			c_\gamma y_v^{-(2\alpha-1)}(1+o(1))
			=
			c_\gamma r^{2\alpha-1}
			v^{-\frac{2\alpha-1}{\alpha}}(1+o(1)).
		\end{align*}
		The prefactor satisfies $(1+v^2)^{\beta+1}=v^{\frac{3\alpha-1}{\alpha}}(1+o(1))$.
		Therefore, we have $v'=c_\gamma r^{2\alpha-1}v(1+o(1))$. Dividing by $v>0$, we obtain
		\begin{align*}
			\frac{d}{dr}\log v(r)
			=
			c_\gamma r^{2\alpha-1}(1+o(1)).
		\end{align*}
		Integrating from a fixed large radius $R$ to $r$ gives
		\begin{align*}
			\log v(r)
			=
			\log v(R)
			+
			c_\gamma\int_R^r s^{2\alpha-1}(1+o(1))\,ds.
		\end{align*}
		The error term integrates to $o(r^{2\alpha})$: for every $\varepsilon>0$,
		the error is bounded by $\varepsilon s^{2\alpha-1}$ for large $s$, and
		$\int^r s^{2\alpha-1}\,ds=\frac{1}{2\alpha}r^{2\alpha}+O(1)$.  Hence
		$\log v(r)=\frac{c_\gamma}{2\alpha}r^{2\alpha}+o(r^{2\alpha})$.
	\end{proof}
	
	\begin{proof}[Proof of Theorem~\ref{thm:positive-bowl-asymptotics}\textup{(iv)}]
		For notational economy, we write $k=k_\gamma$, $c=c_\gamma$, $c_1=c_{\gamma,1}$,
		$d=d_\gamma$, and $A=A_\gamma$.  By
		Theorem~\ref{thm:positive-bowl-asymptotics}\textup{(ii)},
		$v(r)=Ar^d(1+o(1))$.  We first refine the positive equation.
		
		\medskip
		
		The large-slope expansion of the phase variable gives $y_v=r^{-1}v^{\frac{1}{\alpha}}
			\left(1-\beta v^{-2}+O(v^{-4})\right)$. Therefore
		\begin{align*}
			y_v^{-(2k+1)}
			=
			r^{2k+1}v^{-\frac{2k+1}{\alpha}}(1+o(1)).
		\end{align*}
		Using the assumed second-order expansion of $g_+$, we obtain
		\begin{align*}
			g_+(y_v)
			&=
			c\,y_v^{-k}
			+
			c_1y_v^{-(2k+1)}
			+
			o(y_v^{-(2k+1)}) \\
			&=
			c\,r^k v^{-\frac{k}{\alpha}}
			\left[
			1
			+
			k\beta v^{-2}
			+
			\frac{c_1}{c}r^{k+1}v^{-\frac{k+1}{\alpha}}
			+
			o(r^{-2d})
			\right].
		\end{align*}
		Here, we used $v=Ar^d(1+o(1))$, which gives
		$r^{k+1}v^{-\frac{k+1}{\alpha}}=O(r^{-2d})$.
		
		\medskip
		
		The prefactor satisfies
		\begin{align*}
			(1+v^2)^{\beta+1}
			=
			v^{\frac{3\alpha-1}{\alpha}}
			\left(1+(\beta+1)v^{-2}+O(v^{-4})\right).
		\end{align*}
		Multiplying the last two expansions gives
		\begin{align}
			\label{eq:second-order-degenerate-refined-ode}
			v'
			=
			c\,r^k v^{\frac{3\alpha-1-k}{\alpha}}
			\left[
			1
			+
			(1+\beta(k+1))v^{-2}
			+
			\frac{c_1}{c}r^{k+1}v^{-\frac{k+1}{\alpha}}
			+
			o(r^{-2d})
			\right].
		\end{align}
		Substituting the leading asymptotic in the correction terms yields
		\begin{align*}
			(1+\beta(k+1))v^{-2}
			+
			\frac{c_1}{c}r^{k+1}v^{-\frac{k+1}{\alpha}}
			=
			\mathcal C_\gamma r^{-2d}
			+
			o(r^{-2d}),
		\end{align*}
		where $\mathcal C_\gamma=(1+\beta(k+1))A^{-2}+c_1c^{-1}A^{-\frac{k+1}{\alpha}}$. Hence
		\begin{align}
			\label{eq:second-order-degenerate-reduced-ode}
			v'
			=
			c\,r^k v^{\frac{3\alpha-1-k}{\alpha}}
			\left[
			1+\mathcal C_\gamma r^{-2d}+o(r^{-2d})
			\right].
		\end{align}
		
		\medskip
		
		Next, we normalize around the leading term by setting
		$v=Ar^d(1+z)$, with $z\to0$.  The constants $A$ and $d$ satisfy
		$dA=cA^{\frac{3\alpha-1-k}{\alpha}}$.  Substituting
		$v=Ar^d(1+z)$ into \eqref{eq:second-order-degenerate-reduced-ode} and
		expanding in $z$ gives
		\begin{align}
			\label{eq:second-order-z-equation}
			z'
			=
			-\frac{k+1}{r}z
			+
			d\,\mathcal C_\gamma r^{-1-2d}
			+
			O\!\left(\frac{z^2}{r}\right)
			+
			o(r^{-1-2d}).
		\end{align}
		Indeed, the coefficient of the linear term is
		$d\left(\frac{3\alpha-1-k}{\alpha}-1\right)=-(k+1)$.
		
		\medskip
		
		First, we consider the nonresonant case $k+1>4\alpha$, equivalently
		$k+1>2d$.  Before rescaling by $r^{2d}$, we record a preliminary bound.  Let $C_1>0$ control the
		quadratic term in \eqref{eq:second-order-z-equation}.  Because $z\to0$, for
		every $\delta>0$ we may take $R_\delta$ so large that $|z(r)|\le\delta$ for
		$r\ge R_\delta$ and the $o(r^{-1-2d})$ term has absolute value bounded by
		$\delta r^{-1-2d}$.  Choose $\delta$ so small that
		$\lambda=k+1-C_1\delta>2d$.  A standard comparison applied to $|z|$ gives,
		for $r\ge R_\delta$,
		\begin{align*}
			\frac{d}{dr}|z(r)|
			\le
			-\frac{\lambda}{r}|z(r)|
			+
			C_2 r^{-1-2d}
		\end{align*}
		at points where $z(r)\ne0$, and hence by continuity on the whole tail.  After
		multiplying by $r^\lambda$ and integrating, we obtain
		\begin{align*}
			|z(r)|
			\le
			C r^{-2d}
			+
			C r^{-\lambda}
			=
			O(r^{-2d}).
		\end{align*}
		Thus $Y=r^{2d}z$ is bounded.
		
		\medskip
		
		Set $z=r^{-2d}Y$.  From \eqref{eq:second-order-z-equation},
		\begin{align*}
			Y'
			=
			-\frac{k+1-2d}{r}Y
			+
			\frac{d\,\mathcal C_\gamma}{r}
			+
			O\!\left(\frac{Y^2}{r^{2d+1}}\right)
			+
			o\left(\frac{1}{r}\right).
		\end{align*}
		Equivalently,
		\begin{align*}
			Y'
			=
			-\frac{k+1-2d}{r}
			\left(
			Y-
			\frac{d\,\mathcal C_\gamma}{k+1-2d}
			\right)
			+
			O\!\left(\frac{Y^2}{r^{2d+1}}\right)
			+
			o\left(\frac{1}{r}\right).
		\end{align*}
		Because $Y$ is bounded, the quadratic term is $O(r^{-2d-1})=o(r^{-1})$.
		The same crossing argument used in Lemma~\ref{lem:stable-scalar-bootstrap},
		with the non-integrable weight $r^{-1}$, gives
		\begin{align*}
			Y(r)
			\to
			Y_*
			=
			\frac{d\,\mathcal C_\gamma}{k+1-2d}.
		\end{align*}
		Consequently, $z(r)=Y_*r^{-2d}+o(r^{-2d})$, and therefore
		\begin{align*}
			v(r)
			=
			Ar^d(1+z(r))
			=
			Ar^d
			+
			AY_*r^{-d}
			+
			o(r^{-d}).
		\end{align*}
		Using $\frac{d}{k+1-2d}=\frac{\alpha}{k+1-4\alpha}$,
		we obtain $B_\gamma=AY_*=\frac{\alpha A_\gamma\mathcal C_\gamma}{k_\gamma+1-4\alpha}$. This proves the nonresonant second-order expansion.
		
		\medskip
		
		Finally, we turn to the resonant case $k+1=4\alpha$, equivalently
		$k+1=2d$.  We set $Y=r^{2d}z$.  Multiplying
		\eqref{eq:second-order-z-equation} by $r^{2d}$ gives
		\begin{align}
			\label{eq:second-order-resonant-Y}
			Y'=d\,\mathcal C_\gamma r^{-1}
			+
			O\!\left(Y^2r^{-(2d+1)}\right)
			+
			o\left(r^{-1}\right).
		\end{align}
		First, we obtain a rough bound for $Y$.  Since $z=r^{-2d}Y=o(1)$, for every
		$\delta>0$ there is $R_\delta$ such that $|z(r)|\le\delta$ for every $r\ge R_\delta$. Let $C_0>0$ be such that the quadratic term in
		\eqref{eq:second-order-resonant-Y} satisfies $\left|
			O\!\left(Y^2r^{-(2d+1)}\right)\right|\le C_0Y^2r^{-(2d+1)}$ for all sufficiently large $r$.  Because $Y=r^{2d}z$, we have after increasing $R_\delta$ if necessary, $\left|O\!\left(\frac{Y^2}{r^{2d+1}}\right)\right|
			\le C_0\delta\frac{|Y|}{r}$ for every $r\ge R_\delta$.
		
		The radius $R_\delta$ is also chosen so large that the $o(r^{-1})$ term in
		\eqref{eq:second-order-resonant-Y} has absolute value at most
		$\frac{\delta}{r}$ for $r\ge R_\delta$.  Hence, for $r\ge R_\delta$,
		\begin{align*}
			|Y'(r)|
			(\le|d\,\mathcal C_\gamma|+\delta)r^{-1}
			+
			C_0\delta|Y(r)|r^{-1}.
		\end{align*}
		Integrating from $R_\delta$ to $r$ gives
		\begin{align*}
			|Y(r)|
			\le
			|Y(R_\delta)|
			+
			(|d\,\mathcal C_\gamma|+\delta)\log\frac{r}{R_\delta}
			+
			C_0\delta\int_{R_\delta}^r\frac{|Y(s)|}{s}\,ds.
		\end{align*}
		The hypotheses of Gronwall's inequality are satisfied on every compact
		interval $[R_\delta,r]$: $|Y|$ is continuous, $s^{-1}$ is integrable on
		compact subintervals of $(R_\delta,\infty)$, and the right-hand side is
		nonnegative.  Therefore
		\begin{align*}
			|Y(r)|
			\le
			\left(
			|Y(R_\delta)|
			+
			(|d\,\mathcal C_\gamma|+\delta)\log\frac{r}{R_\delta}
			\right)
			\exp\left(C_0\delta\int_{R_\delta}^r\frac{ds}{s}\right).
		\end{align*}
		Thus $|Y(r)|\le C_\delta r^{C_0\delta}\log r$ for all sufficiently large
		$r$.  Now choose $\delta>0$ so small that $C_0\delta<d$.  Taking any
		$\varepsilon$ with $C_0\delta<\varepsilon<d$, we obtain
		\begin{align*}
			Y(r)=O(r^\varepsilon)
			\mbox{for some }\varepsilon<d.
		\end{align*}
		This rough estimate improves the quadratic term:
		\begin{align*}
			Y^2r^{-(2d+1)}=
			O(r^{2\varepsilon-2d-1})=
			o\!\left(r^{-1}\right).
		\end{align*}
		Equation \eqref{eq:second-order-resonant-Y} therefore reduces to
		$Y'=d\,\mathcal C_\gamma r^{-1}+o\left(r^{-1}\right)$. Integrating from a fixed large radius to $r$ gives $Y(r)=d\,\mathcal C_\gamma\log r+o(\log r)$. Recalling that $Y=r^{2d}z$, we obtain $z(r)=d\,\mathcal C_\gamma r^{-2d}\log r+o(r^{-2d}\log r)$. Therefore
		\begin{align*}
			v(r)
			=
			Ar^d(1+z(r))
			=
			Ar^d
			+
			Ad\,\mathcal C_\gamma r^{-d}\log r
			+
			o(r^{-d}\log r).
		\end{align*}
		Thus $L_\gamma=A_\gamma d_\gamma\mathcal C_\gamma$, and the proof is
		complete.
	\end{proof}
	

\section{Catenoidal translators}
\label{sec:catenoidal-translators}

In this section, we assemble the selected signed neck charts of
Section~\ref{sec:curvature-functions-branches} with the comparison theory of
Section~\ref{sec:barrier-method}.  The selected signed neck chart is $x=g(y,z)$,
$\widehat\gamma(g(y,z),y)=z$, defined on its maximal level domain $\mathcal D$, and $U_g$ denotes the
associated conical sheet. Recall that for the neck parameter, we define the selected neck-radius set by $\mathcal I_0=\{R>0:(R^{-1},0)\in\mathcal D\}$. 

\medskip

Furthermore, for the lower graphical side, we read the reoriented phase through the reflected chart $g^\sharp(y,z)=-g(-y,-z)$, defined in $\mathcal D^\sharp=\{(y,z):(-y,-z)\in\mathcal D\}$. Moreover, whenever a lower branch appears in this section, it means a selected $g^\sharp$-lower branch in the sense of Definition~\ref{def:selected-lower-branch}: $I_-\subset\{y:(y,-1)\in\mathcal D^\sharp\}$, and $g_-(y)=g^\sharp(y,-1)$ on $I_-$.
Throughout the section, we keep the standing assumption $\alpha>\frac{1}{2}$ as in
Section~\ref{sec:barrier-method}.

\subsection{Local necks and the reflected lower branch}\,
\label{subsec:local-catenoidal-neck}

Let $\mathtt c(s)=(r(s),u(s))$ be arclength-parametrized. Then, we have that on the selected sheet the signed phase equation is equivalent to
\begin{align}
	\begin{cases}\label{eq:local-neck-phase-system}
		\dot r=\cos\theta,\quad
		\dot u=\sin\theta,\quad
		\dot\theta=g\left(r^{-1}\sin\theta,\cos\theta\right), 
		\\
		\left(r^{-1}\sin\theta,\cos\theta\right)\in\mathcal D.
	\end{cases}	
\end{align}
Here the phase variables are $x=\dot\theta$, $y=r^{-1}\sin\theta$, $z=\cos\theta$, so that the signed level equation reads $\widehat\gamma(x,y)=z$.

\begin{proposition}
	\label{prop:local-catenoidal-neck}
	For every $R\in\mathcal I_0$, there exist $\varepsilon_R>0$ and a unique
	solution
	\begin{align*}
		\mathtt c_R(s)=(r_R(s),u_R(s)),
		\qquad
		|s|<\varepsilon_R,
	\end{align*}
	of \eqref{eq:local-neck-phase-system} satisfying $\mathtt{c}_R(0)=(R,0)$ and  $\theta_R(0)=\frac{\pi}{2}$. Moreover, $\dot\theta_R(0)=-\rho_0R^{-1}$, $\ddot r_R(0)=\rho_0R^{-1}>0$. Finally, after decreasing $\varepsilon_R$, if necessary, the two sides $s>0$ and
	$s<0$ are vertical graphs over $r>R$ close to $R$, denoted by $u_R^+$ and
	$u_R^-$.
\end{proposition}

\begin{proof}
	To begin, we fix $R\in\mathcal I_0$.  Then $(R^{-1},0)\in\mathcal D$ and since $\mathcal D$ is open,
	there is a neighborhood $\mathcal V\subset\mathcal D$ of $(R^{-1},0)$.  By
	the regularity of the selected signed neck data and
	Proposition~\ref{prop:conical-sheet-maximal-chart}, the function $g(y,z)$ is of
	class $\mathcal C^1$ on $\mathcal V$.
	
	\medskip
	
	Next, we consider the map 
	\begin{align}\label{map for local necks}
		(r,u,\theta)
		\mapsto
		\left(
		\cos\theta,
		\sin\theta,
		g\left(r^{-1}\sin\theta,\cos\theta\right)
		\right),
	\end{align}
which it is defined on the open set
	\begin{align*}
		\mathcal O_R
		=
		\left\{
		(r,u,\theta)\in(0,\infty)\times\mathbb R\times\mathbb R:
		\left(r^{-1}\sin\theta,\cos\theta\right)\in\mathcal V
		\right\}
	\end{align*}
	and it contains the point $\left(R,0,\frac{\pi}{2}\right)$. Moreover, since the map
	$(r,u,\theta)\mapsto\left(r^{-1}\sin\theta,\cos\theta\right)$ is of class $\mathcal C^1$ for $r>0$, we have that the map in \eqref{map for local necks} is $\mathcal C^1$ on
	$\mathcal V$.  Hence the right-hand side of  \eqref{eq:local-neck-phase-system} is of class $\mathcal C^1$ on
	$\mathcal O_R$, and therefore locally Lipschitz in the phase variables near
	$\left(R,0,\frac{\pi}{2}\right)$. Now, we used  Picard--Lindelöf theorem
	\cite[Theorem~2.2]{teschl2012ordinary} to obtain $\varepsilon_R>0$ and a
	unique solution of the initial value problem
	\begin{align*}
		\begin{cases}
			\dot r=\cos\theta,\quad 
			\dot u=\sin\theta,\quad\dot\theta=
			g\left(r^{-1}\sin\theta,\cos\theta\right),
			\\
			r(0)=R,\quad u(0)=0,\quad \theta(0)=\frac{\pi}{2},
		\end{cases}, \mbox{ on }(-\varepsilon_R,\varepsilon_R).
	\end{align*}
    After decreasing $\varepsilon_R$, the
	solution remains in $\mathcal O_R$, and we denote it by
	$(r_R,u_R,\theta_R)$. Then, by evaluating  in the upper-third equation at $s=0$ toguether with the zero trace, we obtain $\dot\theta_R(0)=g(R^{-1},0)=-\rho_0R^{-1}$. Furthermore, since $\dot r_R=\cos\theta_R$, we have $\ddot r_R=-\sin\theta_R\,\dot\theta_R$. In particular, at $s=0$, $\theta_R(0)=\frac{\pi}{2}$, and therefore $\ddot r_R(0)=\rho_0R^{-1}>0$.  Consequently, we have that $r_R$ has a strict local minimum at the neck.
	
	\medskip
	
	The graphical description follows from the Taylor expansion.  Because
	$\dot r_R(0)=\cos\left(\frac{\pi}{2}\right)=0$ and
	$\dot u_R(0)=\sin\left(\frac{\pi}{2}\right)=1$, Taylor's formula gives
	\begin{align*}
		r_R(s)
		=
		R
		+
		\frac{\rho_0}{2R}s^2
		+
		o(s^2),
		\qquad
		\dot r_R(s)
		=
		\frac{\rho_0}{R}s+o(s),
		\qquad
		u_R(s)
		=
		s+o(s).
	\end{align*}
	After decreasing $\varepsilon_R$ if necessary, we have
	$\dot r_R(s)>0$ for $0<s<\varepsilon_R$ and
	$\dot r_R(s)<0$ for $-\varepsilon_R<s<0$.  Hence $r_R(s)>R$ for
	$0<|s|<\varepsilon_R$, and each side of the neck can be written as a
	vertical graph over $r>R$ close to $R$.  These two local graphs are denoted
	by $u_R^+$ and $u_R^-$.
\end{proof}

\begin{figure}[htbp]
	\centering
	\begin{tikzpicture}[scale=0.6, >=Latex]
		\draw[->] (-0.3,0) -- (5.8,0) node[right] {$r$};
		\draw[->] (0,-2.5) -- (0,2.5) node[above] {$u$};
		
		\draw[dashed] (1.0,-2.2) -- (1.0,2.2);
		\draw[dashed] (1.8,-2.2) -- (1.8,2.2);
		\draw[dashed] (2.9,-2.2) -- (2.9,2.2);
		
		\node[below] at (1.0,0) {$R_1$};
		\node[below] at (1.8,0) {$R_2$};
		\node[below] at (2.9,0) {$R_3$};
		
		\draw[thick]
		plot[domain=-1.8:1.8, samples=120, variable=\t]
		({1.0 + 0.22*(0.5*(exp(\t)+exp(-\t))-1)},{\t});
		
		\draw[thick]
		plot[domain=-1.8:1.8, samples=120, variable=\t]
		({1.8 + 0.34*(0.5*(exp(\t)+exp(-\t))-1)},{\t});
		
		\draw[thick]
		plot[domain=-1.8:1.8, samples=120, variable=\t]
		({2.9 + 0.50*(0.5*(exp(\t)+exp(-\t))-1)},{\t});
		
		\draw[thick,->] (1.8,-0.7) -- (1.8,0.8);
		
		\node[right] at (2.25,1.15) {$u_{R_2}^+$};
		\node[right] at (2.25,-1.15) {$u_{R_2}^-$};
				
		\node[align=left, anchor=west] at (3.65,-1.7)
		{$\ddot r_{R_i}(0)=\rho_0R_i^{-1}>0$\\[5pt]
			$(i=1,2,3)$};
	\end{tikzpicture}
	\caption{\footnotesize Local neck profiles for different values of the neck radius $R$.
		Each profile has a strict radial minimum at $r=R$, and the conical zero
		trace gives $\dot\theta_R(0)=-\rho_0R^{-1}$, hence
		$\ddot r_R(0)=\rho_0R^{-1}$.}
	\label{fig:local-catenoidal-neck}
\end{figure}

Next, we pass from the lower side to its reoriented graphical description.  For
$s>0$ we set $\overline{\mathtt c}_R(s)=\mathtt c_R(-s)$ with $\overline\theta_R(s)=\theta_R(-s)+\pi$. Then $\dot{\overline\theta}_R(s)=-\dot\theta_R(-s)$, $\sin\overline\theta_R(s)=-\sin\theta_R(-s)$, and $\cos\overline\theta_R(s)=-\cos\theta_R(-s)$. Thus the phase variables $x=\dot\theta$, $y=r^{-1}\sin\theta$, and
$z=\cos\theta$ transform under reorientation as $(\overline x,\overline y,\overline z)=(-x,-y,-z)$.

\medskip

We notice that by the oddness condition in the selected signed data, the reflected phase still
satisfies the signed translator equation, $\widehat\gamma(\overline x,\overline y)=\overline z$. Thus, if $x=g(y,z)$ represents the selected sheet, the reoriented lower sheet is
represented by
\begin{align*}
	\overline x
	=
	-g(-\overline y,-\overline z)
	=
	g^\sharp(\overline y,\overline z).
\end{align*}
This is the same selected signed phase diagram read after reversing the
orientation of the lower side. Then we select a connected fixed-level component as
\begin{align*}
	I_-
	\subset
	\{y:(y,-1)\in\mathcal D^\sharp\},
\end{align*}
and we write $g_-(y)=g^\sharp(y,-1)$ for $y\in I_-$.  On every interval where the
reoriented lower side is a vertical graph, with slope $v=\frac{du}{dr}$ and
$Q(r,v)\in I_-$, the lower graphical equation is
\begin{align}
	\label{eq:reoriented-lower-graphical-equation}
	v'
	=
	(1+v^2)^{\beta+1}g_-(Q(r,v)),
	\qquad
	Q(r,v)\in I_-.
\end{align}

\subsection{Upper entry}\,
\label{subsec:upper-entry-sector}

At the upper side, we first describe the branch issued from the neck.  The local initial value
problem fixes the outgoing phase; the role of the selected sheet is to identify
the component through which this phase enters the positive branch.  For a
graphical slope $v$, we write
\begin{align*}
	x_v
	=
	\frac{v'}{(1+v^2)^{\beta+1}},
	\qquad
	y_v
	=
	Q(r,v).
\end{align*}
The upper-entry sector condition from
Definition~\ref{def:upper-entry-sector-condition} says that the selected
angular component contains the sector $-\rho_0y<x<0$, $y>0$.  On the
upper side of the neck, the relevant levels satisfy $y>1$.

\begin{lemma}
	\label{lem:upper-entry-sheet-from-zero-sector}
	Assume that the selected angular data satisfy the upper-entry sector
	condition in Definition~\ref{def:upper-entry-sector-condition}.  Assume also that the signed profile is $\mathcal C^1$ on this
	sector, that $\partial_x\widehat\gamma>0$, and that the one-sided trace on
	the positive boundary agrees with the normalized positive trace: $\lim\limits_{x\to0^-}\widehat\gamma(x,y)=\widetilde\gamma(0,y)=y^\alpha$, for  $y\ge1$. Then the level $\widehat\gamma(x,y)=1$ defines a unique
	$\mathcal C^1$ graph $x=h_+(y)$ for $y>1$, satisfying
	\begin{align*}
		-\rho_0y<h_+(y)<0,
		\mbox{ with }
		\lim\limits_{y\to1^+}h_+(y)=0.
	\end{align*}
\end{lemma}

\begin{proof}
	For a fixed $y>1$, the conical zero trace gives
	\begin{align*}
		\widehat\gamma(-\rho_0y,y)
		=
		y^\alpha\widehat\gamma(-\rho_0,1)
		=
		0,
	\end{align*}
	whereas the one-sided positive trace gives
	\begin{align*}
		\lim\limits_{x\to0^-}\widehat\gamma(x,y)
		=
		y^\alpha
		>
		1.
	\end{align*}
	The function $x\mapsto\widehat\gamma(x,y)$ is continuous and strictly
	increasing on the sector, so the level $1$ is crossed exactly once.  This
	defines $h_+(y)\in(-\rho_0y,0)$ with
	$\widehat\gamma(h_+(y),y)=1$.  The inequality
	$\partial_x\widehat\gamma>0$ gives the local implicit-function structure, and
	uniqueness patches the local graphs.  Hence $h_+$ is of class $\mathcal C^1$.
	
	\medskip
	
	Let $y_j\to1^+$.  Because $-\rho_0y_j<h_+(y_j)<0$, after passing to a
	subsequence we may assume that $h_+(y_j)\to x_\infty$, with
	$x_\infty\in[-\rho_0,0]$.  Passing to the limit in the level equation gives
	$\widehat\gamma(x_\infty,1)=1$.  The selected zero ray gives value $0$ at
	$x=-\rho_0$, while the positive boundary trace gives value $1$ at $x=0$.
	Strict monotonicity in the interior of the sector therefore forces
	$x_\infty=0$.  Thus $h_+(y)\to0$ as $y\to1^+$.
\end{proof}

\begin{remark}
We notice that the sheet $x=h_+(y)$ is not the positive branch $g_+$.  It is the
upper-entry sheet of the signed phase diagram, lying in the sector
$-\rho_0y<x<0$ for $y>1$.  The positive branch $g_+$ is reached only after
the phase crosses the endpoint $(0,1)$ and enters $I_+$.
\end{remark}

\begin{lemma}
	\label{lem:upper-branch-enters-positive-regime}
	Let $v$ be a positive graphical branch and set $y_v=Q(r,v)$.  Assume that,
	while $y_v>1$, the phase point lies on the upper-entry sheet of
	Lemma~\ref{lem:upper-entry-sheet-from-zero-sector}.  Then $y_v'(r)<-r^{-1}y_v(r)<0$, as long as $y_v>1$. In particular, if $y_v(r)\to+\infty$ at the neck side, then the branch
	reaches the level $y_v=1$ at a finite radius.
\end{lemma}

\begin{proof}
	On the upper-entry sheet,
	\begin{align*}
		x_v
		=
		\frac{v'}{(1+v^2)^{\beta+1}}
		=
		h_+(y_v)
		<
		0
		\mbox{ while }y_v>1.
	\end{align*}
	Using \eqref{eq:eta-derivative-general}, we obtain
	\begin{align*}
		y_v'
		=
		\frac{1+(1-2\beta)v^2}
		{r(1+v^2)^{\beta+1}}v'
		-
		\frac{y_v}{r}.
	\end{align*}
	Because $1-2\beta=\alpha{-1}>0$ and $x_v<0$, the first term is
	negative.  Thus
	\begin{align*}
		y_v'<-r^{-1}y_v<0\mbox{ while }y_v>1.
	\end{align*}
	
	Finally, for a fixed $r_0$ with $y_v(r_0)>1$, integration of
	$\frac{d}{dr}\log y_v(r)<-r^{-1}$ gives
	$y_v(r)<y_v(r_0)r_0r^{-1}$ as long as $y_v>1$.  Hence every level
	$m\in(1,y_v(r_0))$ is reached before
	$r=y_v(r_0)r_0m^{-1}$.  Letting $m\downarrow1$ and using the monotonicity
	of $y_v$, we obtain a finite radius at which $y_v$ reaches $1$.
\end{proof}

\begin{proposition}
	\label{prop:upper-side-global-continuation}
	Assume the hypotheses of
	Lemma~\ref{lem:upper-entry-sheet-from-zero-sector}.  Then the upper side of
	each local neck with $R\in\mathcal I_0$ reaches the positive branch at a
	finite radius.  More precisely, if $v_R^+=(u_R^+)'$, then there exists
	$r_R^{\mathrm{ent}}>R$ such that
	\begin{align*}
		Q(r,v_R^+(r))\in I_+
		\qquad
		\text{for every }r\ge r_R^{\mathrm{ent}}.
	\end{align*}
	From that radius onward, the upper side has the global positive-branch
	continuation described in
	Proposition~\ref{prop:positive-end-leading-order}.  When the hypotheses of
	Theorem~\ref{thm:positive-bowl-asymptotics} hold, its slope has the
	corresponding positive-end asymptotics.  Consequently, its height has the
	corresponding bowl-type expansion up to an additive constant
	$C_R^+$ determined by the entry data at $r_R^{\mathrm{ent}}$.
\end{proposition}

\begin{proof}
	Near the upper side of the neck, the local graphical slope satisfies
	$v_R^+(r)\to+\infty$ as $r\downarrow R$.  Since $R>0$, the large-slope
	relation for $Q$ gives $Q(r,v_R^+(r))\to+\infty$ as $r\downarrow R$.
	While the phase remains in the upper-entry sector, the level
	$\widehat\gamma=1$ is represented by the sheet $x=h_+(y)$ from
	Lemma~\ref{lem:upper-entry-sheet-from-zero-sector}.  Hence
	Lemma~\ref{lem:upper-branch-enters-positive-regime} applies and shows that
	$y_v=Q(r,v_R^+(r))$ reaches the level $1$ at a finite radius.
	
	\medskip
	
	At the entry point, the upper-entry sheet closes at the common phase point
	$(0,1)$.  For larger radii the phase is read in the positive branch
	$x=g_+(y)$ on $I_+$.  After increasing the entry radius slightly, if
	necessary, the upper side solves
	\begin{align*}
		(v_R^+)'
		=
		(1+(v_R^+)^2)^{\beta+1}
		g_+(Q(r,v_R^+(r))),
		\qquad
		Q(r,v_R^+(r))\in I_+.
	\end{align*}
	Proposition~\ref{prop:positive-end-leading-order} gives the global
	positive-branch continuation and the leading order.  The refined slope
	asymptotics follow from Theorem~\ref{thm:positive-bowl-asymptotics} under its
	hypotheses.
	
	\medskip
	
	Integrating the slope from the entry radius gives the corresponding height
	expansion.  The integration introduces only an additive constant, fixed by the
	entry data at $r_R^{\mathrm{ent}}$.  Hence there is a constant $C_R^+$, depending
	on the entry radius and on the height at entry, such that the upper height
	$u_R^+$ is asymptotic to the corresponding bowl-type height plus $C_R^+$.
\end{proof}

\subsection{Lower entry mechanisms}\,
\label{subsec:lower-entry-mechanisms}

Next, we record the lower mechanisms used in the classification.  The reflected
positive-entry mechanism captures the case where the reoriented lower side
enters $I_+$ in finite radius.  The conical mechanism captures an infinite
radial end approaching the selected zero direction.  The finite-output
mechanism is read from a fixed lower component and is controlled by power
barriers.

\begin{lemma}
	\label{lem:reflected-entry-sheet}
	Assume that, in the reoriented lower phase variables $(x,y)$, the selected
	signed profile is defined on the sector $y>1$, $-\rho_0y<x<0$.
	Assume also that $\widehat\gamma$ is of class $\mathcal C^1$ on this sector,
	that $\partial_x\widehat\gamma>0$, and that the traces on the two boundary
	rays satisfy $\widehat\gamma(-\rho_0y,y)=0$, $\lim\limits_{x\to0^-}\widehat\gamma(x,y)=y^\alpha$ for $y>1$. Then the level $\widehat\gamma(x,y)=1$ is represented in this sector by a
	unique graph $x=h_-^\sharp(y)$ for $y>1$, with
	$-\rho_0y<h_-^\sharp(y)<0$.  Moreover, $h_-^\sharp$ is of class
	$\mathcal C^1$, and $\lim\limits_{y\to1^+}h_-^\sharp(y)=0$.
\end{lemma}

\begin{proof}
	For a fixed $y>1$, the two boundary values are
	\begin{align*}
		\widehat\gamma(-\rho_0y,y)=0,
		\qquad
		\lim\limits_{x\to0^-}\widehat\gamma(x,y)=y^\alpha>1.
	\end{align*}
	The function $x\mapsto\widehat\gamma(x,y)$ is continuous and strictly
	increasing on $(-\rho_0y,0)$.  Therefore the level $1$ is crossed exactly
	once, and the crossing point defines
	$h_-^\sharp(y)\in(-\rho_0y,0)$.  The condition
	$\partial_x\widehat\gamma>0$ gives the local implicit-function structure, and
	uniqueness patches the local graphs.  Hence $h_-^\sharp$ is of class
	$\mathcal C^1$.
	
	\medskip
	
	Let $y_j\to1^+$.  Because $-\rho_0y_j<h_-^\sharp(y_j)<0$, after
	passing to a subsequence we may assume
	$h_-^\sharp(y_j)\to x_\infty$, with $x_\infty\in[-\rho_0,0]$.  Passing to
	the limit in the level equation gives $\widehat\gamma(x_\infty,1)=1$.  The
	selected zero ray gives value $0$ at $x=-\rho_0$, while the positive boundary
	trace gives value $1$ at $x=0$.  Strict monotonicity in the interior forces
	$x_\infty=0$.  Thus $h_-^\sharp(y)\to0$ as $y\to1^+$.
\end{proof}

\begin{lemma}
	\label{lem:lower-reflected-positive-entry}
	Let $R\in\mathcal I_0$, and let $v_R^-$ be the reoriented lower graphical
	slope issued from the local neck.  Assume that, on a maximal interval
	$(R,T_R)$ on which $y_{v_R^-}(r)>1$, the phase is represented by the
	reflected entry sheet of Lemma~\ref{lem:reflected-entry-sheet}:
	\begin{align*}
		x_{v_R^-}(r)=h_-^\sharp(y_{v_R^-}(r)).
	\end{align*}
	Assume also that $y_{v_R^-}(r)\to+\infty$ as $r\downarrow R$.  Then
	$T_R<\infty$ and $\lim\limits_{r\to T_R^-}y_{v_R^-}(r)=1$. Consequently, after increasing the entry radius slightly, if necessary, the
	reoriented lower side enters the positive branch $I_+$.  From that radius
	onward it satisfies
	\begin{align}
		\label{eq:lower-reflected-positive-branch-equation}
		(v_R^-)'
		=
		(1+(v_R^-)^2)^{\beta+1}
		g_+(Q(r,v_R^-)),
		\qquad
		Q(r,v_R^-)\in I_+.
	\end{align}
	When the hypotheses of Theorem~\ref{thm:positive-bowl-asymptotics} hold, the
	reoriented lower slope has the corresponding positive-end asymptotics.  Moreover,
	there is a constant $C_R^-$, determined by the lower entry data, such that the
	reoriented lower height is asymptotic to the corresponding bowl-type height plus
	$C_R^-$.
\end{lemma}

\begin{proof}
	On the reflected entry sheet,
	\begin{align*}
		x_{v_R^-}=\frac{(v_R^-)'}{(1+(v_R^-)^2)^{\beta+1}}=h_-^\sharp(y_{v_R^-})<0,\mbox{ whenever }y_{v_R^-}>1.
	\end{align*}
	Using \eqref{eq:eta-derivative-general}, we obtain
	\begin{align*}
		y_{v_R^-}'
		=
		\frac{1+(1-2\beta)(v_R^-)^2}
		{r(1+(v_R^-)^2)^{\beta+1}}
		(v_R^-)'
		-
		\frac{y_{v_R^-}}{r}.
	\end{align*}
	Because $1-2\beta=\alpha^{-1}>0$ and $x_{v_R^-}<0$, the first term is
	negative.  Hence
	\begin{align}
		\label{eq:lower-reflected-eta-decrease}
		y_{v_R^-}'<-r^{-1}y_{v_R^-}<0
		\mbox{ while }y_{v_R^-}>1.
	\end{align}
	
	For a fixed $r_0>R$ with $y_{v_R^-}(r_0)>1$, integrating
	\eqref{eq:lower-reflected-eta-decrease} gives
	$y_{v_R^-}(r)<y_{v_R^-}(r_0)r_0r^{-1}$ as long as
	$y_{v_R^-}>1$.  Thus every level
	$m\in(1,y_{v_R^-}(r_0))$ is reached before
	$r=y_{v_R^-}(r_0)r_0m^{-1}$.  Letting $m\downarrow1$ and using the
	monotonicity of $y_{v_R^-}$ gives a finite limiting radius at which
	$y_{v_R^-}\to1$.
	
	\medskip
	
	The maximal interval cannot end earlier at a level strictly larger than $1$.
	Indeed, if $y_{v_R^-}$ stays in a compact subinterval of $(1,\infty)$ on a
	finite $r$-interval, then $v_R^-$ stays bounded because $Q(r,\cdot)$ is
	strictly increasing and bijective.  The right-hand side of the reflected entry
	equation is locally bounded and locally Lipschitz, since $h_-^\sharp$ is
	$\mathcal C^1$ on compact subintervals of $(1,\infty)$.  The extensibility
	criterion \cite[Corollary~2.15]{teschl2012ordinary} excludes a finite maximal
	endpoint with $y_{v_R^-}>1$.  Therefore the finite endpoint obtained above
	is the first entry level $y=1$.
	
	\medskip
	
	At $y=1$, Lemma~\ref{lem:reflected-entry-sheet} gives
	$h_-^\sharp(y)\to0$, so the phase reaches $(0,1)$.  The positive trace
	compatibility identifies the continuation for larger radii with the positive
	branch $x=g_+(y)$ on $I_+$.  After increasing the entry radius slightly, if
	necessary, \eqref{eq:lower-reflected-positive-branch-equation} holds. At this point the lower reoriented side satisfies the same positive-branch
	equation as a bowl-type end.  Proposition~\ref{prop:positive-end-leading-order}
	gives the global positive-branch continuation, and
	Theorem~\ref{thm:positive-bowl-asymptotics} gives the corresponding refined
	slope asymptotics whenever its hypotheses hold.  Integrating the slope from the
	lower entry radius then gives the height expansion.  The integration introduces
	only an additive constant, fixed by the lower entry data.  Hence the reoriented
	lower height is asymptotic to the corresponding bowl-type height plus a constant
	$C_R^-$.
\end{proof}

Now, we fix the notation for the lower maximal continuation.  Let
$R\in\mathcal I_0$, and let $(I_-,g_-)$ be a selected lower branch.  Once the
reoriented lower side has entered this branch, we write $J_R^-=(R,T_R)$ for the
maximal interval on which it is a vertical graph with slope $v_R^-$ and
satisfies
\begin{align}
	\label{eq:maximal-lower-selected-component}
  (v_R^-)'=(1+(v_R^-)^2)^{\beta+1}g_-(Q(r,v_R^-),\mbox{ with }	Q(r,v_R^-(r))\in I_-.
\end{align}
The endpoint $T_R$ may be finite or infinite.

\medskip

In the finite-output case $I_-^\flat=(\underline y,0)$, with
$\underline y<0$, the reoriented output level belongs to this component
precisely when
\begin{align}
	\label{eq:finite-output-threshold-section-six}
	-R^{-1}\in(\underline y,0)
	\qquad
	\Longleftrightarrow
	\qquad
	R>R_0,
	\qquad
	R_0=-\underline{y}^{-1}.
\end{align}
The threshold $R_0$ is only the entry threshold for this finite-output
component.

\medskip

Subsequently, we threat  the conical zero-origin criterion.  Indeed, for $\chi\in\mathcal X_g$,
recall from Subsection~\ref{subsec:conical-phase-barriers} that
$\lambda_\chi=-\rho_0+\chi$, and that the negative conical barrier
$\varphi_{\chi,-}$ satisfies $x_{\varphi_{\chi,-}}=\lambda_\chi y_{\varphi_{\chi,-}}$, with $\varphi_{\chi,-}<0$. Then, for $\chi$ close to $0$, $\lambda_\chi<0$, and the negative branch is defined
on every sufficiently large tail and satisfies $\varphi_{\chi,-}(r)\to0^-$.

\begin{lemma}
	\label{lem:conical-barriers-zero-origin-return}
	Let $(I_-,g_-)$ be a regular selected lower branch such that
	$0\in\partial I_-$ and $(-\delta_0,0)\subset I_-$ for some $\delta_0>0$.
	Assume that
	\begin{align}
		\label{eq:selected-conical-trace-lower-branch}
		\lim\limits_{y\to0^-}g_-(y)y^{-1}=-\rho_0.
	\end{align}
	Let $v$ be a reoriented lower graphical solution, branch-admissible for
	$(I_-,g_-)$ on an infinite tail $[R_1,\infty)$.  Assume that, for every
	sufficiently small $\varepsilon>0$, if
	$\delta_\varepsilon\in(0,\delta_0)$ is chosen so that
	\begin{align}
		\label{eq:lower-branch-ray-squeezing}
		(-\rho_0+\varepsilon)y
		<
		g_-(y)
		<
		(-\rho_0-\varepsilon)y
		\qquad
		\text{for }-\delta_\varepsilon<y<0,
	\end{align}
	then there is $R_\varepsilon\ge R_1$ such that
	\begin{align}
		\label{eq:lower-level-enters-conical-tail}
		Q(R_\varepsilon,v(R_\varepsilon))
		\in
		(-\delta_\varepsilon,0).
	\end{align}
	Then
	\begin{align*}
		x_v(r)\to0,
		\quad
		y_v(r)\to0,
		\quad
		\frac{x_v(r)+\rho_0y_v(r)}
		{y_v(r)}
		\to0,\mbox{ as }r\to\infty.
	\end{align*}
\end{lemma}

\begin{proof}
	For a sufficiently small $\varepsilon>0$, we use
	$\delta_\varepsilon$ and $R_\varepsilon$ from
	\eqref{eq:lower-branch-ray-squeezing} and
	\eqref{eq:lower-level-enters-conical-tail}.  Set
	$\lambda_{\pm\varepsilon}=-\rho_0\pm\varepsilon$. The first step is to choose the two conical barriers.  We set
	$v_\varepsilon=v(R_\varepsilon)$.  Then, since 
	$Q(R_\varepsilon,v_\varepsilon)\in(-\delta_\varepsilon,0)$ and
	$Q(R_\varepsilon,\cdot)$ is strictly increasing, we may choose
	$A_\varepsilon^-<v_\varepsilon<A_\varepsilon^+<0$ so close to
	$v_\varepsilon$ that
	$Q(R_\varepsilon,A_\varepsilon^\pm)\in(-\delta_\varepsilon,0)$.  For a
	prescribed negative value $A<0$, the explicit formula
	\eqref{eq:explicit-conical-slope-barrier} determines the constant in
	$\varphi_{\chi,-}$ uniquely by
	\begin{align}
		\label{eq:conical-constant-choice-section-six}
		C
		=
		\frac{-A\,R_\varepsilon^{-\lambda_\chi}}
		{\sqrt{1+A^2}}.
	\end{align}
	Choose $\varphi_{\varepsilon,-}$ with
	$\varphi_{\varepsilon,-}(R_\varepsilon)=A_\varepsilon^-$ and
	$\varphi_{-\varepsilon,-}$ with
	$\varphi_{-\varepsilon,-}(R_\varepsilon)=A_\varepsilon^+$.  Thus
	\begin{align}
		\label{eq:conical-barriers-initial-order-section-six}
		\varphi_{\varepsilon,-}(R_\varepsilon)
		<
		v(R_\varepsilon)
		<
		\varphi_{-\varepsilon,-}(R_\varepsilon).
	\end{align}
	
	\medskip
	
	The barriers are admissible for the selected lower branch on
	$[R_\varepsilon,\infty)$.  Differentiating
	$x_{\varphi_{\chi,-}}=\lambda_\chi y_{\varphi_{\chi,-}}$ with
	\eqref{eq:eta-derivative-general} gives
	\begin{align*}
		y_{\varphi_{\chi,-}}'
		=
		\frac{1}{r}
		\left[
		\lambda_\chi
		\bigl(1+(1-2\beta)\varphi_{\chi,-}^2\bigr)
		-
		1
		\right]
		y_{\varphi_{\chi,-}}.
	\end{align*}
	Here $\lambda_\chi<0$ and $y_{\varphi_{\chi,-}}<0$, so
	$y_{\varphi_{\chi,-}}'>0$.  Thus
	$y_{\varphi_{\chi,-}}$ increases to $0$ from below.  Because the initial
	levels lie in $(-\delta_\varepsilon,0)$, the barriers remain branch-admissible
	on the whole tail.
	
	\medskip
	
	For $y\in(-\delta_\varepsilon,0)$, the squeezing inequality gives $\lambda_{\varepsilon}y-g_-(y)<0$, and $\lambda_{-\varepsilon}y-g_-(y)>0$. Then, by using \eqref{eq:conical-barrier-residual}, we get $\mathcal R_-[\varphi_{\varepsilon,-}]<0$, and $\mathcal R_-[\varphi_{-\varepsilon,-}]>0$.
	Consequently, the comparison principle and the initial ordering
	\eqref{eq:conical-barriers-initial-order-section-six} give
	\begin{align*}
		\varphi_{\varepsilon,-}(r)<v(r)<\varphi_{-\varepsilon,-}(r)
		\mbox{ for every }r\ge R_\varepsilon.
	\end{align*}
	So, since $\lambda_{\pm\varepsilon}<0$, both conical barriers tend to $0^-$ as
	$r\to\infty$.  Hence, we obtain $v(r)\to0^-$ and $y_v(r)=Q(r,v(r))\to0^-$.
	
	\medskip
	
	On the selected lower branch, we have $x_v(r)=g_-(y_v(r))$. Then,  we observe that for all large $r$,
	$y_v(r)\in(-\delta_\varepsilon,0)$, and with
	\eqref{eq:lower-branch-ray-squeezing} we have 
	\begin{align*}
		(-\rho_0+\varepsilon)y_v(r)
		<
		x_v(r)
		<
		(-\rho_0-\varepsilon)y_v(r).
	\end{align*}
	Subsequently, by dividing by $y_v(r)$, we have reversed inequalities:
	\begin{align*}
		-\rho_0-\varepsilon<x_v(r)y_v(r)^{-1}<-\rho_0+\varepsilon\Leftrightarrow-\varepsilon<\frac{x_v(r)+\rho_0y_v(r)}{y_v(r)}<\varepsilon.
	\end{align*}
	Letting $\varepsilon\downarrow0$ proves the directional convergence.  The
	same squeezing gives $x_v(r)=O(|y_v(r)|)$, so $x_v(r)\to0$.
\end{proof}


Before proving Theorem~\ref{thm:rotational-translators}, the verification of
the branch alternatives in the phase diagram can be summarized as follows.

\begin{enumerate}[label=\textup{(\alph*)}]
	\item The upper-entry phase is verified by
	Lemma~\ref{lem:upper-entry-sheet-from-zero-sector}.  In the phase variables
	$(x,y)$, the sector $\mathcal S_+(\rho_0)$ must lie in the selected sheet, the trace at $x=0$ must agree with
	$\widetilde\gamma(0,y)=y^\alpha$, and
	$\partial_x\widehat\gamma>0$ must hold on the sector.
	
	\item The reflected positive-entry phase is verified by
	Lemmas~\ref{lem:reflected-entry-sheet} and
	\ref{lem:lower-reflected-positive-entry}.  After reorientation, the level
	$\widehat\gamma=1$ is represented for $y>1$ by a sheet
	$x=h_-^\sharp(y)$ with
	$-\rho_0y<h_-^\sharp(y)<0$ and
	$h_-^\sharp(y)\to0$ as $y\to1^+$.
	
	\item The zero-trace radial branch is verified by
	Definition~\ref{def:zero-origin-return} and
	Lemma~\ref{lem:conical-barriers-zero-origin-return}.  At the branch level,
	the selected fixed-level lower branch satisfies $\lim\limits_{y\to0^-}g_-(y)y^{-1}=-\rho_0$. For the associated solution, the conical barriers trap the graphical phase
	$(x_{v_R^-},y_{v_R^-})$ between rays converging to
	$x=-\rho_0y$.
	
	\item The flat complete fixed-level branch is verified by the branch
	condition in Definition~\ref{def:lower-continuation-outcomes}.  When
	$I_-^\flat=(\underline y,0)$, with $\underline y<0$, then the lower output
	level lies in this component exactly when $-R^{-1}\in(\underline y,0)$, or equivalently,  $R>R_0=-\underline{y}^{-1}$. The expansion of $g_-$ at $0$ provides the power scale used in
	Proposition~\ref{prop:power-barrier-flat-component}.
	
	\item When the selected lower branch is a normalized endpoint branch or a
	pole-limited branch in the sense of
	Definition~\ref{def:lower-continuation-outcomes}, the continuation is
	maximal inside the selected component.  The endpoint is determined by the
	boundary behavior of that component.
\end{enumerate}

\begin{lemma}
	\label{lem:embeddedness-from-slope-comparison}
	Let $\mathtt c_R$ be a catenoidal profile obtained by continuing the local
	neck of Proposition~\ref{prop:local-catenoidal-neck}.  Assume that the two
	sides are written as vertical graphs
	\begin{align*}
		u_R^+:(R,T^+)\to\mathbb R,
		\qquad
		u_R^-:(R,T^-)\to\mathbb R,
	\end{align*}
	with slopes $v_R^+=(u_R^+)'$ and $v_R^-=(u_R^-)'$, and with the common neck
	value $\lim\limits_{r\downarrow R}u_R^+(r)=\lim\limits_{r\downarrow R}u_R^-(r)=0$. Set $T=\min{T^+,T^-}$.  Assume that one of the following comparison mechanisms holds on the common radial interval.
	
	\begin{enumerate}[label=\textup{(\alph*)}]
		\item \label{emb:sign-separated}
		The slopes are sign-separated: $v_R^-(r)\le0<v_R^+(r)$, for every $r\in(R,T)$.
		
		\item \label{emb:ordered-branch}
		There exists $r_0\in(R,T)$ such that
		$u_R^+(r)>u_R^-(r)$ for $R<r\le r_0$ and
		$v_R^-(r_0)\le v_R^+(r_0)$.  Moreover, there is a regular selected branch
		$(I_\sigma,g_\sigma)$ such that $v_R^-$ and $v_R^+$ are branch-admissible
		for $(I_\sigma,g_\sigma)$ on $[r_0,T)$, and $\mathcal R_\sigma[v_R^-]\le0\le\mathcal R_\sigma[v_R^+]$ on $[r_0,T)$.
	\end{enumerate}
	Then $u_R^+(r)>u_R^-(r)$ for every $r\in(R,T)$. Consequently the generating curve has no self-intersection on its common
	radial domain.  In each complete alternative of
	Theorem~\ref{thm:rotational-translators} in which the relevant comparison
	mechanism applies on the whole common radial domain, the corresponding
	rotational catenoidal translator is embedded.	
\end{lemma}

\begin{proof}
	By Proposition~\ref{prop:local-catenoidal-neck}, the neck is a strict radial
	minimum and the two local sides are vertical graphs.  The Taylor expansion at
	the neck gives $u_R^+(r)>u_R^-(r)$ for $r>R$ sufficiently close to $R$. In case~\ref{emb:sign-separated}, we have $(u_R^+-u_R^-)'(r)=v_R^+(r)-v_R^-(r)>0$ on $(R,T)$. Then, the strict separation near the neck therefore persists on the whole common interval. In case~\ref{emb:ordered-branch}, the residual inequalities are precisely the
	hypotheses needed to compare the two slopes on the selected branch.  Applying
	Proposition~\ref{prop:comparison-principle-ode} with the initial inequality
	$v_R^-(r_0)\le v_R^+(r_0)$ gives $v_R^-(r)\le v_R^+(r)$ for every $r\in[r_0,T)$.
	Hence $(u_R^+-u_R^-)'\ge0$ on $[r_0,T)$.  Then, since $u_R^+(r_0)>u_R^-(r_0)$, integration gives $u_R^+(r)-u_R^-(r)\ge u_R^+(r_0)-u_R^-(r_0)>0$ for every $r\in[r_0,T)$. Together with the known separation on $(R,r_0]$, this proves the strict vertical separation in both cases.
	
	\medskip
	
	Each side is a vertical graph on its radial interval, so neither side has
	self-intersections.  The strict separation excludes intersections between the
	two sides away from the neck.  The two branches meet only at the common neck.
	Thus the generating curve is embedded on its common radial domain.  Rotating
	an embedded generating curve with $r>0$ away from the neck gives an embedded
	rotational hypersurface.

	
\end{proof}


\begin{theorem}
	\label{thm:rotational-translators}
	Assume that the selected angular neck data $(A_g,\rho_0)$ are regular and
	that the upper-entry hypotheses of
	Lemma~\ref{lem:upper-entry-sheet-from-zero-sector} hold.  Let
	$R\in\mathcal I_0$ be an admissible neck radius.  Then the zero trace
	$g(y,0)=-\rho_0y$ produces a local catenoidal neck
	$\mathtt c_R=(r_R,u_R)$ of radius $R$.  Its upper side reaches the positive
	branch $I_+$ at finite radius and then continues as a positive bowl-type
	end.  In the regimes covered by
	Theorem~\ref{thm:positive-bowl-asymptotics}, this end has the corresponding
	positive-end asymptotics.
	
	After reorienting the lower side, the continuation is determined by the
	selected implicit branch followed by the lower phase.  The following
	alternatives exhaust the selected continuations.
	
	\begin{enumerate}[label=\textup{(\roman*)}]
		\item  The reoriented lower phase follows the reflected entry sheet of
		Lemma~\ref{lem:lower-reflected-positive-entry}.  Equivalently, it reaches
		the common endpoint $(0,1)$ of the reflected signed level and the
		positive branch.  Hence the lower side enters $I_+$ at finite radius and
		continues, from that radius on, as a second positive bowl-type end.
		Together with the upper side, it forms a complete embedded catenoidal
		translator $W_R$ with two positive bowl-type ends.
		
		\item The reoriented lower phase is carried by a selected fixed-level lower
		branch $(I_-,g_-)$ in the sense of Definition~\ref{def:selected-lower-branch}. The lower end is then determined by the type of $(I_-,g_-)$, as recorded
		in Definition~\ref{def:lower-continuation-outcomes}, with the following
		geometric consequences:
		
		\begin{enumerate}[label=\textup{(\alph*)}]
			\item When the selected fixed-level lower branch has the zero-trace endpoint
			behavior of Definition~\ref{def:zero-origin-return}, then the lower
			phase approaches the origin with direction asymptotic to the selected
			zero ray $x=-\rho_0y$.  The lower side is complete as a radial end,
			and together with the upper side it forms a complete embedded
			catenoidal translator $W_R$ with one positive bowl-type end and one
			zero-trace radial end.
			
			\item  When $(I_-,g_-)$ is a flat complete fixed-level branch in the sense of
			Definition~\ref{def:lower-continuation-outcomes}, then
			$I_-=I_-^\flat=(\underline y,0)$ with $\underline y<0$.  The lower
			output level belongs to this component precisely when $-R^{-1}\in(\underline y,0)$, or equivalently $R>R_0=\-\underline {y}^{-1}$. For such radii, the lower side is complete on $I_-^\flat$, and the
			resulting surface is a complete embedded catenoidal translator $W_R$.
			
			\item When $(I_-,g_-)$ is a normalized endpoint branch or a pole-limited
			branch in the sense of
			Definition~\ref{def:lower-continuation-outcomes}, then the lower
			continuation reaches the boundary of the selected lower component.
			These cases give maximal admissible rotational pieces inside the
			selected signed component, rather than complete catenoidal
			translators.
		\end{enumerate}
	\end{enumerate}
\end{theorem}



\begin{proof}[Proof of Theorem~\ref{thm:rotational-translators}]
	The local neck is given by
	Proposition~\ref{prop:local-catenoidal-neck}.  The equality
	$g(y,0)=-\rho_0y$ gives $\dot\theta_R(0)=-\rho_0R^{-1}$,
	and $\ddot r_R(0)=\rho_0R^{-1}>0$, and hence the local profile has a strict neck of radius $R$.  The upper side
	enters the positive branch by
	Lemma~\ref{lem:upper-entry-sheet-from-zero-sector} and
	Proposition~\ref{prop:upper-side-global-continuation}.  Once it reaches
	$I_+$, it is governed by the positive branch equation.  The refined
	asymptotics follow from
	Theorem~\ref{thm:positive-bowl-asymptotics} under the hypotheses of that
	theorem.
	
	\medskip
	
	The reoriented lower side is considered next.  When the lower phase follows the
	reflected entry sheet, then
	Lemma~\ref{lem:lower-reflected-positive-entry} shows that it reaches
	$I_+$ at finite radius.  The continuation is then governed by the positive
	branch equation and gives a second positive bowl-type end.  The two sides
	are separated by the common-branch comparison in
	Lemma~\ref{lem:embeddedness-from-slope-comparison}\textup{(b)}, and
	therefore form a complete embedded catenoidal translator with two positive
	bowl-type ends.
	
	\medskip
	
	The remaining case is the one in which the reoriented lower phase is
	carried by a selected fixed-level lower branch $(I_-,g_-)$:
	\begin{enumerate}
		\item 	In the zero-trace endpoint case, we assume that this branch has the behavior in
		Definition~\ref{def:zero-origin-return}.  The conical barriers of
		Lemma~\ref{lem:conical-barriers-zero-origin-return} trap the graphical phase
		between rays converging to $x=-\rho_0y$.  Hence $x_{v_R^-},y_{v_R^{-}}\to0$,
		and $\frac{x_{v_R^-}(r)+\rho_0y_{v_R^-}(r)}{y_{v_R^-}(r)}\to0$. The lower graph is defined on an infinite radial interval and $v_R^-(r)\to0^-$.  Thus the lower side is complete as a radial end. The	sign separation of the two graphical slopes places the profile under Lemma~\ref{lem:embeddedness-from-slope-comparison}\textup{(a)}, and the two	 sides form a complete embedded catenoidal translator with one positive
		bowl-type end and one zero-trace radial end.
		
		\item In the flat complete fixed-level case, we assume that $(I_-,g_-)$ is a branch
		$I_-^\flat=(\underline y,0)$, with $\underline y<0$.  The lower output level
		belongs to this component exactly when $-R^{-1}\in(\underline y,0)$, or equivalently $R>R_0=-\underline {y}^{-1}$. For such radii, the power-barrier signs in Proposition~\ref{prop:power-barrier-flat-component} trap the lower slope on its maximal lower interval.  The maximal endpoint cannot be finite: on
		finite intervals the comparison bounds keep $v_R^-$ bounded, and near the
		endpoint $0$ the expansion $g_-(y)=\ell y+O(|y|^{1+\mu})$ as $y\to0^-$ gives a continuous extension of the right-hand side of the lower graphical
		equation.  The extensibility criterion \cite[Corollary~2.15]{teschl2012ordinary} excludes a finite maximal
		endpoint.  Hence the lower side is defined on an infinite radial interval,
		and the power barriers give $Q(r,v_R^-(r))\to0^- $.    Again the sign separation of the graphical slopes places the profile under
		Lemma~\ref{lem:embeddedness-from-slope-comparison}\textup{(a)}, and the two  sides form a complete embedded catenoidal translator.
		\item In the normalized endpoint case, we assume that the selected lower branch is a
		branch or a pole-limited branch.  By
		Definition~\ref{def:lower-continuation-outcomes}, the selected continuation
		reaches the boundary of the chosen lower component.  The solution is then
		maximal inside that selected signed phase component.  These alternatives
		produce maximal admissible rotational pieces rather than complete
		catenoidal translators.
	\end{enumerate}
	
\end{proof}


\subsection{Model catenoidal profiles}\,
\label{subsec:model-catenoidal-profiles}

We close the classification section with schematic profiles for three model
situations appearing in the signed phase analysis.  The first family is given by
the adjacent Hessian quotients $\mathcal Q_{k,k-1}=\frac{S_k}{S_{k-1}}$ in
dimension $n=6$.  For $k=3,4,5$, the reflected fixed level has a flat complete
component $I_-^\flat=(-\tau_k,0)$, and the choice $R=3$ gives
$-R^{-1}\in I_-^\flat$ in all three cases.  The second family is given by the
odd-root speeds $S_k^{\frac{1}{k}}$, with $k$ odd; here the selected reflected
component is a normalized endpoint component.  The third example is the signed
fully nonlinear model $\left(4+\frac{24}{10}\right)^{-1}\left(S_3+\frac{1}{10}HS_2\right)$ with $n=5$, for which the selected regular sheet contains the upper-entry sector and the reflected positive-entry mechanism is available.

\begin{figure}[htbp]
	\centering
	
	\begin{subfigure}[t]{0.32\textwidth}
		\centering
		\begin{tikzpicture}
			\begin{axis}[
				width=\linewidth,
				height=0.78\linewidth,
				axis lines=middle,
				xmin=0, xmax=10.8,
				ymin=-4.1, ymax=5.1,
				xlabel={$r$},
				ylabel={$u$},
				xlabel style={font=\scriptsize,yshift=-2pt},
				ylabel style={font=\scriptsize},
				tick label style={font=\scriptsize},
				label style={font=\scriptsize},
				samples=140,
				clip=false
				]
				\addplot[gray, dashed] coordinates {(3,-4.1) (3,5.0)};
				\node[font=\scriptsize, anchor=north] at (axis cs:3,-4.1) {$R=3$};
				
				\addplot[black, thick, domain=3:10.3]
				{0.55*sqrt(x-3)+0.12*(x-3)^1.55};
				
				\addplot[blue, thick, domain=3:10.3]
				{-1.15*(1-exp(-0.62*sqrt(x-3)))};
				
				\addplot[red!75!black, thick, domain=3:10.3]
				{-0.78*ln(1+sqrt(x-3))};
				
				\addplot[orange!90!black, thick, domain=3:10.3]
				{-0.58*((1+x-3)^0.55-1)};
				
				\fill (axis cs:3,0) circle (1.3pt);
				\node[font=\scriptsize, anchor=south west] at (axis cs:3,0) {neck};
				
				\node[font=\scriptsize, anchor=west] at (axis cs:6.5,3.25)
				{$u_R^+$};
				\node[font=\tiny, anchor=west] at (axis cs:5.8,-1.05)
				{$\mathcal Q_{3,2}$};
				\node[font=\tiny, anchor=west] at (axis cs:5.8,-1.85)
				{$\mathcal Q_{4,3}$};
				\node[font=\tiny, anchor=west] at (axis cs:5.8,-2.70)
				{$\mathcal Q_{5,4}$};
			\end{axis}
		\end{tikzpicture}
		\caption{$\mathcal Q_{k,k-1}$: flat complete lower components.}
	\end{subfigure}
	\hfill
	\begin{subfigure}[t]{0.32\textwidth}
		\centering
		\begin{tikzpicture}
			\begin{axis}[
				width=\linewidth,
				height=0.78\linewidth,
				axis lines=middle,
				xmin=0, xmax=6.3,
				ymin=-3.2, ymax=4.2,
				xlabel={$r$},
				ylabel={$u$},
				xlabel style={font=\scriptsize,yshift=-2pt},
				ylabel style={font=\scriptsize},
				tick label style={font=\scriptsize},
				label style={font=\scriptsize},
				clip=false
				]
				\addplot[gray, dashed] coordinates {(1,-3.2) (1,4.1)};
				\node[font=\scriptsize, anchor=north] at (axis cs:1,-3.2) {$R=1$};
				
				\addplot[black, thick, domain=1:6.0, samples=140]
				{0.60*sqrt(x-1)+0.16*(x-1)^1.52};
				
				\addplot[blue, thick, smooth]
				coordinates {
					(1.00, 0.00)
					(1.22,-0.42)
					(1.55,-0.90)
					(1.95,-1.34)
					(2.35,-1.68)
					(2.70,-1.92)
					(2.95,-2.03)
					(3.10,-2.06)
				};
				
				\fill (axis cs:1,0) circle (1.3pt);
				\node[font=\scriptsize, anchor=south west] at (axis cs:1,0) {neck};
				
				\fill[blue] (axis cs:3.10,-2.06) circle (1.2pt);
				
				\node[font=\scriptsize, anchor=west] at (axis cs:4.15,3.10)
				{$u_R^+$};
				\node[font=\tiny, anchor=west] at (axis cs:3.15,-2.15)
				{endpoint};
			\end{axis}
		\end{tikzpicture}
		\caption{$S_k^{\frac{1}{k}}$: normalized endpoint component.}
	\end{subfigure}
	\hfill
	\begin{subfigure}[t]{0.32\textwidth}
		\centering
		\begin{tikzpicture}
			\begin{axis}[
				width=\linewidth,
				height=0.78\linewidth,
				axis lines=middle,
				xmin=0, xmax=8.8,
				ymin=-2.6, ymax=5.1,
				xlabel={$r$},
				ylabel={$u$},
				xlabel style={font=\scriptsize,yshift=-2pt},
				ylabel style={font=\scriptsize},
				tick label style={font=\scriptsize},
				label style={font=\scriptsize},
				clip=false
				]
				\addplot[gray, dashed] coordinates {(1.6,-2.6) (1.6,5.0)};
				\node[font=\scriptsize, anchor=north] at (axis cs:1.6,-2.6) {$R$};
				
				\addplot[black, thick, smooth]
				coordinates {
					(1.60, 0.00)
					(1.85, 0.42)
					(2.25, 0.86)
					(2.85, 1.35)
					(3.60, 1.93)
					(4.45, 2.55)
					(5.35, 3.15)
					(6.35, 3.78)
					(7.45, 4.45)
					(8.50, 5.05)
				};
				
				\addplot[blue, thick, smooth]
				coordinates {
					(1.60, 0.00)
					(1.82,-0.48)
					(2.20,-0.96)
					(2.75,-1.32)
					(3.35,-1.46)
					(3.95,-1.34)
					(4.55,-0.94)
					(5.10,-0.35)
					(5.55, 0.25)
					(6.05, 0.95)
					(6.75, 1.85)
					(7.55, 2.90)
					(8.45, 4.05)
				};
				
				\fill (axis cs:1.6,0) circle (1.3pt);
				\node[font=\scriptsize, anchor=south west] at (axis cs:1.6,0) {neck};
				
				\fill[blue] (axis cs:5.42,0) circle (1.2pt);
				\node[font=\tiny, anchor=south west] at (axis cs:5.45,0.05)
				{entry};
				
				\node[font=\scriptsize, anchor=west] at (axis cs:6.2,4.25)
				{$u_R^+$};
				\node[font=\scriptsize, anchor=west] at (axis cs:6.25,2.05)
				{$u_R^-$};
			\end{axis}
		\end{tikzpicture}
		\caption{$\left(4+\frac{24}{10}\right)^{-1}\left(S_3+\frac{1}{10}HS_2\right)$: reflected positive entry.}
	\end{subfigure}
	
	\caption{\footnotesize Schematic catenoidal profiles for the model curvature functions used
		in the signed phase analysis.  The panels represent, respectively, flat
		complete lower components for adjacent Hessian quotients, a normalized
		endpoint component for odd-root speeds, and a reflected positive-entry
		mechanism for a signed fully nonlinear model.}
	\label{fig:model-catenoidal-profiles}
\end{figure}


%	\begin{table}[htbp]
	%		\centering
	%		\scriptsize
	%		\setlength{\tabcolsep}{3pt}
	%		\renewcommand{\arraystretch}{1.25}
	%		\begin{tabular}{p{0.20\textwidth}p{0.30\textwidth}p{0.16\textwidth}p{0.18\textwidth}}
		%			\hline
		%			\textbf{Lower phase mechanism} &
		%			\textbf{Phase condition} &
		%			\textbf{Neck radii} &
		%			\textbf{Geometric outcome} \\
		%			\hline
		%			Reflected positive entry &
		%			The reoriented lower phase is represented by a reflected entry sheet
		%			$x=h_-^\sharp(y)$ for $y>1$ and satisfies
		%			Lemma~\ref{lem:lower-reflected-positive-entry}. &
		%			Radii for which the reoriented exit lies in the reflected entry sector. &
		%			Complete catenoidal translator with two positive bowl-type ends. \\
		%			\hline
		%			Selected zero-origin asymptotic &
		%			The lower branch satisfies
		%			$\lim\limits_{y\to0^-}g_-(y)y^{-1}=-\rho_0$, and the negative conical
		%			barriers $\varphi_{\chi,-}$ apply on an infinite tail. &
		%			Radii for which the lower solution enters the selected conical tail. &
		%			Complete translator with upper bowl end and selected zero-origin
		%			asymptotic lower end. \\
		%			\hline
		%			Flat complete fixed-level component &
		%			A fixed lower component satisfies $I_-=(\underline y,0)$,
		%			$g_-(y)=\ell y+O(|y|^{1+\mu})$ as $y\to0^-$, and the power barriers trap
		%			the lower slope. &
		%			Those $R$ such that $-R^{-1}\in I_-$. &
		%			Complete translator with upper bowl end and complete fixed-level lower
		%			end. \\
		%			\hline
		%			Normalized endpoint component &
		%			A fixed lower component satisfies $I_-=(-\infty,-1)$ and
		%			$\lim\limits_{y\to-1^-}g_-(y)=0$. &
		%			Typically $0<R<1$ in normalized examples. &
		%			Maximal selected-component piece relative to the fixed-level chart; loss
		%			of admissibility occurs at $Q=-1$. \\
		%			\hline
		%			Pole-limited component &
		%			A fixed lower component satisfies $I_-=(-1,\lambda_+)$, $-1<\lambda_+<0$,
		%			and $\lim\limits_{y\to\lambda_+^-}|y^{-1}g_-(y)|=\infty$. &
		%			Typically $1<R<R_0$, where $R_0=(-\lambda_+)^{-1}$. &
		%			Maximal selected-component piece with finite-radius boundary or singular
		%			endpoint, relative to the fixed-level chart. \\
		%			\hline
		%		\end{tabular}
	%		\caption{Summary of the lower phase mechanisms and their geometric outcomes.
		%			The phase data are defined in Section~\ref{sec:curvature-functions-branches};
		%			the barrier method proving the corresponding alternatives is developed in
		%			Section~\ref{sec:barrier-method}.}
	%		\label{tab:lower-phase-classification}
	%	\end{table}




	
	


\section{Uniqueness and catenoidal barriers}
	\label{sec:uniqueness-catenoidal-barriers}
	
	In this section, we prove two applications of the rotational theory developed above.
	The first one is a uniqueness theorem for strictly convex entire graphical
	$\gamma$-translators which are asymptotic to a rotational bowl-type translator.
	
	\medskip
	
	The second application concerns graphical translators over bounded domains.  We prove that a convex graphical $\gamma$-translator cannot be
	defined over a bounded domain with vertical boundary blow-up, provided that the
	catenoidal classification supplies a complete embedded family whose upper
	graphical sides enter the positive branch at radii tending to zero with the neck
	radius.  	
	\subsection{Uniqueness for bowl-type asymptotics}\,
	
	Following \cite{Paco_2014}, we use the moving-plane notation.  For
	$X=(x,x_{n+1})\in\mathbb R^{n+1}$, with
	$x=(x_1,\ldots,x_n)\in\mathbb R^n$, we set $\mathbf p(X)=x_1$ and
	$\Pi_t=\{\mathbf p=t\}$.  For a set $A\subset\mathbb R^{n+1}$, write
	$A_+(t)=A\cap\{\mathbf p>t\}$ and
	$A_-(t)=A\cap\{\mathbf p<t\}$.  Let $\delta_t$ be the reflection across
	$\Pi_t$,
	\begin{align*}
		\delta_t(x_1,x_2,\ldots,x_n,x_{n+1})
		=
		(2t-x_1,x_2,\ldots,x_n,x_{n+1}),
	\end{align*}
	and set $A_+^*(t)=\delta_t(A_+(t))$.  The vertical projection is denoted by
	$\pi:\mathbb R^{n+1}\to\mathbb R^n$.
	
	\medskip
	
	Let $u_B:[0,\infty)\to\mathbb R$ be the radial profile of the model bowl, let
	$v_B=u_B'$, and write $B=\{(x,u_B(|x|)):x\in\mathbb R^n\}$. Let $M=\operatorname{graph}(u)$ be a strictly convex entire graphical
	$\gamma$-translator.  The graph $M$ is said to be $C^2$-asymptotic to a vertical
	translate of $B$ if there exists $c\in\mathbb R$ such that, writing
	$u(x)=u_B(|x|)+c+\varphi(x)$, one has
	\begin{align*}
		\varphi(x)\to0,
		\qquad
		D\varphi(x)\to0,
		\qquad
		D^2\varphi(x)\to0
		\qquad
		\text{as }|x|\to\infty .
	\end{align*}
	After replacing $B$ by $B+c e_{n+1}$, we assume from now on that $c=0$, so
	\begin{align}
		\label{eq:C2-asymptotic-to-bowl}
		u(x)=u_B(|x|)+\varphi(x),
		\qquad
		\varphi(x)\to0,
		\qquad
		D\varphi(x)\to0,
		\qquad
		D^2\varphi(x)\to0
	\end{align}
	as $|x|\to\infty$.
	
	\begin{remark}
		\label{rem:C2-asymptotic-moving-plane}
		The $C^2$-asymptotic assumption is used only in the moving-plane
		application.  It gives uniform control of the height, gradient, and
		second-order data of $M$ relative to the prescribed bowl.  In particular,
		the gradient estimate is what allows the asymptotic separation at infinity
		to be differentiated in the strip part of the reflection argument, while
		the second-order control keeps both graphs in the same elliptic
		$\Gamma$-admissible regime at finite contact points.
	\end{remark}
	
	For each $t$, the reflected cap $M_+^*(t)$ is a vertical graph.  Indeed,
	$\delta_t$ preserves the vertical direction, and on $\{x_1<t\}$ we may write
	\begin{align*}
		M_+^*(t)=\{(x,u_t^*(x)):x_1<t\},
		\qquad
		u_t^*(x)=u(2t-x_1,x_2,\ldots,x_n).
	\end{align*}
	Thus the issue is the vertical ordering between $u_t^*$ and $u$.
	
	\begin{lemma}
		\label{lem:start-moving-plane-at-infinity}
		Assume that the positive end of $B$ is in one of the following regimes:
		\begin{enumerate}[label=\textup{(\roman*)}]
			\item the normalized $1$-nondegenerate case with $\alpha\ge1$ and
			\begin{align*}
				v_B(r)=r^\alpha-\frac{a}{r^\alpha}+\frac{b}{r^{3\alpha}}+o(r^{-4\alpha});
			\end{align*}
			
			\item the $1$-degenerate power-growth case with
			\begin{align*}
				v_B(r)
				=
				A_\gamma
				r^{
					\frac{\alpha(k_\gamma+1)}
					{k_\gamma+1-2\alpha}
				}
				(1+o(1)),
				\mbox{ with }
				\frac{\alpha(k_\gamma+1)}
				{k_\gamma+1-2\alpha}\ge1;
			\end{align*}
			
			\item the critical $1$-degenerate case with
			$\log v_B(r)=\frac{c_\gamma}{2\alpha}r^{2\alpha}+o(r^{2\alpha})$.
		\end{enumerate}
		Then there exists $t_0>0$ such that $u_t^*(x)>u(x)$ for every
		$t\ge t_0$ and every $x_1<t$.  Moreover, for every fixed $t>0$ and every
		$\sigma\in(0,t)$, there exist $R_{\sigma,t}>0$ and $\mu_{\sigma,t}>0$ such
		that $u_t^*(x)-u(x)\ge\mu_{\sigma,t}$ whenever $|x|\ge R_{\sigma,t}$ and $x_1\le t-\sigma$.
	\end{lemma}
	
	\begin{proof}
		For a fixed $t>0$, set $x^t=(2t-x_1,x_2,\ldots,x_n)$,
		$\rho=|x|$, and $\rho_t=|x^t|$.  Then
		\begin{align}
			\label{eq:reflected-radius-difference}
			\rho_t^2-\rho^2=4t(t-x_1).
		\end{align}
		In particular, $\rho_t>\rho$ whenever $x_1<t$, and
		$u_B(\rho_t)-u_B(\rho)=\int_\rho^{\rho_t}v_B(s)\,ds$.
		
		\medskip
		
	Firstly, we prove the separation at infinity.  Let $\sigma\in(0,t)$ and assume $x_1\le t-\sigma$.  Then $\rho_t^2-\rho^2\ge4t\sigma$. In the normalized $1$-nondegenerate case with $\alpha>1$, the expansion gives $v_B(s)=s^\alpha(1+o(1))$.  Hence, for every fixed $\varepsilon\in(0,1)$ and all large $s$, $v_B(s)\ge(1-\varepsilon)s^\alpha$. For $s\in[\rho,\rho_t]$, the inequality $s^\alpha\ge\rho^{\alpha-1}s$ gives
		\begin{align*}
			u_B(\rho_t)-u_B(\rho)
			&\ge
			(1-\varepsilon)\rho^{\alpha-1}\int_\rho^{\rho_t}s\,ds  \\
			&=
			\frac{(1-\varepsilon)}{2}\rho^{\alpha-1}(\rho_t^2-\rho^2)
			\ge
			2(1-\varepsilon)t\sigma\,\rho^{\alpha-1}.
		\end{align*}
		Thus $u_B(\rho_t)-u_B(\rho)\to+\infty$ uniformly along sequences with
		$|x|\to\infty$ and $x_1\le t-\sigma$. On the other hand, if  $\alpha=1$,
		$v_B(s)=s-as^{-1}+bs^{-3}+o(s^{-4})$.  Integrating from $\rho$ to
		$\rho_t$, we get
		\begin{align*}
			u_B(\rho_t)-u_B(\rho)
			=
			\frac{\rho_t^2-\rho^2}{2}
			-
			a\log\frac{\rho_t}{\rho}
			-
			\frac{b}{2}
			\left(\frac{1}{\rho_t^2}-\frac{1}{\rho^2}\right)
			+
			o(1).
		\end{align*}
		Here $\rho_t\rho^{-1}\to1$ uniformly as $|x|\to\infty$ in
		$\{x_1\le t-\sigma\}$.  Indeed,
		$\rho_t^2-\rho^2=4t(t-x_1)\le4t(t+\rho)$, and hence
		$\rho_t^2\rho^{-2}=1+O(\rho^{-1})$.  Therefore the logarithmic and
		inverse-square terms vanish at infinity, and
		\begin{align*}
			\liminf\limits_{\substack{|x|\to\infty\\ x_1\le t-\sigma}}
			\bigl(u_B(\rho_t)-u_B(\rho)\bigr)
			\ge
			2t\sigma .
		\end{align*}
		Next, in the $1$-degenerate power-growth case, if the exponent
		$\frac{\alpha(k_\gamma+1)}{k_\gamma+1-2\alpha}$ is strictly larger than $1$,
		the preceding superlinear estimate applies with that exponent and gives
		$u_B(\rho_t)-u_B(\rho)\to+\infty$ uniformly on escaping sequences with
		$x_1\le t-\sigma$.  When the exponent is equal to $1$,
		$v_B(s)=A_\gamma s(1+o(1))$, and the same computation as in the linear case
		gives
		\begin{align*}
			\liminf\limits_{\substack{|x|\to\infty\\ x_1\le t-\sigma}}
			\bigl(u_B(\rho_t)-u_B(\rho)\bigr)
			\ge
			2A_\gamma t\sigma .
		\end{align*}
		Finally, in the critical $1$-degenerate case, choose
		$\kappa\in\left(0,\frac{c_\gamma}{2\alpha}\right)$.  The expansion of
		$\log v_B$ gives $v_B(s)\ge\exp(\kappa s^{2\alpha})$ for all large $s$.
		For all large $\rho$ with $x_1\le t-\sigma$, the estimate
		$\rho_t+\rho\le C_t\rho$ gives
		\begin{align*}
			\rho_t-\rho
			=
			\frac{\rho_t^2-\rho^2}{\rho_t+\rho}
			\ge
			\frac{4t\sigma}{C_t\rho}.
		\end{align*}
		Consequently,
		\begin{align*}
			u_B(\rho_t)-u_B(\rho)
			\ge
			(\rho_t-\rho)\exp(\kappa\rho^{2\alpha})
			\ge
			\frac{4t\sigma}{C_t\rho}\exp(\kappa\rho^{2\alpha})
			\to+\infty .
		\end{align*}
		
		In all cases, there exist $R_{\sigma,t}>0$ and $\mu_{\sigma,t}>0$ such that
		$u_B(\rho_t)-u_B(\rho)\ge2\mu_{\sigma,t}$ whenever
		$|x|\ge R_{\sigma,t}$ and $x_1\le t-\sigma$.  From
		\eqref{eq:C2-asymptotic-to-bowl}, after increasing $R_{\sigma,t}$ if
		necessary, we obtain
		\begin{align*}
			u_t^*(x)-u(x)
			=
			u_B(\rho_t)-u_B(\rho)+\varphi(x^t)-\varphi(x)
			\ge
			\mu_{\sigma,t}
		\end{align*}
		on the same region.
		
		\medskip
		
		Next, we obtain the ordering from a far plane.  The expansions above imply that
		there are constants $c>0$ and $r_0>0$ such that $v_B(s)\ge cs$ for
		$s\ge r_0$.  Hence there is $C>0$ such that, for all $0\le\rho<\rho_t$,
		\begin{align*}
			u_B(\rho_t)-u_B(\rho)
			\ge
			\frac{c}{2}(\rho_t^2-\rho^2)-C .
		\end{align*}
		For $x_1\le\frac{t}{2}$, one has $\rho_t^2-\rho^2\ge2t^2$.  The boundedness
		of $\varphi$ gives $u_t^*>u$ on $\{x_1\le\frac{t}{2}\}$ for every
		sufficiently large $t$.
		
		\medskip
		
		It remains for us to consider the strip $\frac{t}{2}<x_1<t$.  For
		$s\in[x_1,2t-x_1]$ and $x'=(x_2,\ldots,x_n)$, we have
		$\sqrt{s^2+|x'|^2}\ge s>\frac{t}{2}$.  Taking $t$ large enough so that this
		radius is at least $r_0$, the lower bound $v_B(r)\ge cr$ gives
		\begin{align*}
			v_B\left(\sqrt{s^2+|x'|^2}\right)
			\frac{s}{\sqrt{s^2+|x'|^2}}
			\ge
			cs
			\ge
			\frac{c}{2}t .
		\end{align*}
		Using $D\varphi\to0$ and increasing $t$ once more, we get
		$\partial_1u(s,x')>0$ on the strip.  Therefore
		\begin{align*}
			u_t^*(x)-u(x)
			=
			\int_{x_1}^{2t-x_1}\partial_1u(s,x')\,ds
			>
			0
		\end{align*}
		whenever $\frac{t}{2}<x_1<t$.  Combining both regions proves the lemma.
	\end{proof}
	
\begin{theorem}
	\label{thm:bowl-asymptotic-uniqueness}
	Let $B$ be the rotational bowl-type $\gamma$-translator introduced above, and
	assume that it extends smoothly through its tip, is strictly convex, and has a
	positive end in one of the regimes of
	Lemma~\ref{lem:start-moving-plane-at-infinity}.  Let
	$M=\operatorname{graph}(u)$ be a strictly convex entire graphical
	$\gamma$-translator which is $C^2$-asymptotic to a vertical translate of $B$.
	Then $M$ is a vertical translate of $B$.
\end{theorem}

\begin{proof}
	After the vertical normalization in \eqref{eq:C2-asymptotic-to-bowl}, it is
	enough to prove that $M=B$.  We first prove symmetry with respect to
	${x_1=0}$.
	
	```
	\medskip
	
	For $t\ge0$ and $x=(x_1,x')\in\mathbb R\times\mathbb R^{n-1}$ with
	$x_1<t$, set $h_t(x)=u_t^*(x)-u(x)=u(2t-x_1,x')-u(x_1,x')$.  Thus
	$h_t\ge0$ on $\{x_1<t\}$ means that the reflected cap $M_+^*(t)$ lies above
	$M_-(t)$.
	
	\medskip
	
	Following the moving-plane formulation in \cite{Paco_2014}, define
	$\mathcal A=\{t\in[0,\infty):h_s\ge0\text{ on }\{x_1<s\}\text{ for every }s>t\}$.
	Thus $t\in\mathcal A$ means that every plane strictly to the right of
	$\Pi_t$ is admissible.  Lemma~\ref{lem:start-moving-plane-at-infinity}
	gives a number $T>0$, chosen beyond the exterior scale where the asymptotic
	expansions of Section~\ref{sec:positive-branch-asymptotics} give the initial
	separation, such that $T\in\mathcal A$.
	
	\medskip
	
	The Alexandrov connectedness step shows that the admissible set reaches the
	limiting plane.  We first note that $\mathcal A$ is closed in
	$[0,\infty)$.  Indeed, let $t_j\in\mathcal A$ and suppose that
	$t_j\to t_0\in[0,\infty)$.  Given $s>t_0$, we have $s>t_j$ for all large
	$j$, and hence $h_s\ge0$ on $\{x_1<s\}$.  Since $s>t_0$ was arbitrary,
	$t_0\in\mathcal A$.
	
	\medskip
	
	The nonempty closed set $\mathcal A\subset[0,\infty)$ therefore has a
	minimum; write $a=\min\mathcal A$.  We prove that $a=0$.  Assume, toward a
	contradiction, that $a>0$.  Since $a\in\mathcal A$, we have
	\begin{align}
		\label{eq:order-for-s-greater-than-a}
		h_s\ge0
		\qquad
		\text{on }\{x_1<s\}
	\end{align}
	for every $s>a$.  Passing to the limit $s\to a^+$ gives
	\begin{align}
		\label{eq:order-at-a}
		h_a\ge0
		\qquad
		\text{on }\{x_1<a\}.
	\end{align}
	Indeed, for $x=(x_1,x')$ with $x_1<a$, the inequality $x_1<s$ holds for all
	$s>a$ sufficiently close to $a$, and the smoothness of $u$ gives
	$h_s(x)\to h_a(x)$.
	
	\medskip
	
	The order at $a$ is strict in the interior:
	\begin{align}
		\label{eq:strict-order-at-a}
		h_a>0
		\qquad
		\text{on }\{x_1<a\}.
	\end{align}
	Assume otherwise.  Then there exists $x_0=(x_{0,1},x_0')$, with
	$x_{0,1}<a$, such that $h_a(x_0)=0$.  At the point
	$P=(x_0,u(x_0))\in M$, the hypersurfaces $M$ and $M_+^*(a)$ are tangent, and
	\eqref{eq:order-at-a} places one locally on one side of the other.  The
	translator equation is invariant under reflection across vertical
	hyperplanes, and therefore the reflected graph is again a
	$\Gamma$-admissible $\gamma$-translator.  The tangency principle for
	$\gamma$-translators \cite[Theorem~1.4]{Yo}, applied to the connected
	graphical pieces over $\{x_1<a\}$, yields
	\begin{align}
		\label{eq:coincidence-at-a}
		h_a\equiv0
		\qquad
		\text{on }\{x_1<a\}.
	\end{align}
	Thus $M$ is symmetric with respect to $\Pi_a$.
	
	\medskip
	
	This symmetry is incompatible with the normalized asymptotics when $a>0$.
	For $R>2a$, \eqref{eq:coincidence-at-a} gives
	$u(Re_1)=u((2a-R)e_1)$.  Using
	$u(x)=u_B(|x|)+\varphi(x)$ and $\varphi(x)\to0$, we obtain
	$u_B(R)-u_B(R-2a)=o(1)$ as $R\to\infty$.  On the other hand, the regimes in
	Lemma~\ref{lem:start-moving-plane-at-infinity} imply that
	$v_B(r)\ge cr$ for all sufficiently large $r$, for some $c>0$.  Hence, for
	large $R$,
	\begin{align*}
		u_B(R)-u_B(R-2a)
		=
		\int_{R-2a}^{R}v_B(r)\,dr
		\ge
		c\int_{R-2a}^{R}r\,dr
		=
		2ca(R-a),
	\end{align*}
	which tends to $+\infty$.  This contradiction proves
	\eqref{eq:strict-order-at-a}.
	
	\medskip
	
	The plane can now be moved slightly to the left of $\Pi_a$.  Choose
	$\delta\in(0,\frac{a}{8})$.  As in the Alexandrov moving-plane proof for
	translators in \cite{Paco_2014}, we divide the domain into three regions: the
	exterior region away from the plane, the compact region away from the plane,
	and the strip adjacent to the plane.
	
	\medskip
	
	In the exterior region, the estimates in
	Lemma~\ref{lem:start-moving-plane-at-infinity} are uniform when the moving
	parameter lies in $[\frac{a}{2},\frac{3a}{2}]$.  Hence there exist
	$R_1>0$, $\mu_1>0$, and $\varepsilon_1>0$ such that
	\begin{align}
		\label{eq:outside-ball-order-near-a}
		h_s(x)\ge\mu_1
	\end{align}
	whenever $|s-a|<\varepsilon_1$, $|x|\ge R_1$, and $x_1\le s-\delta$.
	
	\medskip
	
	In the compact region, set
	$K=\{x\in\mathbb R^n: |x|\le R_1+1,\ x_1\le a-\frac{\delta}{2}\}$.
	By \eqref{eq:strict-order-at-a}, the continuous function $h_a$ has a
	positive minimum on $K$.  Since $h_s\to h_a$ uniformly on $K$ as $s\to a$,
	after decreasing $\varepsilon_1$ if necessary we have
	\begin{align}
		\label{eq:inside-ball-order-near-a}
		h_s>0
		\qquad
		\text{on }K
	\end{align}
	for every $|s-a|<\varepsilon_1$.
	
	\medskip
	
	In the strip adjacent to the plane, the boundary point step gives the needed
	one-sided control.  The function $h_a$ is nonnegative on $\{x_1<a\}$,
	positive in the interior by \eqref{eq:strict-order-at-a}, and vanishes on
	$\{x_1=a\}$.  The boundary point step in the Alexandrov argument of
	\cite{Paco_2014}, together with the tangency principle for
	$\gamma$-translators, gives the strict inward derivative at the boundary of
	the cap.  Since $\partial_1h_a(a,x')=-2\partial_1u(a,x')$ and the inward
	normal to $\{x_1<a\}$ is $-e_1$, this gives
	\begin{align}
		\label{eq:positive-normal-derivative-at-a}
		\partial_1u(a,x')>0
		\qquad
		\text{for every }x'\in\mathbb R^{n-1}.
	\end{align}
	
	\medskip
	
	The positivity is uniform along the unbounded plane.  On compact subsets of
	$\{x_1=a\}$ it follows from continuity and
	\eqref{eq:positive-normal-derivative-at-a}.  For large $|x'|$, using
	\eqref{eq:C2-asymptotic-to-bowl} and the lower bound $v_B(r)\ge cr$, we have
	\begin{align*}
		\partial_1u(a,x')
		=
		v_B\left(\sqrt{a^2+|x'|^2}\right)
		\frac{a}{\sqrt{a^2+|x'|^2}}
		+
		\partial_1\varphi(a,x')
		\ge
		\frac{ca}{2}
	\end{align*}
	for all sufficiently large $|x'|$.  After decreasing $\delta$ if necessary,
	we therefore have
	\begin{align}
		\label{eq:positive-derivative-strip-near-a}
		\partial_1u(r,x')>0
		\qquad
		\text{whenever } |r-a|<2\delta .
	\end{align}
	
	\medskip
	
	The three regions now imply that a plane slightly to the left of $\Pi_a$ is
	still admissible.  Choose $\varepsilon>0$ with
	$\varepsilon<\min\{\varepsilon_1,\frac{\delta}{2}\}$, and let
	$s>a-\varepsilon$.  For $s\ge a$, \eqref{eq:order-for-s-greater-than-a}
	gives the required inequality.  It remains to consider
	$s\in(a-\varepsilon,a)$.
	
	\medskip
	
	Fix such an $s$, and let $x=(x_1,x')$ with $x_1<s$.  In the exterior region
	$x_1\le s-\delta$, $|x|\ge R_1$, \eqref{eq:outside-ball-order-near-a} gives
	$h_s(x)>0$.  In the compact region $x_1\le s-\delta$, $|x|<R_1$, we have
	$x_1<a-\frac{\delta}{2}$, so $x\in K$, and
	\eqref{eq:inside-ball-order-near-a} gives $h_s(x)>0$.  In the remaining
	strip $s-\delta<x_1<s$, both $x_1$ and $2s-x_1$ lie in
	$(a-2\delta,a+2\delta)$, and \eqref{eq:positive-derivative-strip-near-a}
	gives
	$h_s(x)=\int_{x_1}^{2s-x_1}\partial_1u(r,x')\,dr>0$.
	Thus $h_s>0$ on $\{x_1<s\}$ for every $s\in(a-\varepsilon,a)$.  Together
	with \eqref{eq:order-for-s-greater-than-a}, this proves that
	$a-\varepsilon\in\mathcal A$, contradicting the definition of
	$a=\min\mathcal A$.
	
	\medskip
	
	Consequently $a=0$.  Since $0\in\mathcal A$, we have
	$h_s\ge0$ on $\{x_1<s\}$ for every $s>0$.  Passing to the limit
	$s\to0^+$ gives $u(-x_1,x')\ge u(x_1,x')$ for every $x_1<0$.  Repeating the
	same argument in the direction $-e_1$ gives the opposite inequality, and
	hence $u(-x_1,x')=u(x_1,x')$.  Therefore $M$ is symmetric with respect to
	$\{x_1=0\}$.
	
	\medskip
	
	The translator equation is invariant under Euclidean rotations whose linear
	part fixes $e_{n+1}$.  Applying the same reflection argument in every
	horizontal direction, we obtain rotational symmetry about the vertical axis.
	The positive-branch rotational classification identifies the normalized graph
	with the model bowl $B$.  Undoing the vertical normalization, $M$ is a
	vertical translate of $B$.
	
\end{proof}

	
	
	\subsection{Catenoidal barriers and bounded graphical domains}\,
	\label{subsec:catenoidal-barriers-bounded-domains}
	
	The bounded graphical setting for the catenoidal-barrier application is fixed
	next.  Let $\Omega\subset\mathbb R^n$ be a bounded domain.  A graphical
	translator over $\Omega$ will be written as $M=\operatorname{graph}(u)$, and the
	completeness condition used below is vertical boundary blow-up:
   $u(x)\to+\infty$ as $x\to\partial\Omega $. The result proved in this subsection excludes such graphs under a small-neck condition on the complete catenoidal family.
	
	\medskip
	
	The condition is expressed in terms of the radius at which the upper side of a
	catenoidal translator enters the positive branch.
	
	\begin{definition}
		\label{def:small-neck-catenoidal-family}
		Let $\{W_R:R\in(0,\infty)\}$ be a complete embedded catenoidal family given
		by Theorem~\ref{thm:rotational-translators}.  For each $R>0$, let $u_R^+$
		be the upper graphical side of $W_R$, and let $r_R^{\mathrm{ent}}>R$ be the
		first radius at which this side enters the positive branch $I_+$.  Write
		\begin{align*}
			W_{R,\mathrm{cvx}}^+
			=
			\{(x,u_R^+(|x|)):|x|>r_R^{\mathrm{ent}}\}
		\end{align*}
		for the corresponding positive convex tail.
		
		The family $\{W_R:R\in(0,\infty)\}$ is said to satisfy the small-neck
		entry condition if
		\begin{align}
			\label{eq:small-neck-entry-section-seven}
			\lim\limits_{R\to0^+}r_R^{\mathrm{ent}}=0.
		\end{align}
	\end{definition}
	
	\medskip
	
	The signed phase diagram describes the high-curvature direction followed by the
	upper-entry sheet.  At the neck of $W_R$, the rotational curvature is
	$y=\frac{1}{R}$; hence the regime $R\to0^+$ corresponds, in the selected signed
	chart, to $y\to\infty$.  On the upper-entry sheet, the level
	$\widehat\gamma=1$ is represented by the graph $x=h_+(y)$.  Equivalently, in
	terms of the selected signed chart, $h_+(y)=g(y,1)$ and $\gamma(g(y,1),y)=1$. Writing $a(y)=y^{-1}g(y,1)$, homogeneity gives $\Psi_g(a(y))=y^{-\alpha}$.  The small-neck entry condition records the
	geometric consequence needed for comparison: as the neck shrinks, the positive
	convex tail begins arbitrarily close to the rotation axis.
	
\begin{theorem}
	\label{thm:first-contact-convex-catenoidal-tail}
	Assume that there exists a complete embedded catenoidal family
	${W_R:R\in(0,\infty)}$ satisfying the small-neck entry condition.  Then
	there is no complete convex $\Gamma$-admissible graphical
	$\gamma$-translator $M=\operatorname{graph}(u)$, with
	$u:\Omega\to\mathbb R$, over a bounded connected domain
	$\Omega\subset\mathbb R^n$, satisfying $u(x)\to+\infty$ as $x\to\partial\Omega$.
\end{theorem}

The barrier argument is illustrated in
Figure~\ref{fig:catenoidal-barrier-configuration}.  The small-neck entry
condition allows us to choose the catenoidal tail so that, after centering it
outside the vertical cylinder over $\Omega$, its positive convex part is defined
over the whole bounded domain.



\begin{proof}
	To reach a contradiction, assume that such a translator exists.  Choose
	$q\in\mathbb R^n\setminus\overline\Omega$.  By the small-neck entry
	condition, we may choose $R>0$ such that
	$r_R^{\mathrm{ent}}<\operatorname{dist}(q,\overline\Omega)$.  Center $W_R$
	around the vertical line $q+\mathbb R e_{n+1}$.  Then
	$|x-q|>r_R^{\mathrm{ent}}$ for every $x\in\overline\Omega$, and therefore
	$w(x)=u_R^+(|x-q|)$ is defined entirely on the positive convex tail
	$W_{R,\mathrm{cvx}}^+$ for $x\in\Omega$.  Thus $\operatorname{graph}(w)$ is a
	smooth convex $\Gamma$-admissible graphical $\gamma$-translator over
	$\Omega$.  Moreover, $w$ extends smoothly to a neighborhood of
	$\overline\Omega$, since $|x-q|$ ranges in a compact subset of
	$(r_R^{\mathrm{ent}},\infty)$ on $\overline\Omega$.
	
	\medskip
	
	We now move the graph $M$ vertically.  For $t\in\mathbb R$, consider the
	vertical translate $M+t e_{n+1}=\operatorname{graph}(u+t)$, and define $\mathcal T=\{t\in\mathbb R:w\le u+t\text{ in }\Omega\}$. We notice that this set is nonempty.  Indeed, $w$ is bounded on $\overline\Omega$, while
	$u$ is finite in $\Omega$ and tends to $+\infty$ at the boundary; hence
	$u-w$ is bounded from below on $\Omega$.  Therefore $w\le u+t$ in $\Omega$
	for every sufficiently large $t$.
	
	\medskip
	
	The set $\mathcal T$ is bounded from below.  Fixing any $p_0\in\Omega$, every
	$t\in\mathcal T$ satisfies $t\ge w(p_0)-u(p_0)$.  Hence
	$t_*=\inf\mathcal T$ is finite.  Since $w\le u+t$ in $\Omega$ for every
	$t>t_*$, passage to the infimum gives
	\begin{align}
		\label{eq:critical-barrier-order}
		w\le u+t_*
		\qquad
		\text{in }\Omega .
	\end{align}
		
	The critical translate must touch the catenoidal tail at an interior point.
	Suppose otherwise.  Then $u+t_*-w>0$ in $\Omega$.  The boundedness of $w$ on
	$\overline\Omega$ and the boundary blow-up of $u$ give
	$u(x)+t_*-w(x)\to+\infty$ as $x\to\partial\Omega$.  Hence there exists a
	compact set $K_*\subset\Omega$ such that $u+t_*-w\ge1$ in
	$\Omega\setminus K_*$.  On $K_*$, the continuous function $u+t_*-w$ has a
	positive minimum.  Therefore there exists $\varepsilon>0$ such that
	$u+t_*-w\ge\varepsilon$ in $\Omega$.  It follows that
	$w\le u+t_*-\frac{\varepsilon}{2}$ in $\Omega$, contradicting the definition
	of $t_*=\inf\mathcal T$.  Consequently, there exists $p\in\Omega$ such that
	$w(p)=u(p)+t_*$.
	
	\medskip
	
	At the contact point $p$, the function $u+t_*-w$ has a local minimum equal to
	zero.  Therefore $Du(p)=Dw(p)$, and the graphs of $u+t_*$ and $w$ are tangent
	at $(p,w(p))$ with the same upward unit normal.  Moreover,
	\eqref{eq:critical-barrier-order} places $\operatorname{graph}(w)$ locally on
	one side of $M+t_*e_{n+1}$.
	
	\medskip
	
	The positive convex tail is a $\Gamma$-admissible graphical
	$\gamma$-translator.  The vertical translate $M+t_*e_{n+1}$ is also a
	$\Gamma$-admissible $\gamma$-translator, since vertical translations preserve
	the principal curvatures and the upward unit normal.  Thus the tangency
	principle for $\gamma$-translators \cite[Theorem~1.4]{Yo} applies at the
	contact point and gives $w=u+t_*$ in $\Omega$ . However, this is impossible: $w$ is bounded on $\overline\Omega$, whereas $u(x)+t_*\to+\infty$ as $x\to\partial\Omega$.  The contradiction proves the theorem.	
\end{proof}

	\begin{corollary}
		\label{cor:bounded-cylinder-nonexistence}
		Under the hypotheses of
		Proposition~\ref{thm:first-contact-convex-catenoidal-tail}, there is no
		complete convex $\Gamma$-admissible graphical $\gamma$-translator contained
		in a bounded vertical cylinder and complete in the sense of vertical boundary
		blow-up.
	\end{corollary}
	
	\begin{proof}
		Such a translator is a graph over a bounded domain and satisfies the boundary
		blow-up condition stated above.  This is exactly the situation excluded by
		Proposition~\ref{thm:first-contact-convex-catenoidal-tail}.
	\end{proof}
	
	
\begin{figure}[htbp]
	\centering
	\includegraphics[width=0.69\textwidth]{catenoidal_barrier_mcmanus.png}
	\caption{\footnotesize Catenoidal barrier configuration for a bounded graphical domain.
		After choosing a small neck and centering the catenoidal family outside
		the cylinder over $\Omega$, the positive convex tail is defined over the
		whole domain and can be used in the vertical translation argument.  Image
		courtesy of Ignacio McManus.}
	\label{fig:catenoidal-barrier-configuration}
\end{figure}
	
	
	
\section{Conclusions and further directions}
\label{sec:conclusions}

In this paper we have developed a rotational translator theory for fully
nonlinear extrinsic curvature functions through two complementary branch
constructions.  On the positive side, the selected branch gives the asymptotic
structure of bowl-type ends.  On the signed side, the neck datum and the charts
$g$ and $g^\sharp$ organize the continuation of catenoidal profiles through
mixed-sign curvature regions.  The resulting alternatives show that the lower
end of a fully nonlinear catenoidal translator is governed not only by
homogeneity and ellipticity, but also by the topology of the signed level
diagram.

\medskip

We have also shown how this phase analysis interacts with comparison
principles.  The positive-end asymptotics provide the separation at infinity
needed for the Alexandrov moving-plane argument, following the
mean-curvature-flow strategy for translators asymptotic to the translating
paraboloid \cite{Paco_2014}.  The catenoidal-barrier argument uses complete
small-neck families to exclude bounded graphical domains by a first-contact
argument inside the positive elliptic branch.

\medskip

We finalize the article by listening open problems in the $\gamma$-translators theory:

\begin{op}
	The uniqueness theorem in
	Section~\ref{sec:uniqueness-catenoidal-barriers} assumes $C^2$-asymptotics to
	a prescribed rotational bowl.  In mean curvature flow, Wang proved uniqueness
	of the strictly convex entire graphical translator without imposing this
	asymptotic hypothesis \cite{wang2011convex}.  A natural fully nonlinear problem
	is to determine structural conditions on $\gamma$ under which every complete
	strictly convex entire graphical $\Gamma$-admissible translator is a vertical
	translate of the rotational bowl.
\end{op}

\begin{op}
	Convex translators in slabs and $\Delta$-wing translators are fundamental
	non-entire examples in mean curvature flow
	\cite{wang2011convex,bourni2020existence,HMWS-deltawings}.  In the
	$\gamma$-flow setting, the corresponding construction should first be
	considered for curvature functions admitting suitable grim-reaper barriers and
	complete catenoidal families with small necks.  One may then ask whether
	$\Delta$-wing $\gamma$-translators exist in slabs and whether their asymptotics
	are governed by the positive branch and the signed catenoidal alternatives
	developed here.
\end{op}

\begin{op}
	The stability of the translating paraboloid for mean curvature flow was
	studied by Clutterbuck--Schn\"urer--Schulze
	\cite{clutterbuck2007stability}.  A fully nonlinear analogue should be
	formulated in classes where the rotational bowl is smooth, the positive-end
	asymptotics are precise, and small-neck catenoidal barriers are available.  In
	such a class, one may ask whether graphical solutions starting close to the
	bowl in a weighted topology converge, after vertical normalization, to a
	translate of the bowl.
\end{op}


	\appendix
	
\section{Elementary properties of $Q(r,s)$ and $y_v(r)$}
\label{app:q-eta-properties}

Throughout this appendix, we collect the elementary identities used in
Sections~\ref{sec:barrier-method}, \ref{sec:positive-branch-asymptotics}, and
\ref{sec:catenoidal-translators}.  We assume $\alpha>0$ and set
$\beta=\frac{\alpha-1}{2\alpha}$.  Then
\begin{align}
	\label{eq:beta-identities-appendix}
	1-2\beta=\alpha^{-1}>0,
	\mbox{ so }
	\beta<\frac12 .
\end{align}
Recall that
\begin{align}
	\label{eq:q-def-appendix}
	Q(r,s)
	=
	\frac{s}{r(1+s^2)^\beta}.
\end{align}
For a graphical slope $v$, we write $y_v(r)=Q(r,v(r))$.  Thus $y_v$ is the
rotational phase coordinate of the graphical solution, consistently with the
abstract phase variables $(x,y)$, where $x$ denotes the meridional curvature and
$y$ denotes the rotational curvature.

\medskip

For each fixed $r>0$, the map $s\mapsto Q(r,s)$ is strictly increasing.  Indeed,
\begin{align}
	\label{eq:q-partial-s}
	\partial_s Q(r,s)
	=
	\frac{1+(1-2\beta)s^2}
	{r(1+s^2)^{\beta+1}},
\end{align}
and \eqref{eq:beta-identities-appendix} gives
$\partial_s Q(r,s)>0$ for every $r>0$ and every $s\in\mathbb R$.

\medskip

For each fixed $r>0$, we have
\begin{align}
	\label{eq:q-limits}
	\lim\limits_{s\to+\infty}Q(r,s)=+\infty,
	\qquad
	\lim\limits_{s\to-\infty}Q(r,s)=-\infty.
\end{align}
The reason is that
\begin{align}
	\label{eq:q-infinity-asymptotic}
	Q(r,s)
	\sim
	r^{-1}\operatorname{sgn}(s)|s|^{1-2\beta}
	\mbox{ as } |s|\to\infty,
\end{align}
and $1-2\beta>0$.  Together with strict monotonicity, this proves that
\begin{align}
	\label{eq:q-bijection}
	Q(r,\cdot):\mathbb R\to\mathbb R
\end{align}
is a bijection for every $r>0$.

\medskip

The function $Q(r,s)$ has the same sign as $s$, since
$(1+s^2)^\beta>0$.  Thus
\begin{align}
	\label{eq:q-sign}
	Q(r,s)s\ge0,
	\quad
	Q(r,s)=0\quad\Longleftrightarrow\quad s=0.
\end{align}
Moreover,
\begin{align}
	\label{eq:q-partial-r}
	\partial_rQ(r,s)
	=
	-\frac{s}{r^2(1+s^2)^\beta}
	=
	-r^{-1}Q(r,s).
\end{align}

\begin{lemma}
	\label{lem:level-set-control}
	For every $m\in\mathbb R$ and every $r>0$, there exists a unique number
	$w_m(r)\in\mathbb R$ such that
	\begin{align}
		\label{eq:wm-definition}
		Q(r,w_m(r))=m.
	\end{align}
	The function $w_m$ is of class $\mathcal C^1$ on $(0,\infty)$.  Moreover,
	$w_m(r)$ has the same sign as $m$, $w_0(r)=0$, and if $a<b$, then
	$w_a(r)<w_b(r)$ for every $r>0$.  Finally, for every function $v$ defined at
	$r$,
	\begin{align}
		\label{eq:wm-trapping}
		w_a(r)\le v(r)\le w_b(r)
		\quad\Longleftrightarrow\quad
		a\le Q(r,v(r))\le b;
	\end{align}
	in particular,
	\begin{align}
		\label{eq:eta-trapping}
		w_a(r)\le v(r)\le w_b(r)
		\quad\Longleftrightarrow\quad
		a\le y_v(r)\le b .
	\end{align}
\end{lemma}

\begin{proof}
	The existence and uniqueness of $w_m(r)$ follow from the bijectivity of
	$Q(r,\cdot)$.  The sign property follows from \eqref{eq:q-sign}, while the
	ordering and the trapping equivalences follow from the strict monotonicity of
	$s\mapsto Q(r,s)$. The regularity of $w_m$ follows from the implicit function theorem applied to
	$Q(r,w_m(r))=m$, since \eqref{eq:q-partial-s} gives $\partial_sQ>0$.
\end{proof}

Differentiating \eqref{eq:wm-definition} and using
\eqref{eq:q-partial-r}--\eqref{eq:q-partial-s}, we obtain
\begin{align}
	\label{eq:wm-derivative}
	w_m'(r)
	=
	\frac{w_m(r)(1+w_m(r)^2)}
	{r\left(1+(1-2\beta)w_m(r)^2\right)}.
\end{align}
Equivalently, the identity $Q(r,w_m(r))=m$ gives
\begin{align}
	\label{eq:wm-derivative-level}
	w_m'(r)
	=
	(1+w_m(r)^2)^{\beta+1}
	\frac{m}{1+(1-2\beta)w_m(r)^2}.
\end{align}

Let $(I_\sigma,g_\sigma)$ be a selected graphical branch, and let $m\in I_\sigma$.
For the residual
$\mathcal R_\sigma[w]=w'-(1+w^2)^{\beta+1}g_\sigma(Q(r,w))$, the level function
$w_m$ satisfies
\begin{align}
	\label{eq:wm-residual-branch}
	\mathcal R_\sigma[w_m]
	=
	(1+w_m^2)^{\beta+1}
	\left[
	\frac{m}{1+(1-2\beta)w_m^2}
	-
	g_\sigma(m)
	\right].
\end{align}
This follows by substituting \eqref{eq:wm-derivative-level} and
$Q(r,w_m(r))=m$ into the definition of $\mathcal R_\sigma$.

\medskip

Let $v$ be a $\mathcal C^1$ function.  Differentiating
$y_v(r)=Q(r,v(r))$ gives
\begin{align}
	\label{eq:eta-derivative-general}
	y_v'(r)
	=
	\frac{1}{r(1+v^2)^{\beta+1}}
	\left(1+(1-2\beta)v^2\right)v'
	-
	\frac{1}{r}y_v(r).
\end{align}

Assume that $v$ solves the implicit graphical ODE
\begin{align}
	\label{eq:branch-ode-appendix}
	v'
	=
	(1+v^2)^{\beta+1}g(y_v,z),
\end{align}
where $z$ is a fixed level and $g$ is the selected implicit branch.  Substituting
\eqref{eq:branch-ode-appendix} into \eqref{eq:eta-derivative-general} gives
\begin{align}
	\label{eq:eta-derivative-branch}
	y_v'
	=
r^{-1}
	\left(
	\left(1+(1-2\beta)v^2\right)g(y_v,z)
	-
	y_v
	\right).
\end{align}

When $v(r)\to+\infty$,
\begin{align}
	\label{eq:eta-large-slope-asymptotic}
	y_v(r)
	=
	v(r)^{\frac{1}{\alpha}}r^{-1}(1+o(1)).
\end{align}
More precisely, under the same assumption,
\begin{align}
	\label{eq:eta-large-slope-second-order}
	y_v(r)
	=
	v(r)^{\frac{1}{\alpha}}r^{-1}
	\left(1-\beta v(r)^{-2}+O(v(r)^{-4})\right).
\end{align}
Indeed, \eqref{eq:q-def-appendix} gives
$y_v(r)=r^{-1}v(r)^{1-2\beta}(1+v(r)^{-2})^{-\beta}$.  Since
$1-2\beta=\alpha^{-1}$, Taylor expansion gives
$(1+v(r)^{-2})^{-\beta}=1-\beta v(r)^{-2}+O(v(r)^{-4})$, which proves
\eqref{eq:eta-large-slope-asymptotic} and
\eqref{eq:eta-large-slope-second-order}.

\medskip

For every $c>0$, $r>0$, and $s\in\mathbb R$,
\begin{align}
	\label{eq:q-scaling}
	Q(cr,s)=c^{-1}Q(r,s).
\end{align}
Consequently, if $Q(r,w_m(r))=m$, then $Q(cr,w_m(r))=mc^{-1}$.

\subsection{Level-control lemmas used in the catenoidal construction},
\label{subsec:appendix-catenoidal-level-control}

The following elementary consequences are used in
Section~\ref{sec:catenoidal-translators} to keep the phase variable inside a
chosen component.

\begin{lemma}
	\label{lem:level-control-upper-entry}
	Let $v$ be branch-admissible for a selected branch $(I_\sigma,g_\sigma)$ on
	an interval $J\subset(0,\infty)$.  For levels $a<b$ in $I_\sigma$, the
	inequality $w_a\le v\le w_b$ on $J$ is equivalent to
	\begin{align*}
		a\le y_v(r)\le b
		\qquad
		\text{for every }r\in J.
	\end{align*}
\end{lemma}

\begin{proof}
	This is exactly \eqref{eq:eta-trapping}.
\end{proof}

\begin{lemma}
	\label{lem:level-function-scaling}
	For every fixed level $m$ and every $c>0$, the function $r\mapsto w_m(r)$
	satisfies $Q(cr,w_m(r))=\frac{m}{c}$.  In particular, changing the radial scale
	changes the level of a fixed slope by the factor $c^{-1}$.
\end{lemma}

\begin{proof}
	This is the scaling identity \eqref{eq:q-scaling} applied to the defining
	relation $Q(r,w_m(r))=m$.
\end{proof}

\section*{Declarations}
\textbf{Funding:}
No funding was received to assist with the preparation of this manuscript.

\begin{bibdiv}
	\begin{biblist}

\bib{alencar2026halfspace}{article}{
			author={Alencar, H.},
			author={Bessa, G. P.},
			author={Silva Neto, G.},
			title={Half-space theorems for translating solitons of the
				$r$-mean curvature flow},
			journal={Journal of Differential Equations},
			volume={460},
			pages={114115},
			year={2026},
			doi={10.1016/j.jde.2026.114115}
		}
		
		

\bib{andrews-mccoy-zheng-contracting}{article}{
			author={Andrews, B.},
			author={McCoy, J. A.},
			author={Zheng, Y.},
			title={Contracting convex hypersurfaces by curvature},
			journal={Calculus of Variations and Partial Differential Equations},
			volume={47},
			number={3--4},
			pages={611--665},
			year={2013},
			doi={10.1007/s00526-012-0530-3}
		}
		
		

\bib{bourni2020existence}{article}{
			author={Bourni, T.},
			author={Langford, M.},
			author={Tinaglia, G.},
			title={On the existence of translating solutions of mean curvature flow in slab regions},
			journal={Analysis \& PDE},
			volume={13},
			number={4},
			pages={1051--1072},
			year={2020}
		}
		
		

\bib{clutterbuck2007stability}{article}{
			author={Clutterbuck, J.},
			author={Schn{\"u}rer, O. C.},
			author={Schulze, F.},
			title={Stability of translating solutions to mean curvature flow},
			journal={Calculus of Variations and Partial Differential Equations},
			volume={29},
			number={3},
			pages={281--293},
			year={2007},
			doi={10.1007/s00526-006-0033-1}
		}
		
		

\bib{cogo-lynch-vicanek2023rotational}{article}{
			author={Cogo, A.},
			author={Lynch, S.},
			author={Vi{\v c}\'anek Mart\'inez, O.},
			title={Rotational symmetry of ancient solutions to fully nonlinear curvature flows},
			journal={arXiv preprint arXiv:2310.08301},
			year={2023},
			eprint={2310.08301}
		}
		
		

\bib{HIMW-notes}{incollection}{
			author={Hoffman, D.},
			author={Ilmanen, T.},
			author={Mart{\'i}n, F.},
			author={White, B.},
			title={Notes on Translating Solitons for Mean Curvature Flow},
			booktitle={Minimal Surfaces: Integrable Systems and Visualisation},
			editor={Hoffmann, T.},
			editor={Kilian, M.},
			editor={Leschke, K.},
			editor={Mart{\'i}n, F.},
			series={Springer Proceedings in Mathematics \& Statistics},
			volume={349},
			pages={231--280},
			publisher={Springer},
			address={Cham},
			year={2021},
			doi={10.1007/978-3-030-68541-6\_9}
		}
		
		

\bib{hoffman2019graphical}{article}{
			author={Hoffman, D.},
			author={Ilmanen, T.},
			author={Mart{\'i}n, F.},
			author={White, B.},
			title={Graphical translators for mean curvature flow},
			journal={Calculus of Variations and Partial Differential Equations},
			volume={58},
			number={4},
			pages={117},
			year={2019},
			doi={10.1007/s00526-019-1560-x}
		}
		
		

\bib{HMW-scherk}{article}{
			author={Hoffman, D.},
			author={Mart{\'i}n, F.},
			author={White, B.},
			title={Scherk-like translators for mean curvature flow},
			journal={Journal of Differential Geometry},
			volume={122},
			number={3},
			pages={421--465},
			year={2022},
			doi={10.4310/jdg/1675712995}
		}
		
		

\bib{HMW-translating-annuli}{article}{
			author={Hoffman, D.},
			author={Mart{\'i}n, F.},
			author={White, B.},
			title={Translating annuli for mean curvature flow},
			journal={Advances in Mathematics},
			volume={455},
			pages={109875},
			year={2024},
			doi={10.1016/j.aim.2024.109875}
		}
		
		

\bib{HMWS-deltawings}{article}{
			author={Hoffman, D.},
			author={Mart{\'i}n, F.},
			author={White, B.},
			title={Annuloids and $\Delta$-wings},
			journal={Advanced Nonlinear Studies},
			volume={24},
			number={1},
			pages={74--96},
			year={2024},
			doi={10.1515/ans-2023-0111}
		}
		
		

\bib{HMWS-tridents}{article}{
			author={Hoffman, D.},
			author={Mart{\'i}n, F.},
			author={White, B.},
			title={Nguyen's tridents and the classification of semigraphical translators for mean curvature flow},
			journal={Journal f\"ur die reine und angewandte Mathematik (Crelle's Journal)},
			volume={786},
			pages={79--105},
			year={2022},
			doi={10.1515/crelle-2022-0005}
		}
		
		

\bib{lima2025translators}{article}{
			author={de Lima, R. F.},
			author={Pipoli, G.},
			title={Translators to Higher Order Mean Curvature Flows in
				$\mathbb R^n\times\mathbb R$ and $\mathbb H^n\times\mathbb R$},
			journal={The Journal of Geometric Analysis},
			volume={35},
			number={3},
			pages={92},
			year={2025},
			doi={10.1007/s12220-025-01892-3}
		}
		
		

\bib{mantegazza2011lecture}{book}{
			author={Mantegazza, C.},
			title={Lecture Notes on Mean Curvature Flow},
			series={Progress in Mathematics},
			volume={290},
			publisher={Birkh\"auser/Springer Basel AG},
			place={Basel},
			date={2011},
			doi={10.1007/978-3-0348-0145-4}
		}
		
		

\bib{martin2024classification}{article}{
			author={Mart{\'i}n, F.},
			author={S{\'a}ez, M.},
			author={Tsiamis, R.},
			author={White, B.},
			title={Classification of Semigraphical Translators},
			journal={arXiv preprint arXiv:2411.16889},
			year={2024},
			eprint={2411.16889}
		}
		
		

\bib{Paco_2014}{article}{
			author={Mart{\'i}n, F.},
			author={Savas-Halilaj, A.},
			author={Smoczyk, K.},
			title={On the topology of translating solitons of the mean curvature flow},
			journal={Calculus of Variations and Partial Differential Equations},
			volume={54},
			number={3},
			pages={2853--2882},
			year={2015},
			doi={10.1007/s00526-015-0886-2}
		}
		
		

\bib{rengaswami2024classification}{article}{
			author={Rengaswami, S.},
			title={Classification of Bowl-Type Translators to Fully Nonlinear Curvature Flows},
			journal={The Journal of Geometric Analysis},
			volume={34},
			number={6},
			pages={Article 186},
			year={2024},
			doi={10.1007/s12220-024-01641-6}
		}
		
		

\bib{RengaswamiTorresSantaella2023Bowl}{article}{
			author={Rengaswami, S.},
			author={Torres Santaella, J.},
			title={Bowl Soliton Asymptotics and Applications},
			journal={arXiv preprint arXiv:2310.09984},
			year={2023},
			eprint={2310.09984}
		}
		
		

\bib{spruck-xiao}{article}{
			author={Spruck, J.},
			author={Xiao, L.},
			title={Complete translating solitons to the mean curvature flow in
				$\mathbb R^3$ with nonnegative mean curvature},
			journal={American Journal of Mathematics},
			volume={142},
			number={3},
			pages={993--1015},
			year={2020},
			doi={10.1353/ajm.2020.0023}
		}
		
		

\bib{teschl2012ordinary}{book}{
			author={Teschl, G.},
			title={Ordinary Differential Equations and Dynamical Systems},
			series={Graduate Studies in Mathematics},
			volume={140},
			publisher={American Mathematical Society},
			place={Providence, RI},
			year={2012}
		}
		
		

\bib{wang2011convex}{article}{
			author={Wang, X.-J.},
			title={Convex solutions to the mean curvature flow},
			journal={Annals of Mathematics},
			volume={173},
			number={3},
			pages={1185--1239},
			year={2011},
			doi={10.4007/annals.2011.173.3.1}
		}
		
		

\bib{Yo}{article}{
			author={Torres Santaella, J.},
			title={Maximum principles and consequences for $\gamma$-translators in
				$\mathbb R^{n+1}$},
			journal={Bulletin of the London Mathematical Society},
			volume={56},
			number={4},
			pages={1433--1448},
			year={2024},
			doi={10.1112/blms.13004}
		}
		
		

\bib{Yo2}{article}{
			author={Torres Santaella, J.},
			title={Maximum principles and consequences for $\gamma$-translators in
				$\mathbb R^{n+1}$ II},
			journal={Journal of Mathematical Analysis and Applications},
			volume={552},
			number={2},
			pages={129809},
			year={2025},
			doi={10.1016/j.jmaa.2025.129809}
		}
		
	

	\end{biblist}
\end{bibdiv}


\end{document}
